% Consciencisme, panafricanisme et développement de l’Afrique — Philosophie 1re Tech — Cours signé OnBuch+
% Programme officiel MINESEC. Compiler avec : tectonic N11-consciencisme-panafricanisme-developpement-afrique.tex
\newcommand{\DOCMATIERE}{Philosophie}
\newcommand{\DOCNIVEAU}{1re Tech}
\newcommand{\DOCMODULE}{Philosophie – classe de Première technique}
\newcommand{\DOCLECON}{N11}
\newcommand{\DOCTITRE}{Consciencisme, panafricanisme et développement de l’Afrique}
% OnBuch+ — préambule « cours signés » ENRICHI (compilé avec tectonic / XeLaTeX)
% Système visuel « Pop » d'OnBuch+ : fond crème (jamais blanc), bordures encre
% épaisses, ombres « dures » décalées, titres Archivo Black en capitales.
% Version MATHÉMATIQUES : graphes (pgfplots/tikz), boîtes pédagogiques
% (définition, propriété, méthode, attention, le savais-tu, exercice résolu,
% exercice, TP, bilan), page de garde et signature OnBuch+ ancrée en bas.
%
% Macros attendues dans main.tex AVANT \input{preamble} :
%   \DOCMATIERE  \DOCNIVEAU  \DOCMODULE  \DOCLECON (numéro)  \DOCTITRE
\documentclass[11pt]{article}

\usepackage{fontspec}
\usepackage[french]{babel}
\usepackage[a4paper,margin=2cm,top=3.4cm,bottom=2.8cm,headheight=1.4cm,footskip=1.3cm]{geometry}
\usepackage{amsmath,amssymb,amsfonts}
\usepackage{mathtools} % \xmapsto, \coloneqq... souvent utilisés par les modèles
\usepackage{cancel} % simplifications barrées (bilans de forces)
% \abs{z} (module, valeur absolue) et \norm{u} (norme), utilisés par les modèles
\providecommand{\abs}{}\let\abs\relax\DeclarePairedDelimiter{\abs}{\lvert}{\rvert}
\providecommand{\norm}{}\let\norm\relax\DeclarePairedDelimiter{\norm}{\lVert}{\rVert}
\usepackage[table]{xcolor} % [table] : fournit aussi \rowcolors (alternance de lignes)
\usepackage{pagecolor}
\usepackage{tikz}
\usetikzlibrary{arrows.meta,positioning,calc,shapes.geometric,shapes.misc,shapes.arrows,decorations.pathmorphing,decorations.pathreplacing,decorations.markings,patterns,shadows,mindmap,trees,fit,backgrounds,matrix,intersections,angles,quotes}
\usepackage{pgfplots}
\usepackage[version=4]{mhchem} % équations nucléaires \ce{^{235}_{92}U}
\usepackage[european,siunitx]{circuitikz} % schémas électriques
\pgfplotsset{compat=1.18}
\usepgfplotslibrary{fillbetween,groupplots} % aires entre courbes (calcul intégral)
\usepackage{enumitem}
\usepackage{fancyhdr}
% [most] charge toutes les bibliothèques tcolorbox (listings, minted, raster,
% documentation...) et gonfle inutilement la mémoire TeX sur un document long
% (~300 boîtes) : on ne charge que ce qui est réellement utilisé.
\usepackage[skins,breakable]{tcolorbox}
\usepackage{graphicx}
\usepackage{colortbl}
\usepackage{booktabs}
\usepackage{tabularx}
\usepackage{multirow}
\usepackage{diagbox} % case d'en-tête diagonale des tableaux à double entrée
% Colonnes extensibles centrée / gauche / droite (C, L, R), souvent utilisées par les modèles
\newcolumntype{C}{>{\centering\arraybackslash}X}
\newcolumntype{L}{>{\raggedright\arraybackslash}X}
\newcolumntype{R}{>{\raggedleft\arraybackslash}X}
\usepackage{array}
\usepackage{multicol}
\usepackage{siunitx}
% parse-numbers=false : accepte aussi \SI{2,5 \times 10^{-3}}{...}
\sisetup{locale=FR,output-decimal-marker={,},group-minimum-digits=4,parse-numbers=false}
\DeclareSIUnit\ohm{\ensuremath{\symup{\Omega}}} % U+2126 absent de la police
\DeclareSIUnit\division{div} % réglages d'oscilloscope (ms/div, V/div)
\DeclareSIUnit\tr{tr} % tours (vitesse de rotation, ex. \SI{1500}{\tr\per\minute})
\DeclareSIUnit\molar{M} % concentration molaire (ex. \SI{3}{\molar}, \SI{1}{\micro\molar})
\DeclareSIUnit\an{an} % année (le modèle confond parfois avec \year, un compteur TeX) — ex. \SI{5}{\kilo\meter\per\an}
\DeclareSIUnit\FCFA{FCFA} % franc CFA (ex. \SI{500}{\FCFA\per\kilo\gram})
\DeclareSIUnit\degreeBrix{\textdegree Brix} % teneur en sucre d'un sirop/jus (réfractométrie)
\DeclareSIUnit\month{mois} % le modèle utilise parfois \month, un compteur TeX, comme unité de durée
\DeclareSIUnit\microliter{\micro\liter} % le modèle écrit parfois \microliter en un seul token
\DeclareSIUnit\million{millions} % dénombrement (ex. \SI{15}{\million\per\mL} spermatozoïdes)
\usepackage{qrcode}
% \degree : commande fréquemment utilisée par les modèles pour un angle ou une
% température en dehors de \SI{}{} (non fournie par les paquets chargés).
\providecommand{\degree}{^{\circ}}
% \qed (fin de démonstration), utilisé par les modèles sans amsthm
\providecommand{\qed}{\hfill$\square$}
% \barre : double barre des valeurs interdites dans un tableau de variations (tkz-tab)
\providecommand{\barre}{\ensuremath{\|}}
% environnement proof (amsthm non chargé) utilisé par les modèles
\ifdefined\proof\else\newenvironment{proof}[1][Démonstration]{\par\noindent\textit{#1.}\ }{\qed\par}\fi
\usepackage{lastpage}
\usepackage{hyperref}
\hypersetup{hidelinks}

% ============================================================
% Palette « Pop » OnBuch+ — mêmes hex que la classe P (plus_ui.dart)
% ============================================================
\definecolor{poporange}{HTML}{F59321}
\definecolor{popdark}{HTML}{E07A0C}
\definecolor{popcream}{HTML}{FFF6E8}
\definecolor{popink}{HTML}{1C1714}
\definecolor{poppurple}{HTML}{7A5AE0}
\definecolor{popgreen}{HTML}{2E9E6A}
\definecolor{poppink}{HTML}{E0457B}
\definecolor{popblue}{HTML}{2F7DE1}
\definecolor{popgold}{HTML}{C98A12}
\definecolor{popmuted}{HTML}{8A7F76}
\definecolor{popline}{HTML}{EADFD2}
\definecolor{popgray}{HTML}{8A7F76}
\colorlet{popgrey}{popgray}
% Alias tolérants (le modèle nomme parfois les couleurs différemment)
\colorlet{popwhite}{white}
\colorlet{popblack}{popink}
\colorlet{popred}{poppink}
\colorlet{popyellow}{popgold}
% Teintes claires (fonds de tableaux/encadrés) — le modèle nomme parfois
% les couleurs avec un suffixe L (ex. popblueL) sans les avoir définies.
\colorlet{poporangeL}{poporange!20}
\colorlet{popdarkL}{popdark!20}
\colorlet{poppurpleL}{poppurple!20}
\colorlet{popgreenL}{popgreen!20}
\colorlet{poppinkL}{poppink!20}
\colorlet{popblueL}{popblue!20}
\colorlet{popgoldL}{popgold!20}
\colorlet{popredL}{popred!20}
\colorlet{popgrayL}{popgray!20}
% Teintes claires pour fonds de tableaux / zones
\colorlet{popblueL}{popblue!10!white}
\colorlet{poporangeL}{poporange!14!white}
\colorlet{popgreenL}{popgreen!12!white}
\colorlet{poppurpleL}{poppurple!12!white}
\colorlet{poppinkL}{poppink!12!white}
\colorlet{popgoldL}{popgold!14!white}

\pagecolor{popcream}

% ============================================================
% Typographie
% ============================================================
\defaultfontfeatures{Ligatures=TeX}
\setmainfont{PlusJakartaSans-Medium.ttf}[Path=../../../fonts/,BoldFont=PlusJakartaSans-Bold.ttf]
% Maths en sans-serif (Fira Math) avec chiffres et lettres droites de Plus
% Jakarta, pour que les formules se fondent dans le texte.
\usepackage[math-style=ISO,bold-style=ISO]{unicode-math}
\setmathfont{FiraMath-Regular.otf}[Path=../../../fonts/]
\setmathfont{PlusJakartaSans-Medium.ttf}[Path=../../../fonts/,range={up/num,up/{Latin,latin},\mathup}]
\usepackage{icomma} % virgule décimale sans espace en mode math
% Caractères mathématiques Unicode absents des polices de texte (Plus Jakarta,
% Archivo Black) : rendus par leur équivalent mathématique, en texte comme en
% formule — sinon ils s'affichent en carré vide (ex. « Arithmétique dans ℤ »).
\usepackage{newunicodechar}
\newunicodechar{ℝ}{\ensuremath{\mathbb{R}}}
\newunicodechar{ℤ}{\ensuremath{\mathbb{Z}}}
\newunicodechar{ℂ}{\ensuremath{\mathbb{C}}}
\newunicodechar{ℕ}{\ensuremath{\mathbb{N}}}
\newunicodechar{ℚ}{\ensuremath{\mathbb{Q}}}
\newunicodechar{𝔻}{\ensuremath{\mathbb{D}}}
\newunicodechar{α}{\ensuremath{\alpha}}
\newunicodechar{β}{\ensuremath{\beta}}
\newunicodechar{γ}{\ensuremath{\gamma}}
\newunicodechar{δ}{\ensuremath{\delta}}
\newunicodechar{ε}{\ensuremath{\varepsilon}}
\newunicodechar{θ}{\ensuremath{\theta}}
\newunicodechar{λ}{\ensuremath{\lambda}}
\newunicodechar{μ}{\ensuremath{\mu}}
\newunicodechar{σ}{\ensuremath{\sigma}}
\newunicodechar{τ}{\ensuremath{\tau}}
\newunicodechar{φ}{\ensuremath{\varphi}}
\newunicodechar{ψ}{\ensuremath{\psi}}
\newunicodechar{ω}{\ensuremath{\omega}}
\newunicodechar{Σ}{\ensuremath{\Sigma}}
% majuscules produites par \MakeUppercase dans les titres (α-aminés -> Α)
\newunicodechar{Α}{\ensuremath{\alpha}}
\newunicodechar{Β}{\ensuremath{\beta}}
\newunicodechar{Γ}{\ensuremath{\Gamma}}
\newunicodechar{Δ}{\ensuremath{\Delta}}
\newunicodechar{Ε}{\ensuremath{\varepsilon}}
\newunicodechar{Θ}{\ensuremath{\Theta}}
\newunicodechar{Λ}{\ensuremath{\Lambda}}
\newunicodechar{Μ}{\ensuremath{\mu}}
\newunicodechar{Τ}{\ensuremath{\tau}}
\newunicodechar{Φ}{\ensuremath{\Phi}}
\newunicodechar{Ψ}{\ensuremath{\Psi}}
\newunicodechar{Ω}{\ensuremath{\Omega}}
\newunicodechar{↦}{\ensuremath{\mapsto}}
\newunicodechar{∈}{\ensuremath{\in}}
\newunicodechar{∉}{\ensuremath{\notin}}
\newunicodechar{∧}{\ensuremath{\wedge}}
\newunicodechar{⇔}{\ensuremath{\Leftrightarrow}}
\newunicodechar{⟺}{\ensuremath{\Longleftrightarrow}}
\newunicodechar{⇒}{\ensuremath{\Rightarrow}}
\newunicodechar{⟹}{\ensuremath{\Longrightarrow}}
\newunicodechar{∘}{\ensuremath{\circ}}
\newunicodechar{∪}{\ensuremath{\cup}}
\newunicodechar{∩}{\ensuremath{\cap}}
\newunicodechar{≡}{\ensuremath{\equiv}}
\newunicodechar{⊂}{\ensuremath{\subset}}
\newunicodechar{⊥}{\ensuremath{\perp}}
\newunicodechar{∃}{\ensuremath{\exists}}
\newunicodechar{∀}{\ensuremath{\forall}}
\newunicodechar{∛}{\ensuremath{\sqrt[3]{\,}}}
\newunicodechar{∜}{\ensuremath{\sqrt[4]{\,}}}
\newunicodechar{⃗}{\ensuremath{{}^{\to}}}
\newunicodechar{ⁿ}{\ifmmode^{n}\else\textsuperscript{\ensuremath{n}}\fi}
\newunicodechar{ˣ}{\ifmmode^{x}\else\textsuperscript{\ensuremath{x}}\fi}
\newunicodechar{ᵇ}{\ifmmode^{b}\else\textsuperscript{\ensuremath{b}}\fi}
\newunicodechar{ʳ}{\ifmmode^{r}\else\textsuperscript{\ensuremath{r}}\fi}
\newunicodechar{ᵘ}{\ifmmode^{u}\else\textsuperscript{\ensuremath{u}}\fi}
\newunicodechar{ʸ}{\ifmmode^{y}\else\textsuperscript{\ensuremath{y}}\fi}
\newunicodechar{ᵃ}{\ifmmode^{a}\else\textsuperscript{\ensuremath{a}}\fi}
\newunicodechar{ᵉ}{\ifmmode^{e}\else\textsuperscript{\ensuremath{e}}\fi}
\newunicodechar{ᵏ}{\ifmmode^{k}\else\textsuperscript{\ensuremath{k}}\fi}
\newunicodechar{ᵖ}{\ifmmode^{p}\else\textsuperscript{\ensuremath{p}}\fi}
\newunicodechar{ᴺ}{\ifmmode^{N}\else\textsuperscript{\ensuremath{N}}\fi}
\newunicodechar{ᶜ}{\ifmmode^{c}\else\textsuperscript{\ensuremath{c}}\fi}
\newunicodechar{ʰ}{\ifmmode^{h}\else\textsuperscript{\ensuremath{h}}\fi}
\newunicodechar{ᵠ}{\ifmmode^{q}\else\textsuperscript{\ensuremath{q}}\fi}
\newunicodechar{ₙ}{\ifmmode_{n}\else\textsubscript{\ensuremath{n}}\fi}
\newunicodechar{ₐ}{\ifmmode_{a}\else\textsubscript{\ensuremath{a}}\fi}
\newunicodechar{ᵢ}{\ifmmode_{i}\else\textsubscript{\ensuremath{i}}\fi}
\newunicodechar{ᵣ}{\ifmmode_{r}\else\textsubscript{\ensuremath{r}}\fi}
\newunicodechar{ₖ}{\ifmmode_{k}\else\textsubscript{\ensuremath{k}}\fi}
\newunicodechar{ᵦ}{\ifmmode_{\beta}\else\textsubscript{\ensuremath{\beta}}\fi}
\newunicodechar{ₓ}{\ifmmode_{x}\else\textsubscript{\ensuremath{x}}\fi}
\newfontfamily\popdisplay{ArchivoBlack-Regular.ttf}[Path=../../../fonts/]
\newfontfamily\poplogo{SpaceGrotesk-Bold.ttf}[Path=../../../fonts/]
\newfontfamily\popsemi{PlusJakartaSans-SemiBold.ttf}[Path=../../../fonts/]
\newfontfamily\popxbold{PlusJakartaSans-ExtraBold.ttf}[Path=../../../fonts/]
\newfontfamily\popxboldls{PlusJakartaSans-ExtraBold.ttf}[Path=../../../fonts/,LetterSpace=3.0]

\setlist[enumerate]{leftmargin=1.7em,itemsep=5pt,topsep=4pt}
\setlist[itemize]{leftmargin=1.7em,itemsep=4pt,topsep=4pt,label={\color{poporange}\textbullet}}
\setlength{\parindent}{0pt}
\setlength{\parskip}{5pt}
\renewcommand{\arraystretch}{1.25}

% pgfplots : style Pop par défaut
\pgfplotsset{
  every axis/.append style={axis line style={popink,line width=1pt},tick style={popink},
    grid style={popline,line width=0.5pt},label style={font=\small},tick label style={font=\footnotesize}},
  popcurve/.style={poporange,line width=1.8pt,no markers,smooth},
}
% Même style disponible dans un \draw TikZ simple (hors pgfplots)
\tikzset{popcurve/.style={poporange,line width=1.8pt}}
\tikzset{popgrid/.style={popline,very thin}}
\colorlet{popcurve}{poporange} % le nom du style est parfois utilisé comme couleur

% ============================================================
% Pill / squiggle
% ============================================================
\newcommand{\poppill}[3]{%
  \tikz[baseline=-0.55ex]\node[draw=popink,line width=1.2pt,rounded corners=2.6pt,%
    fill=#2,inner xsep=6.5pt,inner ysep=3pt]{%
    \popxboldls\fontsize{7}{7}\selectfont\color{#3}\MakeUppercase{#1}};%
}
\newcommand{\popsquiggle}[1]{%
  \tikz[baseline=-0.4ex]{
    \draw[#1,line width=2.6pt,line cap=round]
      plot [smooth] coordinates {(0,0.09) (0.34,0.01) (0.68,0.21) (1.02,0.01) (1.36,0.10)};
  }%
}
% Mot-clé surligné (marqueur orange sous le texte)
\newcommand{\cle}[1]{{\popxbold\color{popdark}#1}}

% ============================================================
% En-tête / pied
% ============================================================
\pagestyle{fancy}
\fancyhf{}
\renewcommand{\headrulewidth}{0pt}
\renewcommand{\footrulewidth}{0pt}
\lhead{\raisebox{-0.15ex}{\poplogo\fontsize{13}{13}\selectfont\color{popink}OnBuch\color{poporange}+}}
\rhead{\poppill{\DOCMATIERE\ \textbullet\ \DOCNIVEAU\ \textbullet\ Leçon \DOCLECON}{white}{popink}}
\lfoot{\popxbold\fontsize{7.6}{7.6}\selectfont\color{popmuted}\MakeUppercase{Cours signé OnBuch+}\ \textbullet\ {\popsemi\DOCTITRE}}
\rfoot{\tikz[baseline=-0.6ex]\node[rounded corners=8.5pt,fill=popink,minimum height=17pt,minimum width=17pt,inner xsep=5pt,inner sep=0]{\popxbold\fontsize{8}{8}\selectfont\color{popcream}\thepage\,/\,\pageref*{LastPage}};}

% ============================================================
% Titres
% ============================================================
\newcommand{\chaptitre}[2]{%
  \begin{tcolorbox}[enhanced,colback=poporange,colframe=popink,boxrule=2.6pt,arc=11pt,
      drop shadow=popink,left=15pt,right=15pt,top=12pt,bottom=11pt,before skip=3pt,after skip=18pt]
    \poppill{#2}{popink}{popcream}\par\vspace{9pt}
    {\raggedright\hyphenpenalty=10000\exhyphenpenalty=10000\popdisplay\fontsize{22}{24}\selectfont\color{popcream}\MakeUppercase{#1}\par}
  \end{tcolorbox}
}
% Section numérotée : pastille orange avec le numéro + titre Archivo
\newcounter{popsec}
\newcommand{\coursec}[1]{%
  \stepcounter{popsec}\setcounter{popsub}{0}%
  \par\vspace{16pt}\noindent
  \begin{minipage}{\linewidth}
  \tikz[baseline=-0.7ex]\node[draw=popink,line width=1.6pt,fill=poporange,rounded corners=4pt,minimum size=22pt,inner sep=0]{\popdisplay\fontsize{12}{12}\selectfont\color{popcream}\thepopsec};%
  \hspace{8pt}{\popdisplay\fontsize{15}{17}\selectfont\color{popink}\MakeUppercase{#1}}\par
  \vspace{4pt}\noindent{\color{poporange}\rule{36pt}{3.6pt}}
  \end{minipage}\par\nopagebreak\vspace{8pt}
}
\newcounter{popsub}
\newcommand{\courssub}[1]{%
  \stepcounter{popsub}%
  \par\vspace{9pt}\noindent{\popxbold\fontsize{11.5}{13}\selectfont\color{popink}{\color{poporange}\thepopsec.\thepopsub}\hspace{6pt}#1}\par\nopagebreak\vspace{4pt}
}

% ============================================================
% Boîtes pédagogiques — même recette : fond blanc, bordure épaisse colorée,
% ombre dure encre, badge plein. Titre optionnel [#1].
% ============================================================
\tcbset{popbase/.style={breakable,enhanced,colback=white,boxrule=2.3pt,arc=9pt,
  shadow={3pt}{-3pt}{0pt}{popink},left=14pt,right=14pt,top=11pt,bottom=11pt,before skip=11pt,after skip=15pt,
  fontupper=\popsemi\fontsize{10.6}{14.5}\selectfont\color{popink},
  fontlower=\popsemi\fontsize{10.6}{14.5}\selectfont\color{popink}}}
% Coche dessinée (le glyphe ✓ n'existe pas dans les polices du document)
\newcommand{\popcheck}[1]{\tikz[baseline=-0.5ex]\draw[#1,line width=1.6pt,line cap=round,line join=round](-0.13,0.0)--(-0.04,-0.09)--(0.13,0.1);}
\providecommand{\coche}{\popcheck{popgreen}}
\providecommand{\resultat}[1]{\par\smallskip\noindent\fcolorbox{popgreen}{popgreenL}{\parbox{\dimexpr\linewidth-2\fboxsep-2\fboxrule\relax}{\textbf{Résultat.} #1}}\par\smallskip}
\newcommand{\popboxhead}[3]{\poppill{#1}{#2}{white}\ifx\relax#3\relax\else\hspace{6pt}{\popxbold\fontsize{10.6}{12}\selectfont\color{popink}#3}\fi\par\vspace{7pt}}

\newtcolorbox{aretenir}[1][]{popbase,colframe=popgreen,before upper={\popboxhead{À retenir}{popgreen}{#1}}}
\newtcolorbox{exemplebox}[1][]{popbase,colframe=poporange,before upper={\popboxhead{Exemple}{poporange}{#1}}}
\newtcolorbox{definition}[1][]{popbase,colframe=popblue,before upper={\popboxhead{Définition}{popblue}{#1}}}
\newtcolorbox{propriete}[1][]{popbase,colframe=poppurple,before upper={\popboxhead{Propriété}{poppurple}{#1}}}
\newtcolorbox{methode}[1][]{popbase,colframe=popgold,before upper={\popboxhead{Méthode}{popgold}{#1}}}
\newtcolorbox{attention}[1][]{popbase,colframe=poppink,before upper={\popboxhead{Attention}{poppink}{#1}}}
\newtcolorbox{experience}[1][]{popbase,colframe=popblue,colback=popblueL,before upper={\popboxhead{Expérience}{popblue}{#1}}}
% « Le savais-tu ? » : carte violette pleine (contexte, culture, Cameroun)
\newtcolorbox{savaistu}[1][]{popbase,colback=poppurple,colframe=popink,
  fontupper=\popsemi\fontsize{10.4}{14.2}\selectfont\color{white},
  before upper={\poppill{Le savais-tu\,?}{white}{poppurple}\ifx\relax#1\relax\else\hspace{6pt}{\popxbold\color{white}#1}\fi\par\vspace{7pt}}}
% Exercice résolu : énoncé en haut, \tcblower puis la solution sur fond crème
\newcounter{popexo}\newcounter{popexr}
\newtcolorbox{exoresolu}[1][]{popbase,colframe=poporange,colbacklower=popcream,
  segmentation style={popink,line width=1.2pt,dashed},
  before upper={\stepcounter{popexr}\popboxhead{Exercice résolu \thepopexr}{poporange}{#1}},
  before lower={\poppill{Solution}{popink}{popcream}\par\vspace{6pt}}}
% Exercice (non résolu en place) — la correction est dans la partie Corrigés
\newtcolorbox{exercice}[1][]{popbase,colframe=popink,
  before upper={\stepcounter{popexo}\popboxhead{Exercice \thepopexo}{popink}{#1}}}
% Corrigé
\newtcolorbox{corrige}[1][]{popbase,colframe=popgreen,colback=popgreenL,
  before upper={\popboxhead{Corrigé}{popgreen}{#1}}}
% Objectifs / prérequis / bilan
\newtcolorbox{objectifs}{popbase,colframe=popink,colback=poporangeL,
  before upper={\popboxhead{Objectifs de la leçon}{popink}{}}}
\newtcolorbox{prerequis}{popbase,colframe=popink,colback=popblueL,
  before upper={\popboxhead{Prérequis}{popblue}{}}}
\newtcolorbox{bilan}{popbase,colframe=popink,colback=white,boxrule=2.8pt,
  before upper={\popboxhead{Fiche bilan}{poporange}{L'essentiel à connaître pour l'examen}}}
% Figure encadrée avec légende
\newtcolorbox{popfigure}[1][]{enhanced,breakable=false,colback=white,colframe=popline,boxrule=1.4pt,arc=8pt,
  left=10pt,right=10pt,top=10pt,bottom=8pt,before skip=10pt,after skip=12pt,halign=center,
  lower separated=false,halign lower=center,
  fontlower=\popsemi\fontsize{9}{11.5}\selectfont\color{popmuted},
  #1}
\newcommand{\legende}[1]{\par\vspace{6pt}{\popsemi\fontsize{9}{11.5}\selectfont\color{popmuted}#1}}

% ============================================================
% Signature OnBuch+
% ============================================================
\newcommand{\poplogoinline}[1]{{\poplogo\fontsize{#1}{#1}\selectfont\color{popink}OnBuch\color{poporange}+}}
\newcommand{\signatureonbuch}{%
  \begin{tcolorbox}[enhanced,colback=popink,colframe=popink,boxrule=2.2pt,arc=10pt,
      drop shadow=poporange,left=14pt,right=14pt,top=10pt,bottom=10pt,before skip=0pt,after skip=0pt]
    \tikz[baseline=-0.7ex]\node[circle,fill=popgreen,draw=popcream,line width=1.3pt,minimum size=22pt,inner sep=0]{\popcheck{white}};%
    \hspace{9pt}\begin{minipage}[c]{0.62\linewidth}
      {\popxbold\fontsize{12}{13}\selectfont\color{popcream}Signé\ \textbullet\ {\poplogo OnBuch\color{poporange}+}}\par\vspace{2pt}
      {\raggedright\popsemi\fontsize{8}{10.5}\selectfont\color{popline}\MakeUppercase{Équipe pédagogique OnBuch+}\\\MakeUppercase{Conforme au programme officiel MINESEC}\par}
    \end{minipage}\hfill
    {\poplogo\fontsize{22}{22}\selectfont\color{popcream}OnBuch\color{poporange}+}
  \end{tcolorbox}%
}

% ============================================================
% Page de garde
% #1 titre · #2 matière · #3 niveau · #4 module · #5 n° leçon · #6 durée ·
% #7 description / objectifs (texte ou itemize)
% ============================================================
\newcommand{\pagegarde}[7]{%
  \thispagestyle{empty}
  \newgeometry{a4paper,margin=1.7cm,top=1.6cm,bottom=1.4cm}%
  \begin{tikzpicture}[remember picture,overlay]
    \fill[poporange!18] ([xshift=-3.2cm,yshift=-3.6cm]current page.north east) circle (4.6cm);
    \fill[poppurple!14] ([xshift=2.4cm,yshift=7.6cm]current page.south west) circle (3.8cm);
  \end{tikzpicture}%
  {\poplogo\fontsize{30}{30}\selectfont\color{popink}OnBuch\color{poporange}+}\hfill
  \raisebox{0.6ex}{\poppill{Cours signé \textbullet\ Premium}{popink}{popcream}}\par
  \vspace{1pt}\popsquiggle{poppurple}\par
  \vspace{8pt}
  \begin{tcolorbox}[enhanced,colback=poporange,colframe=popink,boxrule=2.8pt,arc=13pt,
      drop shadow=popink,left=18pt,right=18pt,top=12pt,bottom=13pt,after skip=10pt]
    \poppill{#2 \textbullet\ #3}{popink}{popcream}\hspace{5pt}\poppill{#4}{white}{popink}\par\vspace{8pt}
    {\popdisplay\fontsize{14}{14}\selectfont\color{popink}LEÇON #5}\par\vspace{4pt}
    {\raggedright\hyphenpenalty=10000\exhyphenpenalty=10000\popdisplay\fontsize{25}{27}\selectfont\color{popcream}\MakeUppercase{#1}\par}
    \vspace{8pt}
    {\popxbold\fontsize{9.5}{9.5}\selectfont\color{popink}\MakeUppercase{Durée conseillée : #6}}
  \end{tcolorbox}
  \begin{tcolorbox}[enhanced,colback=white,colframe=popink,boxrule=2.2pt,arc=10pt,
      drop shadow=popink,left=16pt,right=16pt,top=10pt,bottom=10pt,after skip=8pt]
    \poppill{Dans cette leçon}{poppurple}{white}\par\vspace{5pt}
    \popsemi\fontsize{9.3}{12.6}\selectfont\color{popink}#7
  \end{tcolorbox}
  \noindent\begin{minipage}[t]{0.487\linewidth}
    \begin{tcolorbox}[enhanced,colback=popgreen,colframe=popink,boxrule=2.2pt,arc=10pt,
        drop shadow=popink,left=11pt,right=11pt,top=10pt,bottom=10pt,height=3.1cm,valign=center]
      \tcbox[colback=white,colframe=popink,boxrule=1.3pt,arc=5pt,boxsep=0pt,nobeforeafter,
        left=3pt,right=3pt,top=3pt,bottom=3pt,valign=center]{\qrcode[height=1.9cm]{https://play.google.com/store/apps/details?id=com.luvvix.onbuch&hl=fr}}%
      \hspace{8pt}\begin{minipage}[c]{0.5\linewidth}
        {\popdisplay\fontsize{11}{12.5}\selectfont\color{white}\MakeUppercase{Télécharge OnBuch}}\par\vspace{4pt}
        {\raggedright\popsemi\fontsize{8.4}{11}\selectfont\color{white}Gratuit \textbullet\ Google Play\\Tuteur IA, annales, concours, résultats\par}
      \end{minipage}
    \end{tcolorbox}
  \end{minipage}\hfill
  \begin{minipage}[t]{0.487\linewidth}
    \begin{tcolorbox}[enhanced,colback=poppurple,colframe=popink,boxrule=2.2pt,arc=10pt,
        drop shadow=popink,left=11pt,right=11pt,top=10pt,bottom=10pt,height=3.1cm,valign=center]
      \tikz[baseline=-0.7ex]\node[circle,fill=white,draw=popink,line width=1.3pt,minimum size=1.6cm,inner sep=0]{\popdisplay\fontsize{24}{24}\selectfont\color{poppurple}+};%
      \hspace{8pt}\begin{minipage}[c]{0.6\linewidth}
        {\popdisplay\fontsize{11}{12.5}\selectfont\color{white}\MakeUppercase{Passe à OnBuch+}}\par\vspace{4pt}
        {\raggedright\popsemi\fontsize{8.4}{11}\selectfont\color{white}Tous les cours signés, Tuteur IA illimité, fiches PDF — onglet OnBuch+ de l'app\par}
      \end{minipage}
    \end{tcolorbox}
  \end{minipage}
  \vspace{8pt}
  \popcontenu
  \vfill
  \signatureonbuch
  \restoregeometry
  \newpage
}

% Rangée « Ce cours contient » (page de garde)
\newcommand{\poptile}[3]{% #1 couleur #2 gros texte #3 légende
  \begin{tcolorbox}[enhanced,colback=#1,colframe=popink,boxrule=2pt,arc=9pt,shadow={2.5pt}{-2.5pt}{0pt}{popink},
      left=6pt,right=6pt,top=9pt,bottom=9pt,halign=center,valign=center,height=1.75cm,width=\linewidth,nobeforeafter]
    {\popdisplay\fontsize{12.5}{13}\selectfont\color{white}\MakeUppercase{#2}}\par\vspace{3pt}
    {\centering\popsemi\fontsize{7.8}{9.5}\selectfont\color{white}#3\par}
  \end{tcolorbox}}
\newcommand{\popcontenu}{%
  \par\noindent{\popxboldls\fontsize{7.5}{7.5}\selectfont\color{popmuted}CE COURS SIGNÉ CONTIENT}\par\vspace{7pt}
  \noindent\begin{minipage}[t]{0.235\linewidth}\poptile{popblue}{Cours}{complet, illustré, pas à pas}\end{minipage}\hfill
  \begin{minipage}[t]{0.235\linewidth}\poptile{poporange}{Méthodes}{et exercices résolus}\end{minipage}\hfill
  \begin{minipage}[t]{0.235\linewidth}\poptile{poppink}{Exercices}{type examen + corrigés}\end{minipage}\hfill
  \begin{minipage}[t]{0.235\linewidth}\poptile{popgreen}{Fiche bilan}{l'essentiel à retenir}\end{minipage}\par}

% Sommaire
\newtcolorbox{sommaire}{enhanced,colback=white,colframe=popink,boxrule=2.2pt,arc=10pt,
  drop shadow=popink,left=18pt,right=18pt,top=13pt,bottom=13pt,before skip=6pt,after skip=20pt,
  before upper={\poppill{Sommaire}{poppurple}{white}\par\vspace{9pt}\popsemi\fontsize{10.6}{14.5}\selectfont\color{popink}}}

% Signature de fin de cours — poussée tout en bas de la dernière page
\newcommand{\findecours}{%
  \par\vfill\nopagebreak
  \begin{center}\popsquiggle{poporange}\end{center}\vspace{4pt}
  \signatureonbuch
}
\AtBeginDocument{\def\llbracket{\mathopen{[\mkern-3.5mu[}}\def\rrbracket{\mathclose{]\mkern-3.5mu]}}} % intervalles d'entiers (glyphe ⟦ absent de la police)
% Glyphes absents de Fira Math : redéfinis à partir de glyphes disponibles
\AtBeginDocument{%
  \def\setminus{\mathbin{\backslash}}\let\smallsetminus\setminus
  \def\vdots{\mathord{\vbox{\baselineskip3.2pt\lineskiplimit0pt\kern4pt\hbox{.}\hbox{.}\hbox{.}}}}%
  \def\ddots{\mathinner{\mkern1mu\raise6.5pt\hbox{.}\mkern2mu\raise3.6pt\hbox{.}\mkern2mu\raise0.7pt\hbox{.}\mkern1mu}}%
  \def\triangle{\mathord{\scalebox{0.9}{$\symup{\Delta}$}}}%
  \def\nabla{\mathord{\raisebox{\depth}{\scalebox{1}[-1]{$\symup{\Delta}$}}}}%
  \def\checkmark{\popcheck{popdark}}%
  \def\star{\mathbin{\ast}}\def\bigstar{\mathord{\ast}}%
  \def\diamond{\mathbin{\scalebox{0.6}{$\Diamond$}}}%
  \def\bigcup{\mathop{\mathchoice{\raisebox{-0.3em}{\scalebox{1.7}{$\cup$}}}{\scalebox{1.25}{$\cup$}}{\cup}{\cup}}\displaylimits}%
  \def\bigcap{\mathop{\mathchoice{\raisebox{-0.3em}{\scalebox{1.7}{$\cap$}}}{\scalebox{1.25}{$\cap$}}{\cap}{\cap}}\displaylimits}%
  \def\lesseqqgtr{\mathrel{\lessgtr}}\def\bigoplus{\mathop{\oplus}\displaylimits}%
}
\providecommand{\ssi}{si et seulement si} % abréviation fréquente
\AtBeginDocument{\providecommand{\wideparen}{\overparen}} % arc de cercle

\begin{document}

\pagegarde{Consciencisme, panafricanisme et développement de l’Afrique}{Philosophie}{1re Tech}{Philosophie – classe de Première technique}{N11}{2 séances (fiche de progression harmonisée)}{Cette leçon étudie le consciencisme de Kwame Nkrumah, ses fondements idéologiques et philosophiques, ainsi que l'historique du panafricanisme, de ses origines diasporiques à l'Union africaine. Elle invite l'élève à appliquer la notion d'unité africaine à des situations concrètes du Cameroun et à construire une argumentation philosophique rigoureuse.
\vspace{4pt}
\begin{multicols}{2}\small
\begin{itemize}
\item À la fin de cette leçon, je sais définir le consciencisme et le panafricanisme à partir de textes et de documents
\item À la fin de cette leçon, je sais expliquer les fondements idéologiques (marxisme, socialisme africain, anti-impérialisme) du consciencisme
\item À la fin de cette leçon, je sais dégager les fondements philosophiques du consciencisme (matérialisme, humanisme, conscience révolutionnaire)
\item À la fin de cette leçon, je sais retracer les grandes étapes historiques du panafricanisme (congrès, OUA, UA) sur une frise et une carte
\item À la fin de cette leçon, je sais identifier les figures majeures du panafricanisme (Du Bois, Garvey, Nkrumah, Senghor, Césaire, Um Nyobè)
\item À la fin de cette leçon, je sais appliquer la notion d'unité africaine à des situations concrètes de la vie courante au Cameroun (franc CFA, CEMAC, crises régionales)
\end{itemize}\end{multicols}}

\begin{sommaire}
\begin{multicols}{2}
\begin{enumerate}
\item Introduction : problématique et définitions
\item Contexte historique : colonisation, indépendances et quête d'identité
\item Les fondements idéologiques du consciencisme
\item Les fondements philosophiques du consciencisme
\item Historique du panafricanisme
\item Unité africaine et développement : applications concrètes au Cameroun
\item Méthode : rédiger et justifier une démarche argumentative
\item Activité d'intégration
\item Méthodes et exercices résolus
\item Exercices
\item Corrigés des exercices
\item Fiche bilan
\end{enumerate}
\end{multicols}
\end{sommaire}

\begin{objectifs}
\begin{itemize}
\item À la fin de cette leçon, je sais définir le consciencisme et le panafricanisme à partir de textes et de documents
\item À la fin de cette leçon, je sais expliquer les fondements idéologiques (marxisme, socialisme africain, anti-impérialisme) du consciencisme
\item À la fin de cette leçon, je sais dégager les fondements philosophiques du consciencisme (matérialisme, humanisme, conscience révolutionnaire)
\item À la fin de cette leçon, je sais retracer les grandes étapes historiques du panafricanisme (congrès, OUA, UA) sur une frise et une carte
\item À la fin de cette leçon, je sais identifier les figures majeures du panafricanisme (Du Bois, Garvey, Nkrumah, Senghor, Césaire, Um Nyobè)
\item À la fin de cette leçon, je sais appliquer la notion d'unité africaine à des situations concrètes de la vie courante au Cameroun (franc CFA, CEMAC, crises régionales)
\item À la fin de cette leçon, je sais analyser et commenter un extrait de texte de Nkrumah en dégageant la thèse et les arguments
\item À la fin de cette leçon, je sais rédiger et justifier une démarche argumentative (introduction, développement, conclusion) sur une question de développement africain
\item À la fin de cette leçon, je sais évaluer de manière critique l'actualité du panafricanisme face aux défis du développement de l'Afrique
\end{itemize}
\end{objectifs}

\begin{prerequis}
\begin{itemize}
\item Notion de colonisation et de décolonisation de l'Afrique (histoire de 3e et de Seconde)
\item Notion d'idéologie et de conscience (leçons de philosophie de Première sur la conscience)
\item Repères sur les indépendances africaines (années 1960, indépendance du Cameroun en 1960-1961)
\item Notion de liberté et d'aliénation étudiée en début d'année
\item Connaissance générale des organisations internationales (ONU) vue en éducation à la citoyenneté
\end{itemize}
\end{prerequis}

\newpage

\coursec{Introduction : problématique et définitions}

\courssub{Situation d'entrée : un débat au lycée de Nkongsamba}

En 2024, lors d'un débat au lycée de Nkongsamba, des élèves s'interrogent : pourquoi, plus de soixante ans après les indépendances, l'Afrique reste-t-elle dépendante des puissances étrangères alors que des leaders comme Kwame Nkrumah avaient proposé le \emph{consciencisme} comme voie de développement autonome ? Au Cameroun, les discussions sur la monnaie (le franc CFA), sur l'unité nationale et sur l'intégration régionale (CEMAC, Union africaine) renvoient directement à ces idéologies nées de la lutte anticoloniale.

Ce débat de lycéens n'est pas anodin. Il traduit une inquiétude réelle et largement partagée : les États africains sont politiquement indépendants depuis les années 1960, mais beaucoup d'observateurs estiment que leur \cle{indépendance économique} reste inachevée. Les matières premières (cacao, pétrole, bois, minerais) partent souvent brutes vers l'extérieur ; les produits manufacturés reviennent à des prix élevés ; les décisions monétaires dépendent parfois d'institutions situées hors du continent. Face à ce constat, deux attitudes s'affrontent : les uns jugent que l'Afrique doit d'abord renforcer son \cle{unité} pour peser dans le monde ; les autres pensent que chaque État doit régler ses problèmes seul. Ce débat n'est pas seulement économique ou politique : il est \emph{philosophique}, car il porte sur la question de savoir ce que l'Afrique doit être, ce qu'elle peut penser par elle-même et comment elle peut transformer sa situation.

\begin{exemplebox}[La question du franc CFA]
Le franc CFA, utilisé au Cameroun et dans plusieurs pays d'Afrique centrale et de l'Ouest, est arrimé à l'euro et son fonctionnement implique historiquement la France. Certains y voient une garantie de stabilité ; d'autres, un signe de dépendance héritée de la colonisation. Ce débat concret illustre exactement le type de questions auxquelles le consciencisme et le panafricanisme prétendent répondre : un développement authentiquement africain est-il possible dans des cadres hérités de la domination coloniale ?
\end{exemplebox}

\courssub{Problématique de la leçon}

La situation précédente peut être condensée en une question directrice :

\begin{center}
\textbf{Comment penser un développement authentiquement africain après la colonisation ?}
\end{center}

Cette problématique se déploie en plusieurs questions secondaires que la leçon traitera successivement :
\begin{itemize}
\item Quelles sont les origines historiques du consciencisme et du panafricanisme ? (contexte de la colonisation et des indépendances)
\item Sur quels fondements idéologiques et philosophiques repose le consciencisme de Nkrumah ?
\item Comment le panafricanisme s'est-il construit au fil du temps, des premiers congrès à l'Union africaine ?
\item Ces idées peuvent-elles encore éclairer les choix de développement du Cameroun et de l'Afrique aujourd'hui ?
\end{itemize}

\begin{attention}[Ne pas réduire la problématique à un slogan]
La question n'est pas « faut-il être pour ou contre l'unité africaine ? ». Une problématique philosophique exige d'\emph{examiner} les concepts (indépendance, développement, unité, aliénation), de confronter les thèses et de justifier une réponse argumentée. Au Probatoire et au Baccalauréat technique, le correcteur attend une réflexion problématisée, pas une prise de position militante.
\end{attention}

\courssub{Définition du consciencisme}

Commençons par l'intuition. Quand un peuple a été longtemps dominé, il ne suffit pas de chasser le dominateur pour être libre : il faut aussi se libérer les \emph{idées}, c'est-à-dire cesser de penser avec les catégories de l'ancien maître. Le consciencisme part de cette intuition simple : la libération politique doit être prolongée par une libération de la \cle{conscience}.

\begin{definition}[Le consciencisme]
Le \cle{consciencisme} est une doctrine philosophique élaborée par Kwame Nkrumah (1909-1972), premier président du Ghana indépendant, dans son ouvrage \emph{Consciencism: Philosophy and Ideology for Decolonization} publié en 1964. Il se définit comme une synthèse entre le \textbf{matérialisme dialectique} (selon lequel la réalité matérielle et les contradictions sociales expliquent l'histoire) et les \textbf{valeurs traditionnelles africaines} (l'égalitarisme, l'humanisme, la vie communautaire). Nkrumah le présente comme \og l'arme philosophique du panafricanisme \fg{} : une théorie destinée à guider la décolonisation des esprits et la construction d'une Afrique unie et socialiste.
\end{definition}

Précisons la structure de cette doctrine. Pour Nkrumah, la conscience africaine est le produit de trois expériences historiques superposées : la tradition africaine originelle (fondée sur la communauté et l'égalité), l'apport de l'islam et l'apport du christianisme européen lié à la colonisation. Le consciencisme propose de \emph{dépasser} les tensions entre ces trois héritages en les harmonisant dans une philosophie nouvelle, adaptée aux réalités africaines. Le mot lui-même est révélateur : il s'agit de donner à l'Africain une \emph{conscience} claire de sa situation, de sa dignité et de sa mission historique.

\begin{savaistu}[Un mot inventé par Nkrumah]
Le mot \og consciencisme \fg{} vient de l'anglais \emph{consciencism}, un terme forgé par Nkrumah lui-même en 1964 pour le titre de son ouvrage. Il ne s'agit donc pas d'un mot du dictionnaire, mais d'un néologisme philosophique, construit à partir de \emph{conscience}. En créant un mot nouveau, Nkrumah affirme symboliquement que l'Afrique peut produire ses propres concepts au lieu d'emprunter ceux des autres.
\end{savaistu}

\courssub{Définition du panafricanisme}

L'intuition du panafricanisme est tout aussi simple : les peuples africains, et tous les descendants d'Africains dispersés dans le monde par la traite négrière, partagent une histoire commune d'humiliation et de résistance ; ils ont donc intérêt à s'unir. Le préfixe grec \emph{pan-} signifie \og tout, ensemble \fg{} (comme dans \og panorama \fg{}) : le panafricanisme est littéralement la volonté de penser l'Afrique \emph{entière}, et au-delà d'elle, toute la diaspora noire.

\begin{definition}[Le panafricanisme]
Le \cle{panafricanisme} est un mouvement politique et culturel, né à la fin du XIX\textsuperscript{e} siècle dans la diaspora noire (Amériques, Caraïbes), qui vise l'\textbf{unité} et la \textbf{solidarité} des peuples africains et des personnes d'ascendance africaine dans le monde. Sur le plan politique, il revendique la libération du continent de la domination coloniale puis l'union des États africains ; sur le plan économique, la coopération et le développement solidaire ; sur le plan culturel, la revalorisation des cultures et de la dignité noires.
\end{definition}

Contrairement au consciencisme, le panafricanisme n'est pas l'œuvre d'un seul auteur : c'est un \emph{mouvement historique} traversé par de nombreuses figures (W. E. B. Du Bois, Marcus Garvey, George Padmore, Léopold Sédar Senghor, Kwame Nkrumah, entre autres) et par des étapes institutionnelles (congrès panafricains, création de l'Organisation de l'unité africaine en 1963, puis de l'Union africaine en 2002). Nous étudierons cet historique en détail dans la suite de la leçon.

\begin{popfigure}
\begin{center}
\begin{tikzpicture}[scale=1.0]
% Cercle panafricanisme
\draw[fill=popblueL,draw=popblue,thick,opacity=0.85] (-1.3,0) circle (2.4);
% Cercle consciencisme
\draw[fill=poporangeL,draw=poporange,thick,opacity=0.85] (1.3,0) circle (2.4);
% Étiquettes
\node[popblue,font=\bfseries] at (-2.6,1.5) {Panafricanisme};
\node[popblue,font=\small,align=center] at (-2.5,0.4) {mouvement\\historique global\\(politique, culturel,\\économique)};
\node[poporange,font=\bfseries] at (2.6,1.5) {Consciencisme};
\node[poporange,font=\small,align=center] at (2.5,0.4) {doctrine\\philosophique\\de Nkrumah\\(1964)};
\node[font=\small\bfseries,align=center] at (0,-0.7) {unité de\\l'Afrique,\\décolonisation,\\développement\\autonome};
\node[font=\footnotesize\itshape] at (0,-2.9) {Zone de recouvrement : le consciencisme se veut l'instrument philosophique du panafricanisme.};
\end{tikzpicture}
\end{center}
\legende{Diagramme de Venn : le panafricanisme (mouvement large) et le consciencisme (doctrine précise) se recouvrent sur l'objectif commun de l'unité et du développement de l'Afrique, sans se confondre.}
\end{popfigure}

\courssub{Distinction et lien entre les deux notions}

Le diagramme ci-dessus visualise l'essentiel : les deux notions se \emph{recouvrent} sans se \emph{confondre}. Distinguons-les rigoureusement.

\begin{center}
\begin{tabularx}{\linewidth}{|l|X|X|}\hline
\rowcolor{popblueL} \textbf{Critère} & \textbf{Panafricanisme} & \textbf{Consciencisme} \\ \hline
Nature & Mouvement historique, politique et culturel & Doctrine philosophique d'un auteur \\ \hline
Auteur & Collectif (Du Bois, Garvey, Nkrumah, etc.) & Kwame Nkrumah seul \\ \hline
Naissance & Fin du XIX\textsuperscript{e} siècle (diaspora) & 1964 (ouvrage \emph{Consciencism}) \\ \hline
Contenu & Unité et solidarité de l'Afrique et de la diaspora & Synthèse du matérialisme dialectique et des valeurs africaines \\ \hline
Portée & Large : plusieurs courants et générations & Précise : un système de pensée structuré \\ \hline
\end{tabularx}
\end{center}

\begin{attention}[Erreur fréquente à l'examen]
Ne confonds jamais le \textbf{panafricanisme} (mouvement historique global, porté par de nombreux acteurs sur plus d'un siècle) et le \textbf{consciencisme} (doctrine philosophique d'un seul auteur, Kwame Nkrumah, exposée en 1964). Écrire \og le panafricanisme de Nkrumah publié en 1964 \fg{} est une erreur : en 1964, Nkrumah publie le \emph{consciencisme}, qui s'inscrit \emph{dans} le courant panafricain sans s'y réduire.
\end{attention}

Le lien entre les deux notions est un lien d'\emph{instrument à projet}. Le panafricanisme définit un objectif : l'unité et la libération de l'Afrique. Mais un mouvement politique a besoin d'une théorie qui le fonde et l'oriente. C'est précisément le rôle que Nkrumah assigne au consciencisme : fournir au panafricanisme sa \emph{philosophie}, c'est-à-dire une analyse rationnelle de la situation africaine et une méthode d'action. C'est pourquoi il l'appelle \og l'arme philosophique du panafricanisme \fg{}. On peut dire, en schématisant : le panafricanisme dit \emph{ce qu'il faut faire} (unir l'Afrique) ; le consciencisme explique \emph{pourquoi et comment} le faire (en transformant les consciences et en construisant une société socialiste adaptée aux réalités africaines).

\begin{popfigure}
\begin{center}
\begin{tikzpicture}[scale=1.0]
% Cadre évoquant une photographie
\draw[thick,popdark,fill=popcream] (0,0) rectangle (8.6,5.4);
\draw[popline] (0.25,0.25) rectangle (8.35,5.15);
% Tribune et foule stylisées
\draw[fill=popmuted!60,draw=popdark] (1.0,0.6) rectangle (7.6,1.6);
\draw[fill=popmuted!40,draw=popdark] (1.0,1.6) rectangle (7.6,2.0);
% Micros
\draw[popdark,thick] (4.0,2.0) -- (4.0,2.7); \fill[popdark] (4.0,2.8) circle (0.12);
\draw[popdark,thick] (4.6,2.0) -- (4.6,2.6); \fill[popdark] (4.6,2.7) circle (0.12);
% Orateur bras levé
\fill[popdark] (4.3,3.35) circle (0.28);
\draw[popdark,line width=2pt] (4.3,3.1) -- (4.3,2.0);
\draw[popdark,line width=2pt] (4.3,2.9) -- (5.2,3.9);
\draw[popdark,line width=2pt] (4.3,2.9) -- (3.6,3.5);
% Drapeau
\draw[popdark,thick] (6.6,2.0) -- (6.6,4.6);
\draw[fill=popgold,draw=popdark] (6.6,4.6) -- (7.9,4.35) -- (6.6,4.1) -- cycle;
% Foule : têtes
\foreach \x in {1.4,1.9,2.4,2.9,5.6,6.1,7.0,7.4} {\fill[popmuted] (\x,0.95) circle (0.16);}
% Légende interne
\node[font=\footnotesize\itshape,align=center] at (4.3,0.35) {Accra, Ghana, 6 mars 1957 : minuit, la proclamation de l'indépendance.};
\end{tikzpicture}
\end{center}
\legende{Représentation schématique de la photographie célèbre de Kwame Nkrumah proclamant l'indépendance du Ghana le 6 mars 1957, à Accra, devant une foule immense : \og Notre indépendance n'a pas de sens si elle n'est pas liée à la libération totale de l'Afrique. \fg{} Le Ghana devient le premier pays d'Afrique subsaharienne colonisée à accéder à l'indépendance.}
\end{popfigure}

\courssub{Un exemple appliqué : lire une situation camerounaise avec ces notions}

\begin{exoresolu}[Le débat sur l'unité africaine]
\textbf{Énoncé.} Dans un débat en classe, un élève déclare : \og Le consciencisme, c'est le mouvement qui a créé l'Union africaine pour unir tous les Noirs du monde. \fg{} Cette affirmation contient deux erreurs. Identifie-les et corrige-les en rédigeant une définition exacte de chaque notion.
\tcblower
\textbf{Solution détaillée.}
\begin{enumerate}
\item \textbf{Première erreur : la nature du consciencisme.} Le consciencisme n'est pas un \emph{mouvement} mais une \emph{doctrine philosophique}, élaborée par un auteur précis, Kwame Nkrumah, dans \emph{Consciencism} (1964). Il adapte le matérialisme dialectique aux réalités africaines en y intégrant les valeurs traditionnelles du continent (communauté, égalité, humanisme).
\item \textbf{Deuxième erreur : la confusion des notions.} Le mouvement qui vise l'unité des peuples africains et de la diaspora noire, et dont l'Union africaine (2002, succédant à l'OUA créée en 1963) est une réalisation institutionnelle, c'est le \emph{panafricanisme}, né à la fin du XIX\textsuperscript{e} siècle dans la diaspora.
\item \textbf{Correction rédigée.} \og Le panafricanisme est un mouvement politique et culturel visant l'unité et la solidarité des peuples africains et de la diaspora noire ; il a abouti institutionnellement à l'OUA puis à l'Union africaine. Le consciencisme est la doctrine philosophique de Kwame Nkrumah (1964), qui se présente comme l'instrument théorique de ce mouvement : il fournit au panafricanisme une fondation rationnelle en synthétisant le matérialisme dialectique et les valeurs africaines. \fg{}
\end{enumerate}
\end{exoresolu}

\courssub{Méthode : bien définir une notion philosophique}

\begin{methode}[Définir une notion en quatre étapes]
Pour définir correctement une notion philosophique au Probatoire ou au Baccalauréat technique, procède ainsi :
\begin{enumerate}
\item \textbf{L'étymologie} : décompose le mot et donne son origine (ex. : \emph{pan-} = tout ; \emph{consciencism} = néologisme forgé à partir de \emph{conscience}).
\item \textbf{L'auteur ou le contexte} : précise qui a forgé ou porté la notion, quand et dans quelle situation historique (ex. : Nkrumah, 1964, contexte des indépendances).
\item \textbf{Le contenu} : énonce les traits essentiels de la notion, sans les confondre avec ceux d'une notion voisine.
\item \textbf{Un exemple concret} : illustre par une situation précise (ex. : la politique d'unité africaine de Nkrumah, le débat camerounais sur le franc CFA).
\end{enumerate}
Une définition qui réunit ces quatre éléments est complète et valorisée par le correcteur.
\end{methode}

\courssub{Annonce du plan de la leçon et organisation des séances}

Cette leçon se déroulera sur \textbf{deux séances}.

\textbf{Première séance.} Après cette introduction (problématique et définitions), nous étudierons le \emph{contexte historique} qui a vu naître ces idées : la colonisation, les indépendances et la quête d'identité africaine. Nous examinerons ensuite les \emph{fondements idéologiques} du consciencisme, puis ses \emph{fondements philosophiques} proprement dits.

\textbf{Deuxième séance.} Nous retracerons l'\emph{historique du panafricanisme}, des premiers congrès de la diaspora à l'Union africaine. Nous appliquerons ensuite ces notions à des situations concrètes : l'unité africaine et le développement, avec des exemples camerounais (intégration régionale, CEMAC, débats monétaires). Enfin, une séquence de \emph{méthode} vous entraînera à rédiger et justifier une démarche argumentative, en vue de la dissertation et du commentaire de texte.

\begin{aretenir}[L'essentiel de cette introduction]
\begin{itemize}
\item Le \cle{consciencisme} est la doctrine philosophique de Kwame Nkrumah (\emph{Consciencism}, 1964) : synthèse du matérialisme dialectique et des valeurs africaines, \og arme philosophique du panafricanisme \fg{}.
\item Le \cle{panafricanisme} est un mouvement politique et culturel, né dans la diaspora noire à la fin du XIX\textsuperscript{e} siècle, visant l'unité et la solidarité des peuples africains et de la diaspora.
\item Les deux notions se recouvrent (unité, décolonisation, développement) sans se confondre : l'une est une doctrine d'auteur, l'autre un mouvement historique global.
\item La problématique de la leçon : comment penser un développement authentiquement africain après la colonisation ?
\end{itemize}
\end{aretenir}

\coursec{Contexte historique : colonisation, indépendances et quête d'identité}

Dans la section précédente, nous avons défini le \cle{consciencisme} comme la philosophie du développement proposée par Kwame Nkrumah, et le \cle{panafricanisme} comme le mouvement d'unité des peuples africains et de la diaspora noire. Mais une philosophie ne naît jamais de rien : elle répond à une situation historique précise. Pour comprendre pourquoi des penseurs africains ont éprouvé le besoin de forger une philosophie de l'unité et du développement, il faut d'abord comprendre le traumatisme auquel ils répondaient : la traite négrière, la colonisation, la négation de l'humanité africaine. C'est ce contexte historique que nous allons maintenant reconstituer, en suivant une démarche d'historien : partir des faits, les organiser, puis dégager le problème philosophique qu'ils posent.

\courssub{La traite négrière et la naissance de la diaspora noire}

Commençons par une observation simple : pourquoi le panafricanisme est-il né \emph{hors} d'Afrique, dans les Amériques, avant de revenir sur le continent ? La réponse tient en un mot : la \cle{diaspora}.

\begin{definition}[Diaspora noire]
La \cle{diaspora} noire désigne l'ensemble des populations d'origine africaine dispersées à travers le monde (Amériques, Caraïbes, Europe) à la suite de la traite négrière. Du XV\ieme{} au XIX\ieme{} siècle, environ 12 à 15 millions d'Africains ont été déportés comme esclaves vers les plantations d'Amérique.
\end{definition}

L'intuition est la suivante : arrachés à des ethnies et des langues différentes, mélangés sur les plantations, les esclaves ont perdu leurs appartenances particulières (Yoruba, Bamiléké, Wolof...) pour ne plus être désignés que par une seule étiquette : « Noirs ». Paradoxalement, c'est donc l'oppression elle-même qui a créé la conscience d'une \emph{identité commune} : puisqu'on les opprimait tous comme Noirs, c'est comme Noirs qu'ils allaient se défendre. C'est pourquoi les premières idées panafricanistes apparaissent chez des intellectuels noirs américains et caribéens (W. E. B. Du Bois, Marcus Garvey) à la fin du XIX\ieme{} siècle : la diaspora est l'origine lointaine du panafricanisme.

\courssub{La colonisation européenne de l'Afrique}

\textbf{Le partage de Berlin (1884-1885) : l'Afrique découpée comme un gâteau}\par

À la fin du XIX\ieme{} siècle, les puissances européennes, en pleine révolution industrielle, cherchent des matières premières (caoutchouc, coton, minerais), des débouchés pour leurs produits et des territoires stratégiques. Pour éviter une guerre entre elles, elles se réunissent à Berlin, de novembre 1884 à février 1885, sous la direction du chancelier allemand Bismarck.

\begin{aretenir}[La conférence de Berlin (1884-1885)]
La conférence de Berlin réunit quatorze puissances européennes (plus les États-Unis en observateurs) pour organiser le partage de l'Afrique. \textbf{Aucun Africain n'y participe.} Elle fixe les règles de la colonisation : une puissance qui occupe effectivement un territoire côtier peut revendiquer l'arrière-pays. En moins de trente ans, toute l'Afrique, à l'exception de l'Éthiopie et du Libéria, passe sous domination européenne.
\end{aretenir}

Les conséquences sont considérables et se font encore sentir aujourd'hui :

\begin{itemize}
\item Des \textbf{frontières artificielles}, tracées à la règle sur des cartes, qui ignorent les peuples réels : un même groupe ethnique se retrouve séparé entre deux États (les Ewondo entre Cameroun et Gabon, les peuples peuls éclatés entre plusieurs pays).
\item Une \textbf{économie d'exploitation} : les colonies sont organisées pour extraire les richesses vers l'Europe (travail forcé, impôts écrasants, cultures imposées), non pour développer l'Afrique.
\item Une \textbf{domination politique et culturelle} : les langues, les administrations et les systèmes éducatifs européens sont imposés, tandis que les institutions africaines sont détruites ou méprisées.
\end{itemize}

\begin{popfigure}
\begin{center}
\begin{tikzpicture}[scale=0.9, every node/.style={font=\small}]
% Contour très simplifié de l'Afrique
\draw[thick, popdark, fill=popcream] (1.2,7.2) -- (3.2,7.6) -- (5.2,7.4) -- (6.6,6.6) -- (7.4,5.4) -- (7.0,4.2) -- (7.8,3.2) -- (7.2,1.8) -- (6.0,0.6) -- (4.6,0.2) -- (3.6,0.8) -- (2.8,2.0) -- (2.2,3.0) -- (1.0,3.4) -- (0.2,4.6) -- (0.6,6.0) -- cycle;
% Zones coloniales schématiques
\fill[popblue!35] (1.0,6.2) -- (4.6,7.2) -- (4.8,5.6) -- (3.0,4.8) -- (1.2,5.0) -- cycle;
\fill[poporange!35] (4.8,5.6) -- (6.4,6.4) -- (7.2,5.2) -- (6.8,4.0) -- (5.0,4.2) -- cycle;
\fill[popgreen!35] (3.0,4.8) -- (5.0,4.2) -- (5.2,2.6) -- (3.4,2.4) -- (2.4,3.4) -- cycle;
\fill[poppink!35] (3.4,2.4) -- (5.2,2.6) -- (5.8,1.0) -- (4.4,0.4) -- (3.6,1.0) -- cycle;
% Point Cameroun
\fill[popink] (3.6,4.4) circle (0.09);
\node[popink, font=\small\bfseries] at (4.9,4.7) {Cameroun (1960)};
% Légendes des zones
\node[popblue!80!black, font=\footnotesize] at (2.6,6.0) {Afrique de l'Ouest};
\node[popblue!80!black, font=\footnotesize] at (2.6,5.6) {et du Nord (France)};
\node[poporange!80!black, font=\footnotesize] at (6.0,5.2) {Est (GB)};
\node[popgreen!60!black, font=\footnotesize] at (3.9,3.5) {Centre (Fr., Belg.)};
\node[poppink!80!black, font=\footnotesize] at (4.5,1.6) {Sud (GB)};
% Flèche du nord
\draw[->, thick, popdark] (-0.6,5.6) -- (-0.6,6.8);
\node[popdark] at (-0.6,7.1) {N};
% Quelques dates d'indépendance
\node[font=\footnotesize, popdark] at (1.8,6.8) {1960};
\node[font=\footnotesize, popdark] at (6.6,3.4) {1963};
\node[font=\footnotesize, popdark] at (4.9,0.9) {1990};
\end{tikzpicture}
\end{center}
\legende{Carte muette schématique de l'Afrique coloniale (à compléter en TP) : 1. zone d'influence française ; 2. zone britannique ; 3. Afrique centrale (France, Belgique) ; 4. Afrique australe. Le point noir situe le Cameroun, indépendant le 1\ier{} janvier 1960. Contour indicatif, sans précision cartographique.}
\end{popfigure}

\courssub{Le cas du Cameroun : trois colonisations successives}

L'histoire coloniale du Cameroun est particulièrement riche, car le pays a connu \emph{trois} dominations étrangères successives :

\begin{enumerate}
\item \textbf{La colonisation allemande (1884-1916)} : le 12 juillet 1884, les chefs Duala (Bell, Akwa, Deïdo) signent avec les commerçants allemands un traité qui fait du « Kamerun » un protectorat allemand. Les Allemands imposent le travail forcé pour les plantations et les chantiers, ce qui provoque de nombreuses révoltes.
\item \textbf{La défaite allemande (1916) et le partage franco-britannique} : pendant la Première Guerre mondiale, les troupes françaises et britanniques chassent les Allemands. En 1919, la Société des Nations (SDN) partage le Cameroun : environ les quatre cinquièmes reviennent à la France (Cameroun oriental), le reste à la Grande-Bretagne (Cameroun occidental). En 1946, ces territoires deviennent des territoires sous \emph{tutelle} de l'ONU, avec l'obligation théorique de les conduire vers l'autonomie.
\item \textbf{Les indépendances (1960-1961)} : le Cameroun français accède à l'indépendance le \textbf{1\ier{} janvier 1960}. Le Cameroun britannique du Sud choisit par référendum (février 1961) de rejoindre la République du Cameroun : c'est la \textbf{réunification} du 1\ier{} octobre 1961, qui donne naissance à la République fédérale du Cameroun. Le Cameroun britannique du Nord, lui, rejoint le Nigeria.
\end{enumerate}

Ce cas camerounais illustre parfaitement le problème des frontières artificielles : le Cameroun actuel est né de découpages décidés à Berlin, à Versailles et à New York, jamais par les Camerounais eux-mêmes.

\courssub{La négation de l'humanité africaine : la provocation intellectuelle}

La colonisation ne s'est pas contentée de prendre des terres : elle a prétendu se justifier par une théorie. C'est là un point capital pour le philosophe. La pensée coloniale a élaboré des \emph{thèses de l'infériorité} des Noirs : les Africains seraient des êtres sans histoire, sans écriture, sans philosophie, incapables de se gouverner, que l'Europe aurait le devoir de « civiliser ». C'est ce qu'on a appelé la « mission civilisatrice ».

\begin{attention}[Un préjugé déguisé en science]
Les thèses coloniales de l'infériorité se présentaient comme « scientifiques » (hiérarchie des races, mesures des crânes...). Or elles reposent sur un raisonnement circulaire : on part du préjugé que le Noir est inférieur, puis on sélectionne les « faits » qui le confirment. Ce n'est pas de la science, c'est une \emph{idéologie} : un système d'idées construit pour justifier une domination. Le philosophe doit apprendre à démasquer ce type de raisonnement.
\end{attention}

Cette négation de l'humanité africaine fonctionne comme une \textbf{provocation intellectuelle} : elle oblige les intellectuels africains et noirs à répondre, à prouver — ou plutôt à affirmer — la valeur de leur civilisation. Toute la philosophie africaine du XX\ieme{} siècle naît de cette réplique.

\courssub{Les réponses culturelles : la négritude}

La première grande réponse est culturelle et littéraire : le mouvement de la \cle{négritude}, né à Paris dans les années 1930 autour des étudiants noirs Léopold Sédar Senghor (Sénégal), Aimé Césaire (Martinique) et Léon-Gontran Damas (Guyane).

\begin{definition}[Négritude]
La \cle{négritude} est le mouvement littéraire et philosophique qui affirme la valeur de la civilisation noire, face à sa négation par la pensée coloniale. Pour Senghor, la négritude est « l'ensemble des valeurs culturelles du monde noir » : la sensibilité, le rythme, la communion avec la nature, la vie communautaire. Césaire, lui, en fait une arme de combat : assumer fièrement son identité noire pour résister à l'assimilation.
\end{definition}

L'intuition est puissante : puisque le colonisateur disait « vous n'êtes rien », la négritude répond « nous sommes quelque chose, et ce quelque chose a sa propre beauté ». La formule célèbre de Senghor — « L'émotion est nègre, comme la raison est hellène » — valorise une voie africaine de connaissance par la sensibilité et la participation, différente de la raison analytique occidentale. On pourra discuter cette formule (risque-t-elle d'enfermer le Noir dans l'émotion ?), mais sa fonction historique est claire : \emph{restaurer la dignité} blessée par la colonisation. La négritude prépare ainsi le terrain philosophique du panafricanisme et du consciencisme : avant de s'unir et de se développer, l'Africain doit d'abord cesser de se mépriser lui-même.

\courssub{Les luttes de libération et les indépendances}

\textbf{Des révoltes aux partis nationalistes}\par

La décolonisation n'est pas un cadeau : elle est le résultat de luttes. Après la Seconde Guerre mondiale, plusieurs facteurs se combinent : les puissances coloniales sont affaiblies ; les soldats africains qui ont combattu pour la liberté de la France exigent la leur ; l'ONU proclame le droit des peuples à disposer d'eux-mêmes ; et surtout, des \textbf{partis nationalistes} africains organisent la revendication : le RDA en Afrique française, le Convention People's Party de Nkrumah au Ghana, et au Cameroun l'\cle{UPC} (Union des populations du Cameroun), fondée en 1948, qui réclame l'indépendance immédiate et la réunification des deux Cameroun.

\begin{savaistu}[Ruben Um Nyobè à la tribune de l'ONU]
Ruben Um Nyobè (1913-1958), secrétaire général de l'UPC, est le premier leader nationaliste camerounais à porter la question de l'indépendance du Cameroun devant l'Organisation des Nations unies, à New York, en \textbf{1952}. Il y dénonce la répression coloniale et réclame l'indépendance et la réunification. Pourchassé, il est tué par l'armée française le 13 septembre 1958, deux ans avant l'indépendance qu'il avait préparée. Longtemps effacé de la mémoire officielle, il est aujourd'hui reconnu comme un héros national.
\end{savaistu}

\begin{attention}[Erreur fréquente : les indépendances ne sont pas des « dons »]
Il est faux et dangereux de présenter les indépendances comme des cadeaux généreusement accordés par les colonisateurs. Elles résultent de \textbf{luttes menées par les Africains eux-mêmes} : syndicats, partis, manifestations, grèves, parfois luttes armées (Cameroun, Algérie, Kenya). Dire que « la France a donné l'indépendance », c'est effacer le sacrifice de militants comme Um Nyobè, et c'est reprendre à son compte l'idéologie coloniale de l'Africain passif. Dans une dissertation, montrez toujours l'activité des colonisés.
\end{attention}

\textbf{1960 : « l'année de l'Afrique »}\par

\begin{exemplebox}[Dix-sept indépendances en une seule année]
L'année \textbf{1960} est surnommée « l'année de l'Afrique » : dix-sept États africains accèdent à l'indépendance en quelques mois, dont le Nigeria, le Sénégal, la Côte d'Ivoire, le Congo, la Somalie, Madagascar... et le \textbf{Cameroun, le 1\ier{} janvier 1960}, qui ouvre la série. Le Ghana de Nkrumah les avait précédés dès 1957, première colonie d'Afrique noire à devenir indépendante. En 1963, les États indépendants créent l'Organisation de l'unité africaine (OUA) à Addis-Abeba.
\end{exemplebox}

\begin{popfigure}
\begin{center}
\begin{tikzpicture}[scale=1.0]
% Axe
\draw[->, thick, popdark] (0,0) -- (13.2,0);
% Graduations et événements
\foreach \x/\date/\txt/\col in {
  0.8/{1885}/{Conférence de Berlin :\\ partage de l'Afrique}/popink,
  3.6/{1945}/{5\ieme{} congrès panafricain\\ de Manchester}/popblue,
  6.2/{1957}/{Indépendance\\ du Ghana (Nkrumah)}/popgreen,
  8.8/{1960}/{« Année de l'Afrique » :\\ 17 indépendances, dont\\ le Cameroun (1\ier{} janv.)}/poporange,
  11.6/{1963}/{Création de l'OUA\\ à Addis-Abeba}/poppurple}
{
  \draw[\col, thick] (\x,0) -- (\x,0.25);
  \fill[\col] (\x,0) circle (0.09);
  \node[\col, font=\small\bfseries, above] at (\x,1.15) {\date};
  \node[font=\footnotesize, align=center, above, text=popdark] at (\x,0.28) {\txt};
}
\end{tikzpicture}
\end{center}
\legende{Frise chronologique : de la colonisation à l'unité africaine. 1. Berlin (1885) : partage colonial ; 2. Manchester (1945) : le panafricanisme exige l'indépendance ; 3. Ghana (1957) : première indépendance en Afrique noire ; 4. 1960 : « année de l'Afrique » ; 5. OUA (1963) : organisation de l'unité.}
\end{popfigure}

\courssub{Le problème posé aux nouveaux États : le néocolonialisme}

L'indépendance politique est conquise. Mais très vite, les nouveaux dirigeants découvrent une réalité amère : le drapeau a changé, pas les structures économiques. C'est le problème du \cle{néocolonialisme}.

\begin{definition}[Néocolonialisme]
Le \cle{néocolonialisme} désigne la domination économique et politique qui survit à l'indépendance formelle. L'État est officiellement souverain, mais son économie reste contrôlée de l'extérieur : monnaies arrimées aux anciennes métropoles (le franc CFA), entreprises étrangères détenant les matières premières, dettes, accords de « coopération » déséquilibrés, interventions politiques et militaires. Nkrumah le résume : c'est « l'État qui, en théorie, est indépendant et possède tous les signes extérieurs de la souveraineté, mais dont le système économique et la politique intérieure sont dirigés de l'extérieur ».
\end{definition}

À cela s'ajoutent deux autres handicaps hérités de la colonisation :

\begin{itemize}
\item Les \textbf{frontières artificielles}, sources de conflits entre États et de tensions entre groupes ethniques à l'intérieur de chaque État ;
\item Le \textbf{sous-développement} : l'économie coloniale n'a construit ni industrie, ni routes intérieures cohérentes, ni système éducatif de masse — seulement des axes d'évacuation des matières premières vers les ports.
\end{itemize}

\begin{exoresolu}[Analyser une situation concrète]
\textbf{Énoncé.} Un pays africain est indépendant depuis 1960. Sa monnaie est garantie par l'ancienne métropole, ses principales richesses (pétrole, cacao) sont exploitées par des sociétés étrangères qui rapatrient leurs bénéfices, et son budget dépend de l'aide extérieure. Peut-on dire que ce pays est pleinement souverain ? Justifiez votre réponse à l'aide de la notion de néocolonialisme.
\tcblower
\textbf{Solution.} Formellement, ce pays est souverain : il a un drapeau, un siège à l'ONU, un gouvernement élu. Mais l'analyse des faits montre que les trois leviers essentiels de la souveraineté lui échappent : (1) la \emph{monnaie}, garantie et contrôlée de l'extérieur ; (2) les \emph{ressources}, exploitées au profit d'entreprises étrangères ; (3) le \emph{budget}, dépendant de l'aide, donc des conditions posées par les bailleurs. On peut donc conclure, avec Nkrumah, qu'il s'agit d'une indépendance formelle masquant une dépendance réelle : c'est précisément la définition du \cle{néocolonialisme}. La souveraineté juridique existe, la souveraineté économique fait défaut.
\end{exoresolu}

\courssub{Naissance du besoin d'une philosophie africaine du développement}

Nous pouvons maintenant dégager la conclusion de cette enquête historique, en procédant par étapes :

\begin{enumerate}
\item La colonisation a laissé des États \emph{artificiels, divisés et dépendants} ;
\item La pensée coloniale a laissé une \emph{blessure identitaire} : le doute sur la valeur même de l'Africain ;
\item La négritude a répondu sur le plan culturel, mais la culture seule ne construit ni routes ni usines ;
\item Le néocolonialisme montre que l'indépendance politique sans indépendance économique est une coquille vide.
\end{enumerate}

Il faut donc une pensée qui soit à la fois une \emph{philosophie de l'identité} (rendre à l'Africain la conscience de sa dignité), une \emph{philosophie de l'unité} (surmonter les divisions héritées de Berlin) et une \emph{philosophie de l'action} (organiser le développement économique et social). Tel est précisément le \textbf{terreau du consciencisme} : la philosophie que Nkrumah va proposer comme « conscience de soi » collective de l'Afrique, capable de transformer l'indépendance formelle en libération réelle. C'est ce que nous étudierons dans la section suivante, consacrée aux fondements idéologiques du consciencisme.

\begin{exercice}[Pour s'entraîner]
1. Expliquez en quelques lignes pourquoi le panafricanisme est né dans la diaspora avant de revenir en Afrique.
2. Reconstituez les trois étapes de la domination étrangère au Cameroun, avec les dates.
3. Sujet de réflexion : « Les indépendances africaines ont-elles été données ou conquises ? » Rédigez un paragraphe argumenté (thèse, deux arguments, exemples camerounais).
\end{exercice}

\begin{corrige}[Exercice]
1. La traite négrière a dispersé des millions d'Africains dans les Amériques ; réduits tous au même statut de « Noirs » opprimés, ils ont pris conscience d'une identité commune et ont forgé les premières idées d'unité (Du Bois, Garvey), avant que le mouvement ne revienne sur le continent.
2. Colonisation allemande (1884-1916) ; partage franco-britannique sous mandat SDN puis tutelle ONU (1919-1960) ; indépendance du Cameroun français le 1\ier{} janvier 1960 et réunification avec le Cameroun britannique du Sud le 1\ier{} octobre 1961.
3. Thèse : les indépendances ont été conquises par la lutte. Argument 1 : l'action des partis nationalistes (UPC au Cameroun, CPP au Ghana) et les sacrifices de militants comme Um Nyobè, tué en 1958. Argument 2 : le contexte international (ONU, droit des peuples) n'a fait qu'entériner un rapport de forces créé par la mobilisation africaine. Conclusion : parler de « don » revient à reprendre le mythe colonial de l'Africain passif.
\end{corrige}

\coursec{Les fondements idéologiques du consciencisme}

La section précédente a montré comment la colonisation a brisé les structures traditionnelles et imposé une domination économique et culturelle, suscitant chez les élites africaines une double exigence : recouvrer la souveraineté politique et refonder une identité propre. C'est dans ce creuset que Kwame Nkrumah, premier président du Ghana indépendant (1957), élabore le \cle{consciencisme}, une philosophie d'action qui prétend offrir à l'Afrique les armes théoriques de sa libération totale. Nous analysons ici les piliers idéologiques de cette doctrine.

\courssub{L'héritage marxiste : matérialisme historique et critique de l'impérialisme}

Nkrumah ne cache pas sa dette à l'égard de Karl Marx. Pour lui, le matérialisme historique fournit l'outil d'analyse le plus rigoureux pour comprendre la société : ce ne sont pas les idées qui changent le monde, mais les conditions matérielles d'existence (mode de production, rapports de classe) qui déterminent la conscience. Il retient trois apports majeurs :

\begin{itemize}
    \item La \cle{lutte des classes} comme moteur de l'histoire : la société coloniale oppose le prolétariat et la paysannerie africains à la bourgeoisie compradore et au capital étranger.
    \item La \cle{critique de l'économie politique} : le capitalisme impérialiste pille les ressources africaines (matières premières, main-d'œuvre bon marché) pour accumuler du capital dans les métropoles.
    \item La \cle{nécessité d'une rupture révolutionnaire} : le réformisme est illusoire ; seul un renversement des structures de propriété permet l'émancipation.
\end{itemize}

\begin{attention}[Erreur fréquente : le consciencisme n'est pas une simple copie du marxisme]
Dire que Nkrumah « applique » le marxisme à l'Afrique est une contre-vérité pédagogique. Nkrumah \textbf{adapte} et \textbf{transforme} l'outil marxiste. Il refuse le \textbf{dogmatisme} (l'application mécanique de schémas européens) pour deux raisons :
\begin{enumerate}
    \item L'Afrique précoloniale ne connaît ni la féodalité ni le capitalisme industriel ; elle ignore l'antagonisme de classes tel que Marx le décrit en Europe.
    \item La priorité absolue n'est pas la lutte des classes \emph{interne}, mais la lutte \emph{anti-impérialiste} pour la souveraineté nationale et continentale.
\end{enumerate}
Le consciencisme est une \emph{synthèse}, non une copie.
\end{attention}

\courssub{L'adaptation aux réalités africaines : le communalisme comme base du socialisme}

Si Nkrumah emprunte à Marx l'analyse de l'exploitation capitaliste, il puise dans les \cle{traditions africaines} le modèle d'une société sans classes, sans exploitation, fondée sur l'entraide. C'est la thèse centrale de son ouvrage \emph{Consciencisme} (1964) : le socialisme africain ne s'importe pas, il \emph{revient} aux sources.

\begin{definition}[Communalisme africain]
Organisation sociale traditionnelle fondée sur la \textbf{propriété collective de la terre} (la terre appartient aux ancêtres, aux vivants et aux générations futures ; nul ne peut l'aliéner) et sur la \textbf{solidarité du groupe} (le clan, la lignée, le village). L'individu ne s'accomplit que dans la communauté : « Je suis parce que nous sommes » (formule associée à la philosophie bantoue de l'\emph{ubuntu}).
\end{definition}

Pour Nkrumah, ce communalisme est l'équivalent fonctionnel du « communisme primitif » décrit par Engels, mais il a survécu là où il a disparu en Europe. Il constitue donc la \textbf{base matérielle et morale} du socialisme africain moderne. La tâche n'est pas de « construire » le socialisme \emph{ex nihilo}, mais de \textbf{moderniser} le communalisme traditionnel par la science, la technique et la planification d'État.

\begin{savaistu}[Le parcours intellectuel de Kwame Nkrumah (1909--1972)]
Né à Nkroful, en Gold Coast britannique (actuel Ghana), Nkrumah étudie à l'Achimota College, puis part aux \textbf{États-Unis} (1935--1945) où il obtient des diplômes de philosophie, sociologie et éducation à Lincoln University et à l'Université de Pennsylvanie. Il y découvre le marxisme, le panafricanisme (Du Bois, Garvey) et la pensée de Descartes, Kant, Hegel. Il poursuit à \textbf{Londres} (1945--1947), au cœur des réseaux anticoloniaux (George Padmore, Jomo Kenyatta). C'est cette double formation — occidentale et panafricaine — qui forge le consciencisme.
\end{savaistu}

\courssub{Anti-impérialisme, anticolonialisme et la thèse du néocolonialisme}

Le consciencisme est d'abord une \cle{philosophie de la libération}. Nkrumah distingue trois stades de la domination :
\begin{enumerate}
    \item Le \textbf{colonialisme classique} : occupation militaire, administration directe, exploitation ouverte (avant 1957 pour le Ghana).
    \item L'\textbf{indépendance formelle} : drapeau, hymne, siège à l'ONU, mais économie extravertie, monnaie liée au franc ou à la livre, accords de défense.
    \item Le \textbf{néocolonialisme} : « \emph{dernière étape de l'impérialisme} » (titre de son ouvrage majeur, \emph{Neo-Colonialism: The Last Stage of Imperialism}, 1965). L'État est nominalement souverain, mais son économie et sa politique sont téléguidées par des puissances étrangères via les multinationales, l'aide conditionnée, la dette, les coups d'État.
\end{enumerate}

\begin{aretenir}[Thèse centrale du consciencisme sur le néocolonialisme]
Le néocolonialisme est plus dangereux que le colonialisme classique car il donne l'illusion de la liberté tout en maintenant la servitude économique. Il ne se combat pas seulement par des discours, mais par : (1) l'unité politique africaine (État fédéral continental), (2) la planification socialiste, (3) le non-alignement actif.
\end{aretenir}

\courssub{L'unité politique comme condition \emph{sine qua non} du développement}

La formule célèbre de Nkrumah — « \emph{Seek ye first the political kingdom and all else shall be added unto you} » (Cherchez d'abord le royaume politique et tout le reste vous sera donné par surcroît) — résume sa stratégie. Pour lui, \textbf{aucun développement économique n'est possible dans des États balkanisés, dépendants, aux marchés trop étroits}.

\begin{itemize}
    \item \textbf{Balkanisation} : l'Afrique compte (en 1963) plus de 30 États indépendants, héritage du partage de Berlin (1884--1885). Frontières artificielles, micro-marchés, duplication des infrastructures.
    \item \textbf{Solution} : un \textbf{Gouvernement continental unique} (État fédéral africain) avec une armée commune, une monnaie unique, une planification centrale, une politique étrangère unique.
    \item \textbf{Rôle de l'État} : l'État n'est pas un « veilleur de nuit » libéral ; c'est le \textbf{moteur du développement}. Il planifie, industrialise, éduque, assure la santé. Le secteur privé (national ou étranger) est subordonné au plan national.
\end{itemize}

\courssub{Illustration concrète : la planification industrielle et le barrage d'Akosombo}

Au Ghana, Nkrumah met en œuvre cette vision dès 1957 : plans de développement, création d'entreprises d'État (Black Star Line, Ghana Airways), éducation gratuite, et projet phare : le \textbf{barrage hydroélectrique d'Akosombo} sur la Volta (1961--1965). Objectif : produire l'énergie nécessaire à l'industrialisation (aluminerie de Tema exploitant la bauxite locale) et exporter l'excédent vers les pays voisins (Togo, Bénin, Côte d'Ivoire).

\begin{exemplebox}[Le barrage d'Akosombo : symbole du développement planifié]
Inauguré en janvier 1966 (premiers groupes mis en service dès 1965), le barrage d'Akosombo (lac Volta, l'un des plus grands lacs artificiels du monde), d'une puissance installée initiale de \num{588} MW (portée à \num{912} MW en 1972, puis à \num{1020} MW), a permis au Ghana de :
\begin{itemize}
    \item Alimenter l'aluminerie de Tema (VALCO), transformant localement la bauxite.
    \item Électrifier les villes et de nombreuses zones rurales, faisant fortement progresser le taux d'accès à l'électricité.
    \item Devenir \textbf{exportateur d'électricité} vers le Togo et le Dahomey (actuel Bénin), puis la Côte d'Ivoire, générant des devises.
\end{itemize}
C'est la concrétisation technique de la thèse nkrumahiste : \emph{État fort + planification + énergie = base de la souveraineté économique}.
\end{exemplebox}

\begin{popfigure}
\begin{tikzpicture}
\begin{axis}[
    width=\linewidth,
    height=8cm,
    title={Évolution de la production d'électricité au Ghana (1960--1970)},
    xlabel={Année},
    ylabel={Production (GWh)},
    xmin=1960, xmax=1970,
    ymin=0, ymax=3500,
    xtick={1960,1961,1962,1963,1964,1965,1966,1967,1968,1969,1970},
    ytick={0,500,1000,1500,2000,2500,3000,3500},
    grid=major,
    grid style={popline},
    legend pos=north west,
    legend style={draw=popdark, fill=popcream},
    mark size=2.5pt,
    thick
]
\addplot[popblue, mark=*] coordinates {
    (1960, 350)
    (1961, 400)
    (1962, 480)
    (1963, 550)
    (1964, 700)
    (1965, 1850)
    (1966, 2100)
    (1967, 2300)
    (1968, 2450)
    (1969, 2600)
    (1970, 2800)
};
\addlegendentry{Production totale (GWh)}
\draw[poporange, dashed, thick] (axis cs:1965,0) -- (axis cs:1965,3500) node[above, pos=0.9, rotate=90, anchor=south, font=\small, poporange] {Mise en service d'Akosombo (1965)};
\end{axis}
\end{tikzpicture}
\legende{Source : données indicatives reconstituées d'après la Volta River Authority (VRA) et la Banque mondiale. La rupture de pente en 1965 illustre l'impact immédiat du barrage sur la capacité productive nationale.}
\end{popfigure}

\courssub{Les sources théoriques du consciencisme : une synthèse originale}

Le consciencisme ne se réduit pas au couple Marx/Tradition. Nkrumah revendique explicitement trois sources qu'il synthétise dans une « philosophie matérialiste non-dogmatique » :

\begin{popfigure}
\begin{tikzpicture}[
    node distance=2.5cm and 3cm,
    source/.style={draw=popdark, thick, fill=popcream, rounded corners=8pt, minimum width=3.2cm, minimum height=1.8cm, align=center, font=\small},
    central/.style={draw=popblue, very thick, fill=popblue!10, ellipse, minimum width=4cm, minimum height=2.2cm, align=center, font=\normalsize\bfseries},
    arrow/.style={-Stealth, thick, popdark},
    labelarrow/.style={midway, fill=popcream, rounded corners=2pt, inner sep=2pt, font=\footnotesize, text=popdark}
]
\node[central, align=center] (C){CONSCIENCISME\\(Philosophie de la libération africaine)};

\node[source, align=center] (M) at (-5, 2.5){
    \textbf{MARXISME}\\
    (Matérialisme dialectique)\\
    \textit{Apports :}\\
    -- Analyse de classes\\
    -- Critique du capitalisme\\
    -- Nécessité révolutionnaire
};
\draw[arrow] (M) -- node[labelarrow, above left, align=center]{Outil d'analyse\\ structurelle} (C);

\node[source, align=center] (T) at (-5, -2.5){
    \textbf{TRADITIONS AFRICAINES}\\
    (Communalisme / Humanisme)\\
    \textit{Apports :}\\
    -- Propriété collective\\
    -- Solidarité organique\\
    -- Égalitarisme de base
};
\draw[arrow] (T) -- node[labelarrow, below left, align=center]{Base morale\\ et sociale} (C);

\node[source, align=center] (H) at (5, 0){
    \textbf{PENSÉE HUMANISTE OCCIDENTALE}\\
    (Descartes, Kant, Hegel,\\ Philosophie positive)\\
    \textit{Apports :}\\
    -- Raison critique\\
    -- Idée de progrès\\
    -- Méthode scientifique
};
\draw[arrow] (H) -- node[labelarrow, right, align=center]{Méthode\\ rationnelle} (C);

\draw[popmuted, dashed, thin] (M) -- node[labelarrow, left, popmuted] {Critique de l'aliénation} (T);
\draw[popmuted, dashed, thin] (M) -- node[labelarrow, above right, popmuted] {Dialectique} (H);
\draw[popmuted, dashed, thin] (T) -- node[labelarrow, below right, popmuted] {Universel humain} (H);
\end{tikzpicture}
\legende{Les trois sources du consciencisme selon Nkrumah. Les flèches pleines indiquent l'apport direct à la synthèse ; les flèches pointillées les résonances entre sources.}
\end{popfigure}

\courssub{Dialogue avec les autres idéologies de libération}

Le consciencisme ne vit pas en vase clos. Il dialogue, s'oppose ou converge avec trois grands courants de l'époque :

\begin{center}
\begin{tabularx}{\linewidth}{|l|X|X|X|}\hline
\rowcolor{popblueL}
\textbf{Idéologie} & \textbf{Point de convergence} & \textbf{Point de divergence / Apport spécifique du consciencisme} & \textbf{Position de Nkrumah} \\ \hline
\textbf{Nationalisme} & Refus de la domination étrangère ; fierté identitaire. & Nationalisme étroit = balkanisation, obstacle à l'unité continentale. & \textbf{Nationalisme panafricain} : la nation authentique, c'est l'Afrique. \\ \hline
\textbf{Socialisme (marxiste-léniniste / socialiste africain)} & Propriété collective des moyens de production ; planification. & Socialisme « scientifique » importé vs socialisme « endogène » issu du communalisme. & \textbf{Socialisme africain} : adaptation, non imitation. \\ \hline
\textbf{Non-alignement} (Bandung 1955, Belgrade 1961) & Refus des blocs Est/Ouest ; souveraineté ; paix. & Non-alignement passif vs \textbf{non-alignement actif} : l'Afrique doit peser, proposer, fédérer. & \textbf{Leader du non-alignement radical} : l'unité africaine comme levier mondial. \\ \hline
\end{tabularx}
\end{center}

\courssub{Analyse critique : les limites internes du consciencisme au pouvoir}

Une doctrine se juge aussi à ses résultats politiques. Le régime de Nkrumah (1957--1966) illustre les tensions du consciencisme mis en pratique.

\begin{itemize}
    \item \textbf{Dérive autoritaire} : loi sur la détention préventive (1958), interdiction des partis d'opposition, proclamation de la République à parti unique (1964), culte de la personnalité (\emph{Osagyefo}, « le Rédempteur »).
    \item \textbf{Échec économique relatif} : endettement massif pour financer l'industrialisation (barrage, usines), chute du cours du cacao (ressource principale), corruption d'une partie de la nouvelle élite bureaucratique.
    \item \textbf{Rupture avec la base communaliste} : l'État-parti s'impose \emph{sur} la société au lieu d'en émaner. Le « communalisme moderne » devient dirigisme bureaucratique.
    \item \textbf{Coup d'État du 24 février 1966} : l'armée et la police renversent le régime pendant que Nkrumah est en voyage à Hanoï. Nkrumah y voit la preuve tragique de sa thèse du néocolonialisme (il accuse des puissances extérieures) ; on peut aussi y lire la sanction de l'autoritarisme interne.
\end{itemize}

\begin{methode}[Analyser une doctrine politique : grille de lecture en 4 étapes]
Pour ne pas réduire une pensée à ses caricatures (ni l'idéaliser, ni la diaboliser), appliquez cette démarche :
\begin{enumerate}
    \item \textbf{Identifier les sources} : quelles traditions intellectuelles l'auteur revendique-t-il ? Quelles en sont les ruptures ? (Ex. : Nkrumah assume Marx mais refuse le déterminisme de classes.)
    \item \textbf{Formuler les thèses centrales} : quel est le diagnostic (néocolonialisme) ? Quelle est la thérapie (unité politique + socialisme communaliste) ?
    \item \textbf{Étudier les applications concrètes} : politiques publiques, institutions créées, indicateurs chiffrés (éducation, énergie, PIB, dette, libertés).
    \item \textbf{Évaluer avec esprit critique} : confrontation entre les promesses et les résultats. Distinguer les causes \emph{externes} (néocolonialisme, conjoncture) et \emph{internes} (choix politiques, erreurs de gestion, dérive autoritaire).
\end{enumerate}
\end{methode}

\begin{exoresolu}[Analyse d'un extrait : Nkrumah, \emph{Consciencisme}, 1964]
\textbf{Énoncé :} « La société africaine traditionnelle était égalitaire. Le socialisme africain est le retour à cette égalité, mais sur une base technique et scientifique moderne. »
\textbf{Question :} En quoi cette phrase résume-t-elle la thèse du « communalisme comme base du socialisme » ?
\tcblower
\textbf{Solution détaillée :}
\begin{enumerate}
    \item \textbf{Identification de la thèse} : Nkrumah affirme une \textbf{continuité historique} : le socialisme n'est pas une importation étrangère, c'est la \emph{modernisation} d'un héritage endogène.
    \item \textbf{Analyse des concepts} :
        \begin{itemize}
            \item « Société traditionnelle égalitaire » = communalisme (propriété collective, absence de classes antagonistes).
            \item « Retour » = légitimation culturelle (contre l'accusation d'aliénation idéologique).
            \item « Base technique et scientifique moderne » = rupture avec le traditionalisme figé ; nécessité de l'industrialisation, de la planification, de l'État.
        \end{itemize}
    \item \textbf{Portée idéologique} : cette phrase permet de \textbf{réconcilier} la paysannerie (attachée aux valeurs communautaires) et l'élite modernisatrice (formée à la science occidentale). Elle fonde la \textbf{légitimité interne} du pouvoir nkrumahiste.
    \item \textbf{Limite} : l'égalitarisme traditionnel reposait sur la \textbf{subsistance} et l'autorité des aînés ; le socialisme moderne implique \textbf{productivité industrielle, division technique du travail, État centralisé}. La transition n'est pas sans friction (ex. : expropriation des terres pour le barrage d'Akosombo, environ \num{80000} personnes déplacées).
\end{enumerate}
\end{exoresolu}

\begin{aretenir}[Bilan synthétique : les fondements idéologiques du consciencisme]
\begin{itemize}
    \item \textbf{Sources} : marxisme (analyse), communalisme africain (base éthique et sociale), rationalisme occidental (méthode).
    \item \textbf{Thèses fortes} : néocolonialisme = stade suprême de l'impérialisme ; unité politique continentale = condition du développement ; socialisme africain = communalisme modernisé par l'État planificateur.
    \item \textbf{Réalisations} : Akosombo (énergie, industrie), éducation de masse, diplomatie panafricaine (OUA, 1963).
    \item \textbf{Limites critiques} : dérive autoritaire (parti unique, détention préventive), endettement, déconnexion État/société, renversement en 1966.
    \item \textbf{Héritage} : malgré l'échec politique immédiat, le consciencisme reste la \textbf{matrice théorique} du panafricanisme politique et de l'idée d'intégration continentale (Union africaine, ZLECAf aujourd'hui).
\end{itemize}
\end{aretenir}

\coursec{Les fondements philosophiques du consciencisme}

\courssub{Transition : de l'idéologie à la philosophie}

Dans la section précédente, nous avons analysé les \textbf{fondements idéologiques} du consciencisme : la volonté politique de construire une société africaine libre, unie et socialiste, rompant avec l'héritage colonial. Mais une idéologie, pour être opératoire et durable, ne saurait reposer sur la seule volonté ou l'émotion ; elle requiert une \textbf{assise théorique rigoureuse}, une vision du monde qui explique la réalité et guide l'action. C'est le rôle de la philosophie. Pour Nkrumah, « la philosophie est l'instrument de l'idéologie » : elle en fournit la justification rationnelle, la cohérence interne et la méthode d'analyse. Passons donc du \emph{quoi} (le projet politique) au \emph{pourquoi} et au \emph{comment} (la structure ontologique et épistémologique du consciencisme).

\begin{savaistu}[Nkrumah, philosophe de l'action]
Kwame Nkrumah (1909--1972) n'est pas seulement un homme d'État ; il a étudié la philosophie à Lincoln University (USA) et à la London School of Economics. Il a fréquenté les cercles marxistes et pan-africains londoniens (George Padmore, C.L.R. James). Son ouvrage majeur, \emph{Consciencism : Philosophy and Ideology for De-colonisation} (1964), est une tentative rare en Afrique : construire une philosophie politique \emph{autochtone} utilisant les outils conceptuels du matérialisme dialectique pour penser la spécificité africaine.
\end{savaistu}

\textbf{Le matérialisme philosophique : la primauté de la matière}\par

\textbf{Intuition :} Pourquoi certaines idées dominent-elles une époque ? Pourquoi l'idée de « mission civilisatrice » a-t-elle justifié la colonisation ? Pour le matérialisme, ce ne sont pas les idées qui changent le monde par magie ; ce sont les \textbf{conditions matérielles d'existence} (mode de production, rapports de classe, environnement géographique, niveau technique) qui façonnent les consciences.

\begin{definition}[Matérialisme philosophique]
Doctrine selon laquelle la \textbf{matière} (la réalité objective, existante en dehors et indépendamment de la conscience) est la \textbf{réalité première}, et que la conscience, la pensée, les idées sont le \textbf{reflet} (produit secondaire, superstructure) de cette matière organisée (le cerveau humain, la société). Il n'y a pas d'esprit pur désincarné ; les idées naissent des conditions matérielles d'existence.
\end{definition}

\textbf{Application au consciencisme :} Nkrumah rejette l'idéalisme (hégélien ou religieux) qui ferait de l'« esprit africain » ou de la « négritude » une essence métaphysique, immuable, antérieure à l'histoire. Pour lui, la \cle{personnalité africaine} n'est pas une donnée ontologique figée, mais une \textbf{construction historique} issue des conditions matérielles de vie des peuples africains (économie de subsistance, structures lignagères, traumatisme de la traite, économie de plantation coloniale, lutte pour l'indépendance).

\begin{exemplebox}[Exemple concret : la propriété foncière au Cameroun]
Dans la société traditionnelle bassa ou bamiléké, la terre appartient à la lignée/au clan (propriété collective d'usage). L'idée dominante est : « La terre ne se vend pas, elle se transmet. » La colonisation introduit le Code foncier (1974 au Cameroun, inspiré du droit français) : la terre devient une \emph{marchandise} (titre foncier, vente, hypothèque). \textbf{Le changement matériel (mode d'appropriation de la terre) a fait naître de nouvelles idées juridiques et de nouveaux conflits sociaux.} C'est la thèse matérialiste : l'être (économie foncière) détermine la conscience (droit, représentation).
\end{exemplebox}

\begin{attention}[Piège : le déterminisme mécanique]
Le matérialisme de Nkrumah (hérité de Marx/Engels/Lénine) n'est pas un \emph{réductionnisme économique} grossier (« tout est économie »). La conscience a une \textbf{relative autonomie} : elle réagit sur la matière, elle peut accélérer ou freiner l'histoire. C'est tout le sens de la \emph{dialectique} (voir ci-dessous). Ne confondez pas « primauté de la matière » et « toute-puissance de l'économie à chaque instant ».
\end{attention}

\textbf{Le matérialisme dialectique : le mouvement par contradictions}\par

\textbf{Intuition :} La réalité n'est pas statique. Une société, une cellule, une étoile changent sans cesse. Comment expliquer le changement sans faire appel à une cause externe (Dieu, Destin, Grand Homme) ? Par la \textbf{contradiction interne}.

\begin{definition}[Matérialisme dialectique]
Conception du monde selon laquelle le développement de la matière (nature, société, pensée) s'effectue par la \textbf{lutte des contraires} (contradictions internes) et leur \textbf{dépassement} (résolution). Schéma classique : \textbf{Thèse} (état initial) $\rightarrow$ \textbf{Antithèse} (négation, opposition interne) $\rightarrow$ \textbf{Synthèse} (dépassement qui conserve les acquis et résout la contradiction), laquelle devient une nouvelle thèse.
\end{definition}

\begin{propriete}[Les trois lois de la dialectique (Engels), applicables à l'histoire africaine]
\begin{enumerate}
    \item \textbf{Loi de l'unité et de la lutte des contraires :} Tout phénomène contient des pôles opposés (ex. : travail vs capital, tradition vs modernité, unité vs division) dont l'affrontement est le moteur du mouvement.
    \item \textbf{Loi du passage des changements quantitatifs aux changements qualitatifs :} L'accumulation lente de réformes, de prises de conscience, de forces productives (quantitatif) mène brutalement à une rupture : révolution, indépendance, changement de régime (qualitatif).
    \item \textbf{Loi de la négation de la négation :} Le développement ne revient pas au point de départ (cercle), mais progresse en spirale : on retrouve des formes anciennes (ex. : solidarité communautaire) à un niveau supérieur (ex. : sécurité sociale nationale, coopératives modernes).
\end{enumerate}
\end{propriete}

\begin{experience}[Démarche d'investigation : analyser une contradiction camerounaise]
\textbf{Objet :} La contradiction entre \emph{économie informelle} (survie, débrouillardise, 90\% de l'emploi) et \emph{économie formelle} (État, grandes entreprises, droit du travail, fiscalité).
\textbf{Protocole :}
\begin{enumerate}
    \item Identifier les deux pôles (Thèse : secteur formel structurant mais exclusif ; Antithèse : secteur informel inclusif mais précaire).
    \item Observer leur lutte : répression (contrôles fiscaux, déguerpissements) vs résistance (évasion fiscale, syndicalisme informel).
    \item Chercher la Synthèse possible : politique d'inclusion (statut d'auto-entrepreneur simplifié, accès au crédit, protection sociale universelle) qui \emph{négative la précarité} sans \emph{négative la vitalité} de l'informel.
\end{enumerate}
\textbf{Interprétation :} Le développement camerounais ne viendra pas de l'éradication de l'informel (utopie), mais de sa \textbf{transformation dialectique} en secteur productif protégé. C'est une application concrète du matérialisme dialectique à la politique publique.
\end{experience}

\textbf{Les trois segments de la société africaine : thèse, antithèse, synthèse}\par

C'est le \textbf{cœur original} du consciencisme. Nkrumah applique la dialectique à la structure même de la conscience africaine contemporaine. La société africaine n'est pas un bloc homogène ; elle est traversée par trois \textbf{segments} (courants culturels, religieux, philosophiques) qui s'affrontent et doivent se synthétiser.

\begin{popfigure}
\begin{tikzpicture}[
    node distance=2.5cm and 3cm,
    segment/.style={draw, thick, rounded corners, minimum width=3.5cm, minimum height=2cm, align=center, font=\small, fill=popcream},
    arrow/.style={->, thick, >=Stealth, popblue},
    synth/.style={draw, very thick, rounded corners, minimum width=4cm, minimum height=1.8cm, align=center, font=\small\bfseries, fill=popgreenL, text=popdark},
    label/.style={font=\tiny, text=popmuted, anchor=north}
]
% Nœuds des trois segments
\node[segment, align=center] (trad) at (0,0){
    \textbf{Segment Traditionnel}\\
    (Africanité originelle)\\
    \textit{Valeurs :} Communautarisme, ancêtres,\\
    oralité, terre sacrée, consensus
};
\node[segment, align=center] (islam) at (7,0){
    \textbf{Segment Islamique}\\
    (Apport arabo-musulman)\\
    \textit{Valeurs :} Unité (Oumma), écriture,\\
    droit (Charia), discipline,\\
    commerce transsaharien
};
\node[segment, align=center] (euro) at (3.5,-3.5){
    \textbf{Segment Euro-chrétien}\\
    (Héritage colonial/moderne)\\
    \textit{Valeurs :} Technique, science,\\
    droit positif, individu,\\
    État-nation, christianisme
};
% Flèches de convergence
\draw[arrow, bend left=20] (trad) to node[above, label] {Lutte / Négation} (islam);
\draw[arrow, bend right=20] (islam) to node[below, label] {Lutte / Négation} (trad);
\draw[arrow] (trad) to node[left, label, pos=0.6] {Contradiction} (euro);
\draw[arrow] (islam) to node[right, label, pos=0.6] {Contradiction} (euro);
% Synthèse au centre
\node[synth, align=center] (synth) at (3.5,-1.2){
    \textbf{CONSCIENCE NOUVELLE}\\
    (Consciencisme)\\
    \textit{Synthèse dialectique :}\\
    Humanisme scientifique +\\
    Socialisme africain +\\
    Unité panafricaine
};
% Flèches vers la synthèse
\draw[arrow, dashed, poporange] (trad.north east) -- (synth.south west);
\draw[arrow, dashed, poporange] (islam.north west) -- (synth.south east);
\draw[arrow, dashed, poporange] (euro.north) -- (synth.south);
% Légendes des apports conservés (Aufheben)
\node[label, align=center] at (0, -1.5) {Conserve :\ Solidarité,\ Spiritualité};
\node[label, align=center] at (7, -1.5) {Conserve :\ Discipline,\ Universalisme};
\node[label, align=center] at (3.5, -4.8) {Conserve :\ Science,\ Technique,\ Droit};
\end{tikzpicture}
\legende{Schéma dialectique des trois segments de la société africaine selon Nkrumah. Les flèches pleines indiquent la contradiction/lutte ; les flèches pointillées indiquent l'apport conservé (Aufheben) dans la synthèse : la Conscience nouvelle (Consciencisme).}
\end{popfigure}

\textbf{Analyse des segments :}\par
\begin{itemize}
    \item \textbf{Segment traditionnel (Thèse) :} L'africanité pré-coloniale. Base matérielle : économie de subsistance, propriété collective. Valeur clé : \cle{humanisme communautaire} (« Je suis parce que nous sommes » -- \emph{Ubuntu}). \textbf{Limite :} particularisme ethnique, technicité faible, vulnérabilité face à l'extérieur.
    \item \textbf{Segment islamique (Antithèse partielle) :} Apport par le Nord (Sahélien) et l'Est (Swahili). Base matérielle : commerce longue distance, écriture, État centralisé (empires du Mali, Songhaï, Bornou, Sokoto). Valeur clé : \cle{universalisme} (Oumma transcende la tribu), discipline, livre. \textbf{Contradiction avec le traditionnel :} individualisme religieux vs communautarisme lignager ; écriture vs oralité.
    \item \textbf{Segment euro-chrétien (Antithèse majeure / Négation de la thèse) :} Imposé par la colonisation (esclavage, traite, administration, école missionnaire). Base matérielle : capitalisme industriel, extraction, État bureaucratique. Valeur clé : \cle{technicité scientifique}, droit formel, individu juridique. \textbf{Contradiction violente :} destruction des structures traditionnelles, aliénation culturelle, racisme.
\end{itemize}

\textbf{La Synthèse (Consciencisme) :} Elle ne consiste pas à « choisir » un segment ou à faire un « mélange » éclectique. C'est un \textbf{dépassement dialectique (Aufheben)} : on \emph{conserve} (solidarité, universalisme, science), on \emph{supprime} (particularisme ethnique, obscurantisme, exploitation capitaliste), on \emph{élève} au niveau supérieur : une \textbf{conscience révolutionnaire panafricaine, humaniste et scientifique}.

\begin{aretenir}[Les trois segments et la synthèse]
\begin{itemize}
    \item \textbf{Traditionnel} : Base communautaire (Thèse).
    \item \textbf{Islamique} : Médiation universaliste (Antithèse interne).
    \item \textbf{Euro-chrétien} : Choc modernisateur/aliénant (Antithèse externe).
    \item \textbf{Synthèse (Consciencisme)} = Humanisme scientifique + Socialisme africain + Panafricanisme politique.
\end{itemize}
Le développement africain authentique ne peut venir que de cette synthèse assumée, pas du mimétisme occidental ni du repli identitaire.
\end{aretenir}

\textbf{La conscience révolutionnaire : moteur de l'histoire}\par

\textbf{Du psychologique à l'historique :} En langage courant, « conscience » = connaissance immédiate de soi (psychologie). Chez Nkrumah (et Marx), la \cle{conscience révolutionnaire} est \textbf{collective, de classe, historique}. C'est le passage de la \emph{classe en soi} (position objective dans la production) à la \emph{classe pour soi} (organisation politique consciente de ses intérêts antagonistes).

\begin{definition}[Conscience révolutionnaire (selon Nkrumah)]
Prise de conscience \textbf{collective} par les masses populaires (paysans, ouvriers, intellectuels engagés) de leur \textbf{situation objective d'oppression et d'aliénation} (néocolonialisme, dépendance, sous-développement) et de la \textbf{nécessité historique} de la transformer par l'action politique organisée (parti de masse, État révolutionnaire).
\end{definition}

\begin{exemplebox}[Exemple : La grève des planteurs de cacao au Cameroun (années 1950) / UPC]
Les planteurs (classe en soi : producteurs de rente coloniale, exploités par les compagnies d'achat) prennent conscience (via l'UPC, syndicats) que leur pauvreté n'est pas une fatalité naturelle mais le fruit du \textbf{système colonial} (prix imposés, monopole d'exportation). Ils s'organisent en \textbf{parti politique} (UPC) pour conquérir l'indépendance \emph{réelle} (économique). C'est l'émergence d'une conscience révolutionnaire : de la plainte individuelle à l'action collective structurée.
\end{exemplebox}

\begin{attention}[Erreur fréquente : « Conscience » $\neq$ « Opinion » ou « Sentiment »]
Au Probatoire / Baccalauréat technique, ne définissez jamais la conscience révolutionnaire nkrumahienne comme « le fait d'être fier d'être Africain » (sentiment/négritude) ou « l'opinion que la colonisation est mauvaise » (jugement moral). C'est une \textbf{connaissance scientifique des lois sociales} (matérialisme dialectique) permettant une \textbf{pratique politique efficace}. « La pratique sans la pensée est aveugle ; la pensée sans la pratique est vide. »
\end{attention}

\textbf{L'humanisme : la personne humaine comme valeur suprême}\par

\textbf{Fondement ontologique :} Si la matière est première, l'être humain est \emph{la matière devenue consciente d'elle-même}. Il n'est pas un moyen (main-d'œuvre, chair à canon, « ressource humaine »), il est la \textbf{fin}. Le consciencisme fonde l'éthique sur l'ontologie : l'exploitation de l'homme par l'homme est une \textbf{négation de la nature humaine} (qui est sociale, créative, libre).

\begin{propriete}[Humanisme conscienciste vs Humanisme bourgeois]
\begin{center}
\begin{tabularx}{\linewidth}{|l|X|X|}\hline
\textbf{Dimension} & \textbf{Humanisme bourgeois (libéral)} & \textbf{Humanisme conscienciste (nkrumahien)} \\ \hline
\textbf{Sujet} & Individu abstrait, propriétaire, concurrent & Personne concrète, membre de la communauté, solidaire \\ \hline
\textbf{Égalité} & Formelle (droit), masque inégalités réelles & Réelle (économique, sociale, culturelle) : « À chacun selon ses besoins » \\ \hline
\textbf{Lutte contre le racisme} & Moral/juridique (anti-discrimination) & Structurel : destruction du système capitaliste/impérialiste qui \emph{produit} le racisme \\ \hline
\textbf{Développement} & Croissance PIB, investissement étranger & Épanouissement de la personne humaine (santé, éducation, culture, pouvoir) \\ \hline
\end{tabularx}
\end{center}
\end{propriete}

\textbf{Contre le racisme colonial :} Le racisme n'est pas une « erreur d'appréciation » ou un « préjugé » ; c'est une \textbf{idéologie fonctionnelle} justifiant l'exploitation (esclavage, travail forcé, apartheid, néocolonialisme). L'humanisme conscienciste le combat par la \textbf{preuve matérielle} : l'égalité des potentialités intellectuelles et créatrices démontrée par l'histoire africaine (empires, métallurgie, mathématiques, résistance) et par la construction socialiste présente.

\begin{savaistu}[L'égalité originelle chez Nkrumah]
Nkrumah écrit : « La philosophie du consciencisme repose sur la conviction que tous les hommes sont égaux par nature et que toute inégalité est le produit de conditions sociales historiques. » Cela renvoie à la \emph{Déclaration des droits de l'homme} (1789) mais en radicalisant : l'égalité n'est pas un « droit naturel » métaphysique, c'est un \textbf{fait biologique et anthropologique} que l'histoire (esclavagisme, colonialisme, capitalisme) a nié et qu'il faut \textbf{réaliser concrètement} par la révolution sociale.
\end{savaistu}

\textbf{La philosophie comme instrument d'action : le lien théorie-pratique}\par

C'est la citation la plus célèbre de \emph{Consciencism} : \og \textbf{La pratique sans la pensée est aveugle ; la pensée sans la pratique est vide.} \fg

\textbf{Analyse :}\par
\begin{itemize}
    \item \textbf{Pratique sans pensée = Aveugle :} Action spontanée, émeute, putsch militaire sans projet de société, activismisme stérile. Risque : manipulation par des forces réactionnaires, échec, dictature.
    \item \textbf{Pensée sans pratique = Vide :} Intellectualisme stérile, philosophie de salon, « négritude » purement esthétique ou académique sans prise sur le réel. Risque : complicité objective avec l'ordre établi (néocolonialisme).
    \item \textbf{Unité théorie-pratique :} La philosophie (matérialisme dialectique) fournit la \textbf{boussole} (analyse de la conjoncture, identification des forces motrices, stratégie). L'action révolutionnaire (masses organisées) fournit le \textbf{terrain de vérification} et d'enrichissement de la théorie.
\end{itemize}

\begin{methode}[Articuler Philosophie et Idéologie dans une dissertation]
\textbf{Sujet type :} « L'idéologie panafricaine a-t-elle besoin d'un fondement philosophique ? »
\begin{enumerate}
    \item \textbf{Distinguer :} Idéologie = représentation du monde au service d'un groupe/classe, mobilisatrice, partisane, orientée vers l'action immédiate. Philosophie = recherche rationnelle, critique, systématique des principes premiers (ontologie, épistémologie, éthique), visée de vérité.
    \item \textbf{Montrer le lien (Nkrumah) :} L'idéologie (panafricanisme, socialisme africain) \emph{a besoin} de la philosophie (consciencisme) pour : (a) ne pas être dogmatique ; (b) analyser scientifiquement la réalité africaine (trois segments) ; (c) élaborer une stratégie cohérente (unité politique $\leftarrow$ unité économique $\leftarrow$ base philosophique commune).
    \item \textbf{Nuancer :} La philosophie ne \emph{produit} pas l'idéologie (c'est la lutte des classes qui la produit), elle l'\emph{éclaire} et la \emph{critique}. Une philosophie coupée des masses devient idéalisme.
    \item \textbf{Conclure :} Oui, fondement indispensable pour passer du « panafricanisme sentimental » (Congrès de 1900--1945) au « panafricanisme scientifique » (OUA 1963, UA 2002, ZLECAf 2021).
\end{enumerate}
\end{methode}

\textbf{Étude guidée d'un extrait de \emph{Consciencism} (1964)}\par

\textbf{Texte (adapté, § 12--14) :}
\og La philosophie d'une société n'est pas un luxe intellectuel ; elle est l'expression théorique de sa structure sociale et économique. [...] La société africaine traditionnelle était une société communautaire ; sa philosophie était un humanisme communautaire. L'impact de la société islamique a introduit un élément d'universalisme. L'impact de la société euro-chrétienne a introduit la technique et la méthode scientifique, mais dans un cadre capitaliste et individualiste. La tâche du consciencisme est de synthétiser ces trois courants en une nouvelle philosophie : un humanisme scientifique, base du socialisme africain. \fg

\begin{exoresolu}[Commentaire de texte structuré]
\textbf{1. Situer :} Extrait de \emph{Consciencism} (1964), chapitre sur les « Fondements philosophiques ». Auteur : Kwame Nkrumah, président du Ghana, théoricien du panafricanisme politique. Contexte : Indépendance fraîche (1957), quête d'idéologie propre face au capitalisme et au marxisme-léninisme orthodoxe.
\textbf{2. Dégager la thèse :} La philosophie africaine authentique (consciencisme) doit être une \textbf{synthèse dialectique} des trois segments historiques (traditionnel, islamique, euro-chrétien) pour produire un \textbf{humanisme scientifique} fondant le \textbf{socialisme africain}.
\textbf{3. Analyser les arguments :}
\begin{itemize}
    \item \textit{Arg. ontologique :} « Philosophie = expression de la structure socio-économique » $\rightarrow$ Thèse matérialiste (base détermine superstructure).
    \item \textit{Arg. historique :} Identification des trois segments $\rightarrow$ Preuve empirique de la complexité africaine (contre essentialisme).
    \item \textit{Arg. dialectique :} « Synthétiser » $\rightarrow$ Dépassement (Aufheben) : garder technique/science (euro), universalisme (islam), communauté (tradition) ; rejeter capitalisme, particularisme, obscurantisme.
    \item \textit{Arg. politique final :} Visée = Socialisme africain. La philosophie n'est pas fin en soi, elle est \emph{pour} l'action.
\end{itemize}
\textbf{4. Discuter (Ouverture critique) :}
\begin{itemize}
    \item \textit{Force :} Refus du mimétisme ; prise en compte de l'islam (souvent oublié) ; lien explicite théorie/pratique.
    \item \textit{Limites :} Risque de \emph{volontarisme} : la synthèse est-elle réelle ou décrétée par l'État-parti ? Le « socialisme africain » a-t-il résolu la contradiction classe/ethnie ? L'expérience ghanéenne (CPP, dictature, coup d'État 1966) montre la difficulté de l'application.
    \item \textit{Actualité :} La ZLECAf (Zone de libre-échange continentale) et l'Agenda 2063 de l'UA sont-elles la concrétisation tardive de cet humanisme scientifique ?
\end{itemize}
\end{exoresolu}

\textbf{Arbre conceptuel : de la matière à l'action politique}\par

\begin{popfigure}
\begin{tikzpicture}[
    mindmap,
    grow cyclic,
    every node/.style={concept, execute at begin node=\hskip0pt},
    root concept/.append style={concept color=popblue, minimum size=3.5cm, font=\bfseries\large, text=white},
    level 1 concept/.append style={level distance=4.5cm, sibling angle=90, concept color=popblue!70!poppurple, minimum size=2.8cm, font=\small\bfseries, text=white},
    level 2 concept/.append style={level distance=3.5cm, sibling angle=45, concept color=popgreen!60!popblue, minimum size=2.2cm, font=\small, text=white},
    level 3 concept/.append style={level distance=3cm, sibling angle=30, concept color=poporange!80!popdark, minimum size=2cm, font=\scriptsize, text=white},
]
\node[root concept, align=center] (root){Matérialisme\\ Dialectique}
    child[clockwise from=90] {
        node[level 1 concept, align=center] (matiere){Matière première\\ (Être social)}
        child { node[level 2 concept, align=center]{Mode de production\\ (Base économique)} }
        child { node[level 2 concept, align=center]{Structure de classes\\ (Rapports sociaux)} }
    }
    child[clockwise from=180] {
        node[level 1 concept, align=center] (contra){Contradictions\\ (Moteur)}
        child { node[level 2 concept, align=center] (seg){Trois segments\\ (Trad / Islam / Euro)} }
        child { node[level 2 concept, align=center]{Néocolonialisme\\ vs Souveraineté} }
        child { node[level 2 concept, align=center]{Communautarisme\\ vs Modernité} }
    }
    child[clockwise from=270] {
        node[level 1 concept, align=center] (conscience){Conscience\\ Révolutionnaire}
        child { node[level 2 concept, align=center]{De la classe\\ en soi $\to$ pour soi} }
        child { node[level 2 concept, align=center]{Humanisme\\ scientifique} }
        child { node[level 2 concept, align=center]{Unité panafricaine\\ (Volonté politique)} }
    }
    child[clockwise from=0] {
        node[level 1 concept, align=center] (action){Action Politique\\ (Pratique)}
        child { node[level 2 concept, align=center]{Socialisme africain\\ (Objectif)} }
        child { node[level 2 concept, align=center]{Parti de masse / État\\ (Outil)} }
        child { node[level 2 concept, align=center]{ZLECAf / UA 2063\\ (Réalisations actuelles)} }
    };
\end{tikzpicture}
\legende{Arbre conceptuel du consciencisme : le matérialisme dialectique (racine) structure l'analyse des contradictions africaines (segments), fonde la conscience révolutionnaire (tronc), qui guide l'action politique panafricaine (feuilles). Les flèches orange indiquent les relations de détermination et d'éclairage.}
\end{popfigure}

\begin{aretenir}[Synthèse des fondements philosophiques]
\begin{enumerate}
    \item \textbf{Ontologie :} Matérialisme (primauté des conditions d'existence matérielles).
    \item \textbf{Méthode :} Dialectique (contradictions, dépassement, synthèse des trois segments).
    \item \textbf{Anthropologie :} Humanisme (personne humaine = fin suprême ; égalité originelle).
    \item \textbf{Epistémologie :} Unité théorie-pratique (philosophie = arme de la révolution).
    \item \textbf{Projet :} Socialisme africain + Panafricanisme = Consciencisme.
\end{enumerate}
\end{aretenir}

\begin{exercice}[Entraînement Probatoire / Baccalauréat technique - Dissertation]
\textbf{Sujet :} « La synthèse des trois segments de la société africaine est-elle la condition nécessaire du développement de l'Afrique ? »
\textit{Indications :} Définir les trois segments (thèse/antithèses). Montrer pourquoi le statu quo (fragmentation) bloque le développement (dépendance, conflits identitaires, gaspillage). Analyser la synthèse (consciencisme) comme dépassement dialectique. Discuter : condition \emph{nécessaire} mais est-elle \emph{suffisante} ? (Rôle de la direction politique, du contexte international, de la base productive).
\end{exercice}

\begin{corrige}[Exercice n°11 - Éléments de corrigé]
\textbf{Introduction :} Accroche (Nkrumah, \emph{Consciencism}, 1964). Définition des trois segments. Problématique : la fragmentation culturelle est-elle l'obstacle principal au développement ou le symptôme d'une dépendance économique ? Annonce du plan dialectique.
\textbf{I. Le diagnostic : la fragmentation (trois segments) comme frein structurel au développement.}
A. Héritage colonial : frontières artificielles, économies extraverties, élites aliénées (segment euro-chrétien dominant).
B. Poids des particularismes : ethnicisme politique (segment traditionnel instrumentalisé), tensions inter-religieuses (segment islamique vs chrétien).
C. Conséquence : impossibilité d'une politique industrielle cohérente, de la mutualisation des ressources, de la formation d'un marché intérieur unifié.
\textbf{II. La synthèse conscienciste : condition nécessaire d'un développement endogène.}
A. Dépassement dialectique (Aufheben) : récupérer la solidarité (tradition), l'universalisme (islam), la science (euro) $\rightarrow$ Humanisme scientifique.
B. Base du socialisme africain : planification démocratique, propriété sociale des moyens de production, priorité aux besoins de masse.
C. Base du panafricanisme : marché commun (ZLECAf), monnaie unique (Eco), défense commune, diplomatie unifiée. Sans cette unité philosophique/culturelle, l'unité politique reste formelle (OUA d'hier).
\textbf{III. Condition nécessaire mais non suffisante : les limites du volontarisme philosophique.}
A. La philosophie ne remplace pas l'économie : la synthèse idéologique ne crée pas d'usines, de routes, d'écoles sans lutte des classes concrète et investissement productif.
B. Risque de « philosophisme d'État » : imposer la synthèse par le haut (parti unique, répression) $\rightarrow$ étouffement de la créativité populaire, coups d'État (Ghana 1966).
C. Déterminants externes : ordre impérialiste (FMI, dette, ingérences) qui structure la dépendance bien au-delà des représentations culturelles.
\textbf{Conclusion :} La synthèse des segments est le \emph{préalable subjectif} (conscience) indispensable pour mobiliser les masses et orienter l'État, mais le développement effectif requiert la transformation \emph{objective} de la base productive (industrialisation, souveraineté alimentaire, rupture néocoloniale). Philosophie et économie dialectiquement liées.
\end{corrige}

\coursec{Historique du panafricanisme}

Après avoir analysé les fondements philosophiques du consciencisme — cette synthèse nkrumahiste qui articule matérialisme dialectique, humanisme traditionnel et esprit de révolution —, nous devons maintenant comprendre comment le panafricanisme, comme mouvement historique et politique, a préparé le terrain intellectuel et organisationnel à de telles élaborations. Le consciencisme ne naît pas dans le vide : il est la théorisation philosophique d'une lutte politique séculaire pour l'unité et la libération de l'Afrique. Retracer l'historique du panafricanisme, c'est donc suivre le fil conducteur qui relie les premières revendications de la diaspora à la création de l'Union africaine, en passant par les congrès, les indépendances et les débats stratégiques qui ont façonné l'Afrique contemporaine.

\courssub{Origines diasporiques et premiers pionniers (1900–1914)}

\begin{definition}[Diaspora africaine]
Ensemble des populations d'origine africaine dispersées dans le monde par la traite négrière (XVe–XIXe siècles) et les migrations ultérieures. La diaspora africaine compte aujourd'hui plus de 200 millions de personnes, principalement dans les Amériques, les Caraïbes et l'Europe. C'est dans ce creuset que naît la conscience panafricaine : l'expérience partagée de l'esclavage, du racisme et de la colonisation forge une identité commune au-delà des frontières.
\end{definition}

Le panafricanisme ne naît pas sur le continent africain, mais dans la \cle{diaspora}. Dès 1900, Henry Sylvester Williams (Trinité) organise à Londres la \textbf{première Conférence panafricaine} : 37 délégués, rédaction d'une \emph{Adresse aux nations du monde} dénonçant le racisme et le partage de l'Afrique (Conférence de Berlin, 1884–1885). Des intellectuels comme Edward Blyden (Îles Vierges), Anna Julia Cooper (États-Unis) et W.E.B. Du Bois (qui y participe) posent les jalons d'une solidarité transatlantique.

\medskip
\textbf{Marcus Garvey et l'UNIA (1914).} Né en Jamaïque, Marcus Garvey fonde en 1914 l'\emph{Universal Negro Improvement Association} (UNIA) à Kingston, puis la déplace à Harlem (New York). Son mot d'ordre : \og Africa for the Africans \fg. Garvey prône la fierté noire (\emph{Black Pride}), l'autonomie économique (Black Star Line, usines, commerces) et le \emph{Back to Africa} — le retour physique des Afro-descendants sur la terre ancestrale. En 1920, la convention de l'UNIA à New York rassemble des milliers de délégués et adopte la \emph{Déclaration des droits des peuples nègres du monde}. Bien que son projet de retour massif échoue (opposition coloniale, faillite de la Black Star Line), Garvey inscrit durablement l'idée d'une \cle{nation africaine globale} dans l'imaginaire politique.

\begin{attention}[Erreur fréquente : situer l'origine du panafricanisme en Afrique]
Ne confondez pas le \emph{lieu d'émergence} (diaspora) et le \emph{lieu d'ancrage final} (continent). Le panafricanisme naît \textbf{hors d'Afrique}, porté par des descendants d'esclaves qui n'ont jamais connu le continent (conférence de 1900 à Londres, UNIA 1914 à Kingston/Harlem). Ce n'est qu'avec le 5e congrès de Manchester (1945) que le centre de gravité bascule vers l'Afrique, porté par des leaders africains formés en Occident (Nkrumah, Kenyatta, Padmore). Au Probatoire, cette distinction chronologique et géographique est souvent testée.
\end{attention}

\courssub{L'ère des congrès panafricains : Du Bois et l'institutionnalisation (1919–1945)}

\medskip
\textbf{W.E.B. Du Bois, « père du panafricanisme organisé ».} Sociologue, historien, militant (premier Afro-Américain doctoré à Harvard), Du Bois comprend que la libération de l'Afrique conditionne l'égalité des Noirs partout. Il organise une série de \emph{congrès panafricains} pour porter la cause africaine sur la scène internationale :

\begin{itemize}
    \item \textbf{1919, Paris} (en marge de la Conférence de la Paix) : 57 délégués, 16 pays. Revendication : administration internationale des colonies allemandes, droits politiques pour les Africains.
    \item \textbf{1921, Londres / Bruxelles / Paris} : 110 délégués. Insistance sur l'éducation, la santé, le travail forcé.
    \item \textbf{1923, Londres / Lisbonne} : focus sur le mandat de la SDN.
    \item \textbf{1927, New York} : dernier congrès du cycle Du Bois ; tensions avec Garvey (Du Bois : élitisme intellectuel, intégration ; Garvey : nationalisme de masse, séparation).
\end{itemize}

Ces congrès restent largement diasporiques et intellectuels. Le tournant survient en 1945.

\begin{exemplebox}[Le 5e congrès panafricain de Manchester (octobre 1945)]
Environ \textbf{200 délégués} se réunissent à la Chorlton Town Hall. Pour la première fois, une majorité d'Africains du continent y participent : Kwame Nkrumah (Gold Coast), Jomo Kenyatta (Kenya), Hastings Banda (Nyasaland), Obafemi Awolowo (Nigeria), Peter Abrahams (Afrique du Sud). George Padmore (Trinité) en est l'organisateur principal ; Du Bois, 77 ans, préside honorifiquement.

\textbf{Résolutions clés :}
\begin{itemize}
    \item Revendication de l'\textbf{indépendance politique immédiate} (plus de pétitions, de l'action).
    \item Appel à la lutte des masses : syndicats, partis politiques, action directe.
    \item Affirmation : \og Les peuples des colonies doivent lutter pour leur liberté par tous les moyens à leur disposition \fg.
\end{itemize}
Ce congrès marque le passage du \emph{panafricanisme de pétition} au \emph{panafricanisme de libération}. Nkrumah et Kenyatta rentrent chez eux pour mener la lutte anticoloniale.
\end{exemplebox}

\courssub{Le tournant des indépendances : le Ghana comme base arrière (1957–1958)}

L'indépendance du \textbf{Ghana} (6 mars 1957), première colonie d'Afrique subsaharienne à accéder à la souveraineté, change la donne. Nkrumah, Premier ministre puis président, proclame : \og L'indépendance du Ghana n'a de sens que si elle est liée à la libération totale de l'Afrique \fg. Accra devient la capitale morale du panafricanisme.

Deux conférences majeures s'y tiennent en 1958 :
\begin{itemize}
    \item \textbf{Conférence des États africains indépendants} (avril 1958) : 8 pays (Ghana, Égypte, Éthiopie, Libye, Maroc, Tunisie, Liberia, Soudan). Création d'un \emph{Bureau de libération africain} pour soutenir les mouvements de guérilla (ALN en Algérie, MPLA en Angola, etc.).
    \item \textbf{Conférence de tous les peuples africains} (décembre 1958) : 300 délégués de 28 territoires (partis politiques, syndicats, mouvements de jeunesse). Résolution : \emph{indépendance immédiate et inconditionnelle} ; soutien à la lutte armée si nécessaire. Frantz Fanon y participe.
\end{itemize}

Nkrumah finance, forme et héberge les mouvements de libération. Le panafricanisme devient \textbf{politique d'État}.

\courssub{Les divisions stratégiques : Casablanca contre Monrovia (1961)}

Dès 1960–1961, l'enthousiasme de l'indépendance révèle deux visions incompatibles de l'unité africaine.

\begin{popfigure}
\begin{tikzpicture}[scale=0.85]
% Contour très simplifié de l'Afrique (polygone approximatif)
\draw[popdark, thick] (0,0) -- (10,0) -- (11,2) -- (12,5) -- (11,8) -- (9,10) -- (7,11) -- (4,11) -- (2,9) -- (0,6) -- cycle;
% Groupe de Casablanca (radicaux) - zones colorées
\foreach \x/\y/\name in {
    2.5/3.5/Ghana, 3.5/2.5/Guinée, 4.5/2.0/Mali, 6.5/2.5/Égypte, 7.5/4.5/Algérie, 8.5/7.5/Maroc, 5.5/5.5/Libye} {
    \fill[popred!30] (\x,\y) circle (0.6);
    \node[popdark, font=\tiny, align=center] at (\x,\y-0.9) {\name};
}
% Groupe de Monrovia (modérés) - zones colorées (noms 1961)
\foreach \x/\y/\name in {
    5.5/1.0/Nigeria, 6.5/0.5/Nigeria, 7.0/1.5/Cameroun, 8.0/2.0/Côte d'Ivoire, 8.5/3.0/Sénégal, 9.0/4.0/Sierra Leone, 7.5/6.0/Congo-Léop., 6.0/7.0/Tanganyika, 4.0/8.0/Éthiopie, 3.0/7.5/Soudan, 2.0/6.0/Tchad, 1.5/4.5/Liberia, 3.0/3.0/Togo, 4.0/2.0/Dahomey, 5.0/3.0/Niger, 6.0/4.0/RCA, 7.0/5.0/Gabon, 8.0/6.0/Congo-Brazza, 9.0/7.0/Somalie} {
    \fill[popblue!30] (\x,\y) circle (0.5);
    \node[popdark, font=\tiny, align=center] at (\x,\y-0.8) {\name};
}
% Légende
\node[popdark, font=\small\bfseries] at (13,10) {Groupe de Casablanca};
\fill[popred!30] (13,9.5) circle (0.15);
\node[popdark, font=\small] at (13.5,9.5) {Unité politique immédiate};
\node[popdark, font=\small\bfseries] at (13,8.5) {Groupe de Monrovia};
\fill[popblue!30] (13,8) circle (0.15);
\node[popdark, font=\small] at (13.5,8) {Coopération progressive};
% Flèche Nord
\draw[popdark, ->, thick] (13,6) -- (13,6.8) node[above, font=\tiny] {N};
\end{tikzpicture}
\legende{Carte schématique de l'Afrique (1961) : localisation approximative des pays membres du Groupe de Casablanca (rouge) et du Groupe de Monrovia (bleu). Les noms sont ceux de 1961 (ex. Dahomey, Congo-Léopoldville, Tanganyika). Les positions sont indicatives.}
\end{popfigure}

\begin{center}
\begin{tabularx}{\linewidth}{|l|X|X|}
\hline
\rowcolor{popblueL} \textbf{Critère} & \textbf{Groupe de Casablanca (radicaux)} & \textbf{Groupe de Monrovia (modérés)} \\ \hline
\textbf{Pays fondateurs (1961)} & Ghana, Guinée, Mali, Égypte, Algérie (GPRA), Maroc, Libye & Nigeria, Liberia, Sierra Leone, Sénégal, Côte d'Ivoire, Cameroun, Togo, Dahomey, Niger, Haute-Volta, Tchad, RCA, Congo-Brazzaville, Gabon, Congo-Léopoldville, Éthiopie, Somalie, Tanganyika, etc. \\ \hline
\textbf{Idée d'unité} & \textbf{État fédéral panafricain} immédiat : gouvernement commun, armée unique, monnaie unique, citoyenneté africaine. & \textbf{Coopération interétatique} souple : respect souveraineté, intégration économique progressive, organisations sectorielles. \\ \hline
\textbf{Moyens} & Rupture avec néocolonialisme, soutien actif mouvements libération, non-alignement radical. & Maintien liens avec anciens métropoles (Françafrique, Commonwealth), développement national prioritaire. \\ \hline
\textbf{Figure emblématique} & Kwame Nkrumah (\og Africa must unite \fg, 1963) & Léopold Sédar Senghor, Félix Houphouët-Boigny, Abubakar Tafawa Balewa. \\ \hline
\end{tabularx}
\end{center}

Cette opposition paralyse la première conférence des chefs d'État à Addis-Abeba (mai 1963). Le compromis naît de la médiation de l'empereur Hailé Sélassié et du Togo (Sylvanus Olympio).

\courssub{De l'OUA à l'Union africaine : le compromis institutionnel (1963–2002)}

\begin{savaistu}[Le 25 mai, Journée de l'Afrique]
Le \textbf{25 mai 1963}, 32 chefs d'État africains signent la Charte de l'\emph{Organisation de l'Unité Africaine (OUA)} à Addis-Abeba. Cette date est célébrée chaque année comme la \textbf{Journée de l'Afrique} (ex-Journée de la libération). L'OUA consacre le compromis :
\begin{itemize}
    \item \textbf{Principe intangible} : respect des frontières héritées de la colonisation (\emph{uti possidetis juris}) — victoire de Monrovia.
    \item \textbf{Objectif affiché} : unité politique progressive, libération totale (Afrique du Sud, Rhodésie, colonies portugaises) — victoire de Casablanca.
    \item \textbf{Non-ingérence} : article III, al. 2 — frein à la défense des droits de l'homme et à la gestion des conflits internes.
\end{itemize}
L'OUA réussit la décolonisation (soutien aux mouvements de libération, fonds de libération) mais échoue sur l'intégration économique et la prévention des conflits (guerres civiles, génocides).
\end{savaistu}

\medskip
\textbf{De l'OUA à l'UA (2002) : le tournant de Sirte.} En 1999, à Sirte (Libye), Kadhafi pousse à la transformation. L'\emph{Acte constitutif de l'Union africaine} (Lomé, 2000 ; entrée en vigueur 2002) rompt avec l'OUA sur trois points majeurs :
\begin{enumerate}
    \item \textbf{Droit d'ingérence} : l'UA peut intervenir dans un État membre pour crimes de guerre, génocide, crimes contre l'humanité (article 4h).
    \item \textbf{Institutions renforcées} : Assemblée, Conseil exécutif, Commission, Parlement panafricain, Cour de justice, Conseil de paix et de sécurité.
    \item \textbf{Intégration économique} : marche vers la \emph{Communauté économique africaine} (traité d'Abuja 1991) : union douanière, marché commun, union monétaire, banque centrale africaine.
\end{enumerate}

\medskip
\textbf{La ZLECAf (2021) : le projet économique phare.} La \emph{Zone de libre-échange continentale africaine} (AfCFTA en anglais), lancée officiellement le 1er janvier 2021, vise à créer un marché unique de 1,4 milliard de consommateurs, PIB cumulé ~3 400 milliards USD. Objectifs : éliminer 90\% des droits de douane intra-africains, faciliter la circulation des personnes (protocole sur la libre circulation, 2018), développer les chaînes de valeur régionales. C'est l'aboutissement concret — encore inachevé — du rêve d'intégration économique de la conférence d'Abuja (1991) et des pères fondateurs.

\courssub{Le panafricanisme culturel : négritude, festivals et littérature}

Le panafricanisme n'est pas seulement politique ; il est aussi \textbf{culturel}. La \emph{négritude} (Aimé Césaire, Léopold Sédar Senghor, Léon-Gontran Damas, années 1930–1940) affirme la valeur civilisationnelle des cultures noires face à l'assimilation coloniale. \og La négritude n'est pas une philosophie, c'est une poésie \fg (Senghor) — mais elle fonde une esthétique et une éthique de la reconnaissance.

\begin{itemize}
    \item \textbf{Présence africaine} (revue, 1947–1963, Paris) : lieu de débat entre intellectuels africains et caribéens.
    \item \textbf{Premier Festival mondial des arts nègres, Dakar 1966} : initié par Senghor, 45 pays, 2 000 participants. Affirmation : \og L'art nègre est la manifestation la plus haute de l'âme nègre \fg. Contre-festival d'Alger 1969 (Festival panafricain d'Alger) : art engagé, révolutionnaire, refus de l'esthétisme.
    \item \textbf{Littérature africaine d'expression française/anglaise/portugaise} : Cheikh Hamidou Kane (\emph{L'Aventure ambiguë}), Chinua Achebe (\emph{Things Fall Apart}), Mongo Beti, Ferdinand Oyono, Pepetela, Mia Couto... Ils réécrivent l'histoire, déconstruisent le colonialisme, inventent des formes narratives hybrides.
\end{itemize}

Cette dimension culturelle nourrit le \emph{consciencisme} : Nkrumah y puise la \cle{philosophie de la personnalité africaine} — ni retour au passé, ni mimétisme occidental, mais synthèse créatrice.

\courssub{Méthode : construire une frise chronologique pour mémoriser l'histoire du panafricanisme}

\begin{methode}[Mémoriser une chronologie complexe par périodes]
Pour maîtriser l'historique du panafricanisme (exigé au Probatoire : dissertation, commentaire, questions de cours), ne mémorisez pas une liste brute de dates. \textbf{Regroupez les événements en grandes périodes logiques} et construisez une frise visuelle.
\begin{enumerate}
    \item \textbf{Identifiez les périodes-clés} (ici : 1. Diaspora \& pionniers ; 2. Congrès \& institutionnalisation ; 3. Indépendances \& base arrière ; 4. Divisions \& OUA ; 5. UA \& intégration ; 6. Culturel).
    \item \textbf{Sélectionnez 1 à 2 événements repères par période} (date + nom + enjeu principal).
    \item \textbf{Tracez l'axe horizontal} (flèche du temps) ; placez les périodes comme des blocs colorés.
    \item \textbf{Ajoutez les flèches de causalité} : ex. Manchester 1945 $\rightarrow$ Indépendance Ghana 1957 $\rightarrow$ Conférences Accra 1958 $\rightarrow$ Casablanca/Monrovia 1961 $\rightarrow$ OUA 1963.
    \item \textbf{Révisez en racontant l'histoire} : « De la diaspora au continent, de la pétition à l'action, de la division à l'UA... »
\end{enumerate}
\end{methode}

\begin{popfigure}
\begin{tikzpicture}[x=0.009\linewidth, y=0.6cm, font=\small]
% Axe principal
\draw[popdark, thick, ->] (1900,0) -- (2025,0) node[right] {Années};
% Périodes (blocs colorés)
\fill[popblue!15] (1900,-0.8) rectangle (1920,0.8);
\node[popdark, rotate=90, anchor=south] at (1905,0) {\bfseries Diaspora \& pionniers};
\fill[popgreen!15] (1919,-0.8) rectangle (1946,0.8);
\node[popdark, rotate=90, anchor=south] at (1932,0) {\bfseries Congrès panafricains};
\fill[poporange!15] (1945,-0.8) rectangle (1960,0.8);
\node[popdark, rotate=90, anchor=south] at (1952,0) {\bfseries Indépendances \& Accra};
\fill[popred!15] (1961,-0.8) rectangle (1964,0.8);
\node[popdark, rotate=90, anchor=south] at (1962,0) {\bfseries Divisions \& OUA};
\fill[poppurple!15] (1963,-0.8) rectangle (2003,0.8);
\node[popdark, rotate=90, anchor=south] at (1983,0) {\bfseries OUA (1963–2002)};
\fill[popgold!15] (2002,-0.8) rectangle (2025,0.8);
\node[popdark, rotate=90, anchor=south] at (2013,0) {\bfseries Union africaine};
% Événements repères (tiges + bulles)
\foreach \x/\y/\txt/\col in {
    1900/1.2/{Conf. panafricaine\\Londres (Williams)}/popblue,
    1914/1.6/{UNIA\\Garvey}/popblue,
    1919/2.0/{1er Congrès\\Paris (Du Bois)}/popgreen,
    1945/2.4/{5e Congrès\\Manchester\\Nkrumah, Kenyatta}/popgreen,
    1957/2.8/{Indépendance\\Ghana}/poporange,
    1958/3.2/{Conférences\\Accra}/poporange,
    1961/3.6/{Casablanca /\ Monrovia}/popred,
    1963/4.0/{Création OUA\\Addis-Abeba\\25 mai}/popred,
    2002/4.4/{Acte UA\\Lomé/Sirte}/poppurple,
    2021/4.8/{Lancement\\ZLECAf}/popgold} {
    \draw[\col, thick] (\x,0) -- (\x,\y);
    \node[\col, fill=white, draw=\col, rounded corners=2pt, inner sep=2pt, anchor=south, font=\tiny, align=center] at (\x,\y+0.1) {\txt};
}
% Graduations
\foreach \y in {1900,1920,1940,1960,1980,2000,2020} {
    \draw[popdark] (\y,-0.1) -- (\y,0.1) node[below, font=\tiny] {\y};
}
\end{tikzpicture}
\legende{Frise chronologique synthétique du panafricanisme (1900–2021). Les blocs colorés correspondent aux périodes méthodologiques ; les bulles aux événements repères à connaître pour l'examen.}
\end{popfigure}

\courssub{Exercice résolu : analyser une résolution du 5e congrès de Manchester}

\begin{exoresolu}[Analyse d'un extrait de résolution (Manchester, 1945)]
\textbf{Énoncé :} \og Nous affirmons le droit de tous les peuples coloniaux de contrôler leur propre destin. Tous les colonies doivent être libres de toute domination impériale étrangère, qu'elle soit politique ou économique. Nous exigeons l'indépendance immédiate et totale. \fg (Résolution n° 3, 5e Congrès panafricain, Manchester, octobre 1945).

\textbf{Question :} En quoi cette résolution marque-t-elle une rupture par rapport aux congrès précédents (1919–1927) ? Justifiez en mobilisant le contexte historique.

\tcblower
\textbf{Solution détaillée :}

\medskip
\textbf{1. Identification du document.} Extrait de la résolution n°3 du 5e Congrès panafricain de Manchester (15–21 octobre 1945). Congrès organisé par George Padmore, présidé par W.E.B. Du Bois, avec forte délégation africaine (Nkrumah, Kenyatta, Banda, etc.).

\medskip
\textbf{2. Rupture lexicale et revendicative.}
\begin{itemize}
    \item \textbf{Vocabulaire nouveau} : \og indépendance immédiate et totale \fg, \og contrôler leur propre destin \fg, \og domination impériale étrangère \fg (politique \textbf{et} économique).
    \item \textbf{Disparition} : plus de référence à la Société des Nations, aux mandats, aux pétitions, aux réformes graduelles.
\end{itemize}

\medskip
\textbf{3. Contexte explicatif (causes de la rupture).}
\begin{itemize}
    \item \textbf{Seconde guerre mondiale} : participation massive d'Africains (tirailleurs) ; promesses non tenues (Charte de l'Atlantique 1941 : droit des peuples à disposer d'eux-mêmes) ; affaiblissement des puissances coloniales.
    \item \textbf{Montée des nationalismes} : en Asie (Inde 1947, Indonésie 1945), en Afrique (grèves, émeutes, syndicats).
    \item \textbf{Nouveaux acteurs} : leaders africains formés en Occident (Nkrumah à Lincoln/Penn, Kenyatta à LSE), radicaux, marxistes pour certains (Padmore), qui rejettent le réformisme de Du Bois.
    \item \textbf{Guerre froide naissante} : l'URSS soutient la décolonisation ; le panafricanisme s'internationalise.
\end{itemize}

\medskip
\textbf{4. Conséquences directes.}
Nkrumah et Kenyatta rentrent au pays (1947) et créent/mènent les partis de masse (CPP au Ghana, KANU au Kenya) qui obtiennent l'indépendance (1957, 1963). Manchester est le \textbf{point de bascule} du panafricanisme diasporique vers le panafricanisme continental de libération.

\medskip
\textbf{Conclusion.} La résolution de Manchester opère une \textbf{mutation de nature} : du panafricanisme \emph{de demande} (droits, réformes, internationalisation) au panafricanisme \emph{de conquête} (indépendance, souveraineté, action de masse). C'est l'acte de naissance politique du panafricanisme moderne.
\end{exoresolu}

\begin{exercice}[Entraînement Probatoire — Dissertation]
\textbf{Sujet :} \og L'histoire du panafricanisme est-elle celle d'un échec ou d'une réinvention permanente ? \fg

\textbf{Consignes :}
\begin{enumerate}
    \item Problématisez le sujet (définissez les termes, identifiez la tension).
    \item Construisez un plan dialectique en trois parties (thèse / antithèse / synthèse) en vous appuyant sur \textbf{au moins quatre dates repères} de la frise.
    \item Rédigez l'introduction complète (accroche, définitions, problématique, annonce du plan).
\end{enumerate}
\end{exercice}

\begin{corrige}[Exercice n°1 — Éléments de correction]
\textbf{Problématique :} Le panafricanisme oscille entre l'idéal d'une Afrique unie, souveraine et prospère (réussite) et la réalité d'États fragmentés, dépendants, en proie aux conflits (échec apparent). La question est de savoir si cette tension est constitutive du mouvement (réinvention) ou signe son impuissance.

\textbf{Plan suggéré :}
\begin{itemize}
    \item \textbf{I. Un projet né de l'échec colonial, sans cesse réinventé par la diaspora (1900–1945).} Conférence de 1900 (Williams), Garvey/UNIA 1914 (échec Back to Africa $\rightarrow$ réinvention symbolique), Congrès Du Bois (échec SDN $\rightarrow$ réinvention intellectuelle), Manchester 1945 (rupture radicale).
    \item \textbf{II. L'indépendance conquise mais l'unité trahie : le compromis OUA (1957–1963).} Ghana 1957 (réussite), Accra 1958 (base arrière), Casablanca/Monrovia 1961 (division), OUA 1963 (compromis : souveraineté > unité ; frontières gelées). Échec de l'État fédéral nkrumahiste.
    \item \textbf{III. L'Union africaine et la ZLECAf : une réinvention institutionnelle et économique inachevée (2002–...).} UA 2002 (droit d'ingérence, Parlement, CdPS), NEPAD 2001 (appropriation du développement), ZLECAf 2021 (marché unique). Mais : coups d'État (Mali, Guinée, Burkina, Niger, Gabon), dépendance extraversion, déficit démocratique, mise en œuvre lente de la ZLECAf. Réinvention permanente, pas encore réussite aboutie.
\end{itemize}

\textbf{Introduction type :} \og L'indépendance du Ghana n'a de sens que si elle est liée à la libération totale de l'Afrique \fg (Nkrumah, 1957). Cette phrase résume la promesse fondatrice du panafricanisme : l'unité comme condition de la liberté. Pourtant, soixante ans après la création de l'OUA (25 mai 1963), l'Afrique compte 55 États, des frontières intangibles, une intégration commerciale inférieure à 15\% et des conflits récurrents. Faut-il y voir l'échec d'un rêve ? Ou la preuve que le panafricanisme, né dans la diaspora (Conférence de Londres 1900, Garvey 1914), institutionnalisé par les congrès (Manchester, 1945), porté par les pères des indépendances (Nkrumah, 1957), divisé (Casablanca/Monrovia, 1961), compromis (OUA, 1963), est un \emph{processus de réinvention permanente} face aux défis historiques ? Nous montrerons que le panafricanisme, loin d'avoir échoué, se réinvente à chaque étape — de la revendication identitaire à la construction institutionnelle (UA, 2002) et économique (ZLECAf, 2021) — sans avoir encore réalisé sa promesse originelle.
\end{corrige}

\coursec{Introduction : problématique et définitions}

La leçon qui s'ouvre ici situe la philosophie à l'endroit même où elle devient \emph{arme de construction} : le chantier de l'unité et du développement africains. Le \textbf{consciencisme} (Kwame Nkrumah) et le \textbf{panafricanisme} (mouvement séculaire) ne sont pas des chapelles intellectuelles ; ce sont des \emph{réponses historiques} à la fragmentation coloniale et à la dépendance néocoloniale.

\begin{definition}[Consciencisme]
Philosophie d'action élaborée par Kwame Nkrumah (Ghana), exposée dans \emph{Consciencism} (1964, rev. 1970). Elle vise à synthétiser les trois courants qui traversent l'Afrique (tradition africaine, apport islamo-arabe, héritage euro-chrétien) en une idéologie cohérente : le \textbf{socialisme africain} (humanisme, égalitarisme, justice sociale), pour guider la décolonisation totale (politique, économique, mentale).
\end{definition}

\begin{definition}[Panafricanisme]
Mouvement politique, intellectuel et culturel prônant l'unité, la solidarité et la coopération des peuples africains et de la diaspora, pour leur émancipation collective et leur rayonnement mondial. Il irrigue l'action des pères des indépendances (Nkrumah, Nyerere, Sékou Touré, Keïta, Cheikh Anta Diop) et fonde l'Organisation de l'Unité Africaine (OUA, 1963) puis l'Union Africaine (UA, 2002).
\end{definition}

\begin{aretenir}[Problématique directrice]
\emph{Comment les fondements idéologiques et philosophiques du consciencisme et l'historique du panafricanisme éclairent-ils les défis concrets de l'unité et du développement au Cameroun aujourd'hui ?} Cette question structure toute la leçon : du concept à l'acte citoyen, en passant par les institutions (CEMAC, UA, ZLECAf) et les réalités de terrain (frontières, filières, sécurité).
\end{aretenir}

\coursec{Contexte historique : colonisation, indépendances et quête d'identité}

\begin{experience}[Repères chronologiques : de la conférence de Berlin à l'UA]
\begin{center}
\begin{tabularx}{\linewidth}{|c|X|}\hline
\rowcolor{popblueL}
\textbf{Période / Événement} & \textbf{Signification pour l'unité africaine} \\ \hline
1884-1885 : Conférence de Berlin & Partage de l'Afrique sans Africains ; frontières artificielles, fragmentation ethnique, économies tournées vers l'extérieur. \\ \hline
1900-1945 : Congrès panafricains (Londres, Paris, New York, Manchester) & Éveil de la conscience diasporique (Du Bois, Garvey, Padmore) ; revendication du droit des peuples à disposer d'eux-mêmes. \\ \hline
1945 : Congrès de Manchester (5\textsuperscript{e} congrès) & Bascule : les leaders africains (Nkrumah, Kenyatta, Nyerere) prennent la main ; mot d'ordre : \emph{Indépendance maintenant}. \\ \hline
1958-1960 : Indépendances (Ghana 1957, Cameroun 1960, etc.) & Souveraineté formelle acquise ; mais économies extraverties, armées nationales fragiles, élites formées à l'étranger. \\ \hline
1961 (Casablanca) vs 1961 (Monrovia) & Deux visions : \textbf{fédéralisme immédiat} (Nkrumah, Sékou Touré, Modibo Keïta) vs \textbf{intégration progressive} (Houphouët-Boigny, Senghor, Tafawa Balewa). \\ \hline
25 mai 1963 : Création de l'OUA (Addis-Abeba) & Compromis : charte respectant l'intangibilité des frontières coloniales (\emph{uti possidetis juris}) ; Comité de libération (soutien aux luttes anticoloniales). \\ \hline
1980 : Plan d'action de Lagos ; 1991 : Traité d'Abuja (CEA) & Relance économique : marché commun, union monétaire, parlement panafricain. \\ \hline
2002 : Durban : naissance de l'Union Africaine (UA) & Institution supranationale (Commission, Parlement, Cour de justice, Conseil de paix et sécurité) ; Agenda 2063. \\ \hline
2018-2021 : ZLECAf (AfCFTA) & Projet économique phare : zone de libre-échange continentale, marché unique 1,4 milliard d'habitants. \\ \hline
\end{tabularx}
\end{center}
\end{experience}

\begin{attention}[Piège historiographique : ne pas confondre indépendance et souveraineté réelle]
L'indépendance politique (drapeau, hymne, siège à l'ONU) ne garantit pas la \emph{souveraineté économique} (maîtrise des ressources, transformation locale, monnaie, technologie) ni la \emph{souveraineté mentale} (décolonisation des programmes, des langues officielles, des imaginaires). Le consciencisme naît précisément de ce constat : \emph{le néo-colonialisme est la dernière étape de l'impérialisme} (Nkrumah, 1965).
\end{attention}

\coursec{Les fondements idéologiques du consciencisme}

\begin{propriete}[Les trois sources du consciencisme (Nkrumah, \emph{Consciencism}, chap. 1-3)]
L'idéologie nkrumahiste ne s'invente pas \emph{ex nihilo} ; elle \emph{synthétise} trois strates historiques vécues par les masses africaines :
\begin{enumerate}
    \item \textbf{L'héritage traditionnel africain} : Communautarisme (l'individu dans le clan), propriété collective de la terre, décision par consensus, éthique de la réciprocité (\emph{ubuntu} : « Je suis parce que nous sommes »). \emph{Valeur centrale :} Humanisme.
    \item \textbf{L'apport islamo-arabe} (via le Sahel, le Soudan, la côte est) : Monothéisme, discipline, écriture (ajami), droit, commerce transsaharien, structures étatiques centralisées (empires du Ghana, Mali, Songhaï, Kanem-Bornou, Sokoto). \emph{Valeur centrale :} Égalitarisme (fraternité des croyants).
    \item \textbf{L'héritage euro-chrétien} (colonisation, missions, écoles) : Rationalité scientifique, technologie, droits de l'homme (déclarations), organisation étatique moderne, syndicalisme. \emph{Valeur centrale :} Justice sociale (héritage socialiste/chrétien).
\end{enumerate}
\end{propriete}

\begin{definition}[Synthèse conscienciste : le socialisme africain]
Ce n'est pas l'importation du marxisme-léninisme, ni un retour au passé. C'est une \emph{création} : \textbf{« La société traditionnelle africaine était égalitaire ; le socialisme africain en est la reprise consciente à l'ère moderne »} (Nkrumah). Principes : propriété collective des moyens de production stratégiques, planification démocratique, priorité aux besoins de base (santé, éducation, logement), non-alignement, unité continentale.
\end{definition}

\begin{aretenir}[Conversion catégorique : de la philosophie à l'idéologie]
Pour Nkrumah, la philosophie ne suffit pas ; elle doit subir une \emph{conversion catégorique} en \textbf{idéologie} (système d'idées qui guide l'action politique de masse). Le consciencisme est cette \emph{philosophie devenue idéologie} pour la révolution africaine. \textbf{Exigence :} Le philosophe-technicien doit \emph{traduire} les concepts en programmes concrets (Plan national, ZLECAf, éducation technique).
\end{aretenir}

\coursec{Les fondements philosophiques du consciencisme}

\begin{propriete}[Matérialisme dialectique et épistémologie de la praxis]
Nkrumah ancre le consciencisme dans un \textbf{matérialisme dialectique} non dogmatique :
\begin{itemize}
    \item \textbf{Ontologie :} La matière est première ; l'esprit émerge de l'organisation complexe de la matière (cerveau, société). Pas de dualisme corps/esprit.
    \item \textbf{Épistémologie :} La connaissance naît de la \emph{pratique} (travail, lutte, production). \emph{« La pratique est le critère de la vérité. »} Le savoir technique (ingénierie, agronomie, médecine) est une forme supérieure de conscience.
    \item \textbf{Dialectique :} Contradiction comme moteur de l'histoire : tradition vs modernité, Afrique vs Occident, classes populaires vs élites compradores. La \emph{synthèse} n'est pas un compromis mou, mais un saut qualitatif (unité fédérale, industrialisation).
\end{itemize}
\end{propriete}

\begin{exemplebox}[Le concept de « personnalité africaine »]
Cheikh Anta Diop (\emph{Nations nègres et culture}, 1954) et Nkrumah convergent : l'unité africaine repose sur une \emph{unité culturelle profonde} (matrilinéarité, totémisme, cosmogonie, linguistique négro-égyptienne) masquée par la diversité apparente. Le consciencisme en fait le \emph{socle ontologique} de la fédération : on ne fédère pas des étrangers, on \emph{retrouve} une famille dispersée.
\end{exemplebox}

\begin{attention}[Ne pas réduire le consciencisme à une « idéologie d'État »]
Au Ghana, le consciencisme a servi de doctrine au parti unique (CPP) et à l'autoritarisme croissant de Nkrumah (loi de détention préventive, 1958). \textbf{Leçon pour le technicien :} Une philosophie d'action sans garde-fous démocratiques (séparation des pouvoirs, presse libre, syndicats autonomes) se retourne contre le peuple qu'elle prétend servir. Le panafricanisme d'aujourd'hui intègre \emph{l'État de droit} et la \emph{gouvernance} (Charte africaine de la démocratie, 2007).
\end{attention}

\coursec{Historique du panafricanisme}

\begin{methode}[Trois temps pour retenir l'historique du panafricanisme]
\begin{enumerate}
    \item \textbf{Genèse diasporique (1900-1945) :} Du Bois (premier congrès 1900), Garvey (UNIA, « Back to Africa », Black Star Line), Padmore (syndicalisme international). \emph{Mots-clés :} Race, réparation, retour, solidarité noire mondiale.
    \item \textbf{Appropriation continentale (1945-1963) :} Manchester 1945 (Nkrumah, Kenyatta, Nyerere) ; Indépendances ; Rivalité Casablanca (fédéralisme radical) / Monrovia (coopération souple) ; Création OUA 1963 (compromis : intangibilité frontières + libération restante).
    \item \textbf{Institutionnalisation et relance (1963-aujourd'hui) :} Lagos 1980 (économie) ; Abuja 1991 (Communauté Économique Africaine, 6 étapes sur 34 ans) ; Durban 2002 (UA : organes supranationaux, droit d'ingérence humanitaire) ; Agenda 2063 (vision 50 ans) ; ZLECAf 2018/2021 (marché unique).
\end{enumerate}
\end{methode}

\begin{center}
\begin{tabularx}{\linewidth}{|l|X|X|}\hline
\rowcolor{popblueL}
\textbf{Figure clé} & \textbf{Apport spécifique au panafricanisme} & \textbf{Lien avec le consciencisme} \\ \hline
\textbf{W.E.B. Du Bois} & Intellectuel organisateur (congés 1900-1945) ; concept de « double conscience » ; encyclopédie africaine. & Base intellectuelle : l'unité naît de la prise de conscience historique. \\ \hline
\textbf{Marcus Garvey} & Mobilisation de masse (UNIA, millions d'adhérents) ; fierté noire ; entreprise économique (Black Star Line). & Dimension \emph{populaire} et \emph{économique} : l'unité se construit par le peuple et pour le peuple. \\ \hline
\textbf{George Padmore} & Théoricien du syndicalisme panafricain ; conseiller de Nkrumah ; \emph{Pan-Africanism or Communism?} & Pont vers l'action politique : le panafricanisme comme stratégie de libération nationale. \\ \hline
\textbf{Kwame Nkrumah} & \emph{Africa Must Unite} (1963) ; « Seek ye first the political kingdom » ; architecte OUA ; consciencisme. & \textbf{Fusion} : le panafricanisme devient \emph{philosophie d'État} et projet fédéral concret. \\ \hline
\textbf{Julius Nyerere} & \emph{Ujamaa} (socialisme villageois) ; Arusha Declaration ; soutien inconditionnel aux mouvements de libération. & Éthique du service : l'unité au service des plus pauvres (villagisation, éducation pour l'autonomie). \\ \hline
\textbf{Cheikh Anta Diop} & \emph{Nations nègres et culture} ; unité culturelle négro-africaine ; fédéralisme culturel et linguistique (bantou, peul, arabe comme langues fédérales). & \textbf{Fondement anthropologique/historique} : preuve scientifique de l'unité pour légitimer la fédération. \\ \hline
\end{tabularx}
\end{center}

\begin{savaistu}[Le Cameroun dans l'histoire panafricaine]
Ruben Um Nyobé (UPC) intervient à l'ONU dès 1952 pour l'indépendance du Cameroun : \emph{premier Africain à tribune onusienne} pour son pays. Félix-Roland Moumié (UPC) tisse le réseau panafricain (Ghana, Guinée, Égypte, Algérie, Chine). Ernest Ouandié incarne la résistance armée. L'UPC est \emph{le} parti panafricain camerounais par excellence (membre du Bureau politique panafricain du Caire, 1962). Leur répression (loi 1957, guerre 1955-1971) a retardé l'ancrage panafricain dans l'État post-colonial camerounais, mais leur héritage inspire la société civile actuelle.
\end{savaistu}

\coursec{Unité africaine et développement : applications concrètes au Cameroun}

Après avoir retracé l'\emph{historique du panafricanisme} — des premiers congrès aux indépendances, en passant par le réveil des consciences dans la diaspora — nous mesurons désormais l'ampleur du défi : transformer l'idéal unitaire en levier concret de développement. Le consciencisme de Nkrumah n'est pas une doctrine de salon ; il prétend être une \emph{philosophie d'action}. Comment, dès lors, l'idée d'unité africaine irrigue-t-elle le quotidien des Camerounais, structure leurs institutions et oriente leurs choix économiques ? C'est à cette application concrète que nous consacrons cette section, en articulant échelles locale, sous-régionale et continentale.

\courssub{La monnaie commune : le franc CFA, intégration monétaire et débat sur la souveraineté}

\begin{definition}[Intégration monétaire]
Processus par lequel plusieurs États adoptent une monnaie unique ou lient leurs monnaies par un taux de change fixe irrévocable, géré par une banque centrale commune, afin de faciliter les échanges et la stabilité macroéconomique.
\end{definition}

Le franc CFA (\emph{Franc de la Coopération Financière en Afrique}, devenu \emph{Franc de la Communauté Financière Africaine}) est l'exemple le plus tangible d'intégration vécue au quotidien par les populations de la CEMAC (Communauté Économique et Monétaire de l'Afrique Centrale). Six États — Cameroun, Tchad, République centrafricaine, Congo, Gabon, Guinée équatoriale — partagent cette monnaie, émise par la BEAC (Banque des États de l'Afrique Centrale), dont le siège est à Yaoundé.

\begin{exemplebox}[Le franc CFA dans la vie courante]
Un commerçant de Douala achète des tissus au Nigeria (naira) et les revend à Bangui (franc CFA) ou à Libreville (franc CFA) sans risque de change à l'intérieur de la zone. La parité fixe avec l'euro (1 € = 655,957 FCFA) garantit une stabilité des prix importés (riz, médicaments, carburant) que n'ont pas les pays à monnaie flottante.
\end{exemplebox}

Pourtant, cette intégration suscite un débat philosophique et politique vif, au cœur du consciencisme : \emph{une monnaie commune est-elle compatible avec une réelle souveraineté ?}

\begin{itemize}
    \item \textbf{Thèse de la nécessité (stabilité, crédibilité) :} La mutualisation des réserves de change (garantie de convertibilité par la France, siège au conseil d'administration de la BEAC) protège de l'hyperinflation. C'est une forme de \emph{solidarité monétaire} panafricaine pragmatique.
    \item \textbf{Thèse de l'aliénation (perte de leviers) :} L'impossibilité de dévaluer pour relancer l'exportation, l'ancienne obligation de déposer 50 \% des réserves à la Banque de France (accords de 1973, supprimée par l'accord de 2019 ratifié en 2020), la présence d'un commissaire du gouvernement français au conseil d'administration (supprimée en 2020) sont perçues comme une \emph{servitude volontaire} contradictoire avec l'idéal nkrumahiste d'autonomie totale.
\end{itemize}

\begin{attention}[Piège fréquent : opposer unité africaine et souveraineté nationale]
Ne confondez pas \emph{intégration} et \emph{assimilation}. Le panafricanisme (Nkrumah, Nyerere, Cheikh Anta Diop) vise une \textbf{unité fédérale} où les États cèdent des compétences (monnaie, défense, diplomatie) à un pouvoir supranational \emph{africain}, non à une puissance extérieure. Le débat sur le franc CFA porte précisément sur la nature du partenaire garant : la France ou l'Union africaine ? Critiquer le franc CFA n'est pas rejeter l'unité monétaire africaine, c'en est au contraire une condition pour la rendre \emph{endogène} (projet Eco, banque centrale africaine unique).
\end{attention}

\courssub{Le Cameroun, carrefour institutionnel : CEMAC, CEEAC, UA, Commonwealth}

Le Cameroun incarne par sa géographie et son histoire la \emph{complexité} de l'unité africaine. Il est le seul pays membre à la fois de la \textbf{CEMAC} (zone franc, marché commun centrafricain), de la \textbf{CEEAC} (Communauté Économique des États de l'Afrique Centrale, élargie à l'Angola, au Burundi, au Rwanda, à la RDC, à São Tomé), de l'\textbf{Union africaine} (organisation continentale politique) et du \textbf{Commonwealth} (héritage britannique, réseau anglophone de 56 États).

\begin{savaistu}[Double héritage colonial et appartenances multiples]
Le Cameroun est né de la réunification (1961) entre l'ancien Cameroun français (indépendant en 1960) et le Southern Cameroons britannique (qui a voté l'union au Cameroun plutôt que l'intégration au Nigeria). Cette double colonisation explique son bilinguisme officiel (français/anglais) et sa double appartenance : CEMAC/CEEAC (héritage francophone) et Commonwealth (héritage anglophone). C'est un \emph{microcosme} de l'Afrique : si l'unité y réussit, elle peut réussir partout.
\end{savaistu}

\begin{center}
\begin{tabularx}{\linewidth}{|l|X|X|}\hline
\rowcolor{popblueL}
\textbf{Organisation} & \textbf{Nature de l'intégration} & \textbf{Apport concret pour le Cameroun} \\ \hline
\textbf{CEMAC} & Union monétaire + Union douanière (TEC : Tarif Extérieur Commun) & Monnaie stable, marché de 60 millions de consommateurs, libre circulation des capitaux. \\ \hline
\textbf{CEEAC} & Communauté économique élargie (11 États), sécurité collective (COPAX) & Accès au marché congolais (Kinshasa), corridor Lobito (Angola), force multinationale anti-Boko Haram. \\ \hline
\textbf{Union africaine} & Intégration politique, Agenda 2063, ZLECAf, Passeport africain & Cadre normatif continental, légitimité diplomatique, projets d'infrastructures (PIDA). \\ \hline
\textbf{Commonwealth} & Réseau de coopération linguistique, juridique, éducative, commerciale & Accès aux marchés anglophones (Nigeria, Afrique du Sud, Kenya, Inde, UK), bourses Chevening, harmonisation du droit des affaires (OHADA compatible). \\ \hline
\end{tabularx}
\end{center}

Cette multiplicité n'est pas une contradiction : elle illustre le principe conscienciste de \emph{synthèse}. Le Cameroun ne choisit pas \emph{entre} l'Afrique centrale et l'Afrique anglophone ; il les \emph{relie}. C'est une application concrète de l'unité dans la diversité.

\courssub{L'unité nationale camerounaise : échelle locale du panafricanisme}

\begin{propriete}[Unité nationale comme condition de l'unité continentale]
\emph{« L'unité de l'Afrique ne peut être réalisée que par l'unité effective de ses États constitutifs. »} (Nkrumah, \emph{L'Afrique doit s'unir}). Au Cameroun, l'unité nationale repose sur trois piliers qui sont autant d'exercices quotidiens de panafricanisme \emph{par le bas} :
\end{propriete}

\begin{enumerate}
    \item \textbf{Le bilinguisme officiel (français/anglais) :} Ce n'est pas une simple administration ; c'est une \emph{philosophie du vivre-ensemble}. Chaque acte administratif, chaque panneau routier, chaque émission de la CRTV incarne la reconnaissance de l'Autre culturel. C'est une \emph{éducation permanente à la différence}, prémisse de l'unité continentale.
    \item \textbf{La réunification de 1961 :} Acte volontaire, négocié, ratifié par plébiscite (ONU). Elle prouve que des peuples aux histoires coloniales divergentes peuvent \emph{choisir} l'unité. C'est un précédent juridique et moral pour la fédération africaine.
    \item \textbf{Le « vivre-ensemble » (politique nationale depuis 2010) :} Au-delà du slogan, c'est la gestion quotidienne de la diversité ethnique (250 groupes), religieuse (christianisme, islam, animisme) et régionale. Les mariages interethniques, les marchés mixtes, les équipes de football nationales sont des \emph{faits sociaux} d'unité.
\end{enumerate}

\begin{attention}[Ne pas essentialiser le « vivre-ensemble »]
Le vivre-ensemble n'est pas un état naturel ni acquis définitivement. Les crises anglophones (depuis 2016), les tensions intercommunautaires (Logone-et-Chari, 2021) rappellent que l'unité se \emph{construit} par le droit, la justice distributive et le dialogue. Le consciencisme exige une \emph{veille citoyenne} permanente, pas une béatitude.
\end{attention}

\courssub{Situation concrète 1 : Libre circulation et commerce transfrontalier (Extrême-Nord)}

\begin{experience}[Enquête de terrain : le corridor Kousseri–N'Djamena–Maiduguri]
\textbf{Observation :} Au marché de Kousseri (frontière Tchad) et à la frontière de Banki/Amchidé (Nigeria), des milliers de petits commerçants (majoritairement des femmes) traversent quotidiennement la frontière.
\textbf{Données :} Le commerce informel transfrontalier représenterait 30 à 40 \% du PIB de la région de l'Extrême-Nord (études CILSS/CEEAC).
\textbf{Analyse :} La réglementation CEMAC/CEEAC (carte de séjour, carte de commerçant ambulant) facilite \emph{théoriquement} ce flux. La réalité : harcèlement aux barrages, non-reconnaissance des cartes, fermetures unilatérales de frontières (crise Boko Haram, COVID-19).
\textbf{Interprétation :} L'unité africaine se heurte ici à la \emph{souveraineté régalienne} (contrôle des frontières, sécurité). Le panafricanisme concret exige d'harmoniser les procédures (guichets uniques frontaliers, reconnaissance biométrique) pour transformer la contrainte en opportunité.
\end{experience}

\begin{exoresolu}[Appliquer la notion d'intégration régionale au commerce frontalier]
\textbf{Énoncé :} En vous appuyant sur la notion d'intégration régionale, expliquez pourquoi la fermeture de la frontière camerouano-nigériane en 2019-2020 a constitué un frein au développement local, malgré les impératifs sécuritaires.
\textbf{Solution détaillée :}
\begin{enumerate}
    \item \textbf{Rappel de la notion :} L'intégration régionale (CEEAC, ZLECAf) vise la libre circulation des personnes, des biens et des services pour créer un marché élargi, source d'économies d'échelle et de spécialisation (avantage comparatif).
    \item \textbf{Description de la situation :} La frontière longue de 1 700 km est un espace de vie unique (peuples kanouri, mousgoum, arabes de part et d'autre). Le bétail, le poisson du lac Tchad, le mil, le carburant y circulent historiquement. La fermeture (décision unilatérale du Nigeria puis réciprocité) a stoppé ce flux.
    \item \textbf{Lien notion/situation :} La fermeture a brisé la \emph{division naturelle du travail} : le Nigeria fournit l'essence et les produits manufacturés bon marché ; le Cameroun/Tchad fournissent le bétail, l'oignon, le poisson séché. Résultat : hausse des prix (+ 40 \% pour l'essence à Kousseri), perte de revenus pour les éleveurs, appauvrissement des femmes commerçantes. La sécurité (lutte contre Boko Haram) a primé sur l'intégration économique, révélant la \emph{fragilité} de l'unité africaine quand les États ne mutualisent pas leur sécurité (police frontalière commune, renseignement partagé).
\end{enumerate}
\end{exoresolu}

\courssub{Situation concrète 2 : Solidarité africaine face aux crises (réfugiés, Force multinationale mixte)}

Le Cameroun est parmi les premiers pays d'accueil de réfugiés en Afrique (avec l'Ouganda, le Soudan, l'Éthiopie), avec plus de 480 000 réfugiés et demandeurs d'asile (HCR, 2024), majoritairement Centrafricains (Est) et Nigérians (Extrême-Nord).

\begin{definition}[Solidarité africaine (Charte de l'OUA, 1969 ; Convention de Kampala, 2009)]
Devoir moral et juridique des États africains d'accueillir et protéger les personnes déplacées pour des causes africaines (conflits, catastrophes), en partageant le fardeau (\emph{burden sharing}) plutôt qu'en le externalisant vers l'Europe.
\end{definition}

\begin{exemplebox}[Accueil des réfugiés centrafricains dans l'Est (Garoua-Boulaï, Bertoua)]
Depuis 2013, plus de 300 000 Centrafricains vivent dans des camps ou parmi la population hôte. Le Cameroun applique la \emph{politique de porte ouverte} : accès aux terres, aux centres de santé, aux écoles publiques. C'est une application concrète de l'article 12 de la Charte africaine des droits de l'homme et des peuples (droit d'asile). Coût estimé : milliards de FCFA/an supportés par le budget national et les partenaires (HCR, PAM, UE).
\end{exemplebox}

Sur le plan sécuritaire, la \textbf{Force Multinationale Mixte (FMM)} — Nigeria, Niger, Tchad, Cameroun, Bénin — illustre l'unité \emph{par la contrainte} face à Boko Haram.

\begin{center}
\begin{tabularx}{\linewidth}{|l|X|}\hline
\rowcolor{popblueL}
\textbf{Dimension} & \textbf{Réalité camerounaise} \\ \hline
\textbf{Militaire} & Quartier général de la FMM à N'Djamena ; secteur 1 (Cameroun) commandé par un général camerounais ; opérations conjointes (Opération \emph{Lac Tchad}). \\ \hline
\textbf{Juridique} & Accord de 2015 (Statut des forces) : immunités, règles d'engagement, droits de l'homme. \\ \hline
\textbf{Financière} & Contributions nationales inégales ; appui de l'UA (Fonds pour la paix) et partenaires (UE, USA). \\ \hline
\textbf{Limite} & Commandement non unifié (chaque contingent répond à son Chef d'État) ; méfiance persistante (ex. fermeture frontière Nigeria 2019). \\ \hline
\end{tabularx}
\end{center}

\begin{aretenir}[Unité par la sécurité : un pas vers la fédération ?]
La FMM est une \emph{armée africaine} embryonnaire. Si elle réussit à pacifier le bassin du Lac Tchad, elle aura prouvé que des États souverains peuvent mutualiser la \emph{force légitime} — attribut régalien par excellence — pour un bien commun africain. C'est le cœur du projet nkrumahiste : \emph{« Une armée africaine pour une Afrique unie »}.
\end{aretenir}

\courssub{Situation concrète 3 : La ZLECAf, chance pour la transformation locale (cacao, bois)}

La \textbf{Zone de Libre-Échange Continentale Africaine (ZLECAf)}, entrée en vigueur en 2021, est le projet économique le plus ambitieux de l'UA depuis 1963. Elle vise à créer un marché unique de 1,4 milliard de consommateurs, PIB cumulé 3 400 milliards \$.

\begin{definition}[ZLECAf (AfCFTA)]
Accord instituant une zone de libre-échange continentale : élimination des droits de douane sur 90 \% des lignes tarifaires (produits), puis 97 \% (produits sensibles exclus), libéralisation des services, protocole sur la concurrence, l'investissement, la propriété intellectuelle, le commerce numérique.
\end{definition}

\begin{exemplebox}[Cacao et bois : de l'exportation brute à la transformation locale]
\textbf{Situation actuelle :} Le Cameroun exporte 75 \% de son cacao en fèves brutes (vers Europe/USA) et 60 \% de son bois en grumes ou sciages bruts (vers Chine/UE). Valeur ajoutée captée à l'extérieur.
\textbf{Objectif ZLECAf :} Exporter du \emph{chocolat}, de la \emph{pâte de cacao}, du \emph{beurre de cacao} vers le Nigeria (200 millions de consommateurs), l'Égypte, l'Afrique du Sud, le Kenya. Exporter du \emph{meuble fini}, du \emph{parquet}, du \emph{panneau} vers la CEMAC et l'Afrique de l'Ouest.
\textbf{Leviers :} Tarif préférentiel 0 \% intra-africain vs 10-20 \% pour l'UE (APE) ou la Chine. Règles d'origine : transformation substantielle (ex. fèves $\to$ pâte = origine camerounaise).
\textbf{Obstacles :} Énergie (coupures), logistique (port de Douala saturé, route Douala-N'Djamena), normes (certification ISO, traçabilité déforestation UE), financement (taux d'intérêt 12-15 \%), main-d'œuvre qualifiée.
\end{exemplebox}

\begin{popfigure}
\begin{tikzpicture}[scale=0.8]
% Manual Pie Chart: Intra-African 15%, Rest 85%
\def\angleA{54} % 15% of 360
\def\angleB{306} % 85% of 360
% Slice 1: Intra-African (Green)
\filldraw[popgreen!80] (0,0) -- (3,0) arc (0:\angleA:3) -- cycle;
% Slice 2: Rest of World (Orange)
\filldraw[poporange!70] (0,0) -- (\angleA:3) arc (\angleA:360:3) -- cycle;
% Labels
\node[popdark, font=\small] at (1.5, -1.5) {Intra-africain (15\,\%)};
\node[popdark, font=\small] at (-1.5, 1.5) {Reste du monde (85\,\%)};
% Legend lines
\draw[popdark, thin] (1.5, -1.5) -- (1.2, -0.8);
\draw[popdark, thin] (-1.5, 1.5) -- (-1.0, 0.8);
\end{tikzpicture}
\legende{Part du commerce intra-africain dans le commerce total de l'Afrique (données CNUCED/UA, ~2023). L'objectif de la ZLECAf est de porter ce ratio à 25-30 \% d'ici 2035.}
\end{popfigure}

\begin{methode}[Pour appliquer une notion à une situation concrète : 3 temps]
\begin{enumerate}
    \item \textbf{Rappeler la notion} : Définir précisément le concept (ex. : ZLECAf = élimination droits de douane intra-africains + règles d'origine).
    \item \textbf{Décrire la situation} : Chiffres, acteurs, flux réels (ex. : cacao camerounais : 300 000 t/an, 75 \% brutes, acheteurs UE).
    \item \textbf{Montrer le lien (mécanisme causal)} : Expliquer \emph{comment} la notion modifie la situation (ex. : taux 0 \% + marché nigérian proche $\to$ incitation à installer une usine de pâte à Douala/Bertoua $\to$ emplois + devises retenues).
\end{enumerate}
\end{methode}

\courssub{Développement endogène : valoriser ressources et savoirs locaux}

\begin{definition}[Développement endogène]
Modèle de développement fondé sur la mobilisation prioritaire des \textbf{ressources propres} (naturelles, humaines, culturelles, financières, techniques) et des \textbf{valeurs endogènes} (solidarité, réciprocité, savoir-faire traditionnels), plutôt que sur la reproduction de modèles exogènes (import-substitution, ajustement structurel, aide conditionnée).
\end{definition}

Le consciencisme (Nkrumah) et la pensée de Cheikh Anta Diop convergent : \emph{l'Afrique doit « réinventer l'Afrique »} à partir de son génie propre.

\begin{center}
\begin{tabularx}{\linewidth}{|l|X|X|}\hline
\rowcolor{popblueL}
\textbf{Secteur} & \textbf{Modèle exogène (copié-collé)} & \textbf{Approche endogène (conscienciste)} \\ \hline
\textbf{Agriculture} & Intrants chimiques importés, semences hybrides brevetées, monoculture d'exportation (cacao, coton). & Agroécologie (engrais verts, compost, semences paysannes), polyculture-élevage, souveraineté alimentaire (riz, maïs, manioc, plantain). \\ \hline
\textbf{Santé} & Médecine curative hospitalière, importation de 80 \% des médicaments, évacuations sanitaires. & Pharmacopée traditionnelle normalisée (IMPM, centres de recherche), médecine préventive communautaire, production locale de génériques (CINPHARM). \\ \hline
\textbf{Éducation} & Programmes calqués sur France/USA, langue d'enseignement étrangère, diplôme = visa d'émigration. & Bilinguisme effectif, intégration savoirs endogènes (agriculture, artisanat, histoire locale), formation aux métiers du territoire. \\ \hline
\textbf{Industrie} & Usines « tournevis » (assemblage CKD), dépendance technologique, rapatriement des bénéfices. & Transformation locale matières premières (cacao $\to$ chocolat, bois $\to$ meuble, coton $\to$ textile, bauxite $\to$ alumine), transfert de technologie maîtrisé. \\ \hline
\end{tabularx}
\end{center}

\begin{exemplebox}[Success story endogène : le Ndolé industriel et les jus de fruits « Made in Cameroon »]
Des PME camerounaises (ex. \emph{Les Jus de Fruits du Cameroun}, \emph{Ndolé Pak}) transforment localement mangue, ananas, goyave, ndolé (feuilles de vernonia) en conserves, jus pasteurisés, plats cuisinés surgelés. Elles créent 200-500 emplois directs, achètent aux coopératives paysannes, exportent vers la diaspora (France, USA) \emph{et} vers le marché régional (Gabon, Guinée équatoriale, Tchad) grâce à la CEMAC. C'est le développement endogène en acte : \emph{ressource locale + savoir-faire culinaire + technologie accessible + marché régional}.
\end{exemplebox}

\courssub{Débat philosophique : l'unité africaine, utopie ou nécessité ?}

\begin{center}
\begin{tabularx}{\linewidth}{|l|X|X|}\hline
\rowcolor{popblueL}
\textbf{Thèse} & \textbf{Arguments (Utopie : impossible/irréaliste)} & \textbf{Arguments (Nécessité : vitale/réaliste)} \\ \hline
\textbf{1. Diversité} & 54 États, 2 000 langues, 3 000 ethnies, régimes politiques opposés (démocraties, dictatures, monarchies), niveaux de développement abyssaux (PIB/hab : Seychelles 15 000 \$ vs Burundi 230 \$). L'unité ignore le réel. & La diversité \emph{est} la richesse. L'Europe s'est unie avec 24 langues, des guerres séculaires, des écarts de richesse. L'unité \emph{gère} la différence, elle ne l'efface pas (fédéralisme, subsidiarité). \\ \hline
\textbf{2. Souveraineté} & Les Chefs d'État gardent jalousement leur pouvoir (siège à l'UA, armées, budgets). Aucun n'accepte un « Président des États-Unis d'Afrique » élu au suffrage universel. & La souveraineté \emph{absolue} est un mythe westphalien dépassé. L'interdépendance (climat, pandémies, terrorisme, finance) impose une \emph{souveraineté partagée} (UE, ZLECAf, FMM). Céder du pouvoir pour en gagner (poids mondial). \\ \hline
\textbf{3. Économie} & Marchés fragmentés, infrastructures coloniales (rails vers ports, pas entre capitales), barrières non tarifaires (corruption, paperasserie), devises non convertibles. Coût d'intégration > bénéfices à court terme. & Le marché unique (ZLECAf) = 3 400 Md\$ de PIB. Infrastructures PIDA (barrage Grand Inga, route Transsaharienne, chemin de fer Dakar-N'Djamena) en cours. La fragmentation coûte 2-3 \% PIB/an (CNUCED). \\ \hline
\textbf{4. Histoire/Identité} & L'« Africain » n'existe pas : il y a des Yorubas, des Bantous, des Peuls, des Arabes, des Copte. L'identité panafricaine est une construction intellectuelle d'élites occidentalisées. & L'identité se \emph{construit} par l'action commune (Fanon). La lutte anti-coloniale, la ZLECAf, la CAN, la musique (Afrobeats, Amapiano), la diaspora créent un \emph{« nous »} politique et culturel réel. \\ \hline
\textbf{5. Puissance extérieure} & France (Françafrique), Chine (nouveau partenaire hégémonique), USA (AFRICOM), UE (APE), Russie (Wagner/Africa Corps) divisent pour régner. L'unité dérange les prédateurs. & Précisément : \emph{seule l'unité permet de négocier d'égal à égal}. « Divisés, nous sommes des proies ; unis, nous sommes une puissance » (Nkrumah). L'UA parle d'une voix à l'ONU, au G20 (siège obtenu 2023). \\ \hline
\end{tabularx}
\end{center}

\begin{aretenir}[Synthèse dialectique : l'unité comme \emph{devenir} et non comme \emph{donné}]
Le consciencisme refuse le dilemme binaire. L'unité africaine n'est ni une utopie irréalisable, ni une nécessité qui s'impose d'elle-même. C'est un \emph{projet historique} — au sens hégélien/marxiste : une tâche que les Africains doivent \emph{s'assigner} par la conscience (idéologie), l'organisation (partis, syndicats, société civile), l'action (intégration régionale, ZLECAf, FMM) et la culture (éducation, médias, arts). \textbf{Le Cameroun, par sa position charnière et sa double culture, est un laboratoire de ce devenir.}
\end{aretenir}

\courssub{Bilan : consciencisme, panafricanisme et conscience citoyenne active}

\begin{propriete}[Du concept à l'acte citoyen]
Le consciencisme n'est pas une théorie pour professeurs ; c'est une \emph{arme de construction massive}. Il appelle chaque Camerounais, chaque Africain à :
\begin{enumerate}
    \item \textbf{Penser par soi-même} : refuser le prêt-à-penser (néocolonial ou populiste), analyser les intérêts réels (ex. : qui gagne au franc CFA ? qui perd à la fermeture des frontières ?).
    \item \textbf{Agir localement} : consommer local (cacao transformé, bois local, riz de Ndop), payer ses impôts (souveraineté budgétaire), voter éclairé, participer aux budgets participatifs communaux.
    \item \textbf{S'organiser transversalement} : syndicats panafricains, associations de jeunes transfrontalières (Extrême-Nord/Tchad/Nigeria), réseaux d'entrepreneurs ZLECAf, mouvements féministes africains.
    \item \textbf{Exiger des comptes} : transparence des accords miniers/pétroliers, reddition des comptes sur les fonds ZLECAf/FMM, lutte contre la corruption (CONAC, CONSUPE) comme acte panafricain (les détournements affaiblissent l'Afrique entière).
\end{enumerate}
\end{propriete}

\begin{exercice}[Appliquer le consciencisme à une situation locale]
\textbf{Sujet :} « Le développement de l'Afrique ne peut être que endogène. » Discutez cette affirmation en vous appuyant sur l'exemple de la filière cacao au Cameroun et les opportunités offertes par la ZLECAf.
\textbf{Consignes :} Structurez votre réponse en thèse/antithèse/synthèse. Mobilisez les notions de \emph{valeur ajoutée}, \emph{règles d'origine}, \emph{souveraineté alimentaire/industrielle}, \emph{conscience citoyenne}.
\end{exercice}

\begin{corrige}[Exercice n°1 : Filière cacao et développement endogène]
\textbf{Thèse (Oui, endogène nécessaire) :} L'histoire prouve que l'exportation brute (fèves) enrichit les transformateurs européens (chocolatiers suisses, français, allemands) et laisse le Cameroun à la merci des cours mondiaux (volatilité, termes de l'échange défavorables). Le développement endogène impose de transformer sur place : usines de broyage (pâte, beurre, poudre), chocolateries artisanales/industrielles. La ZLECAf offre le débouché : 1,4 milliard de consommateurs, classe moyenne africaine en croissance, goût pour le chocolat local. Exemple : \emph{Chocolaterie Robert} (Madagascar) ou \emph{Instant Chocolat} (Côte d'Ivoire) prouvent la viabilité.
\textbf{Antithèse (Contraintes exogènes réelles) :} L'endogène pur est impossible : machines (Allemagne/Chine), emballages (Europe), certification (ISO, Rainforest Alliance = normes occidentales), financement (banques internationales, taux élevés), énergie (dépendance hydrocarbures/barrages chinois), logistique (port Douala géré par concessionnaire MSC/Africa Global Logistics). Il faut \emph{négocier} l'exogène, pas le nier.
\textbf{Synthèse (Consciencisme : synthèse dialectique) :} Le développement endogène n'est pas l'autarcie. C'est la \emph{maîtrise du rapport aux autres}. Le Cameroun doit : (1) investir massivement dans l'énergie (barrage Nachtigal, solaire) et la formation technique (ingénieurs agroalimentaires, maintenanciers) ; (2) imposer des clauses de transfert de technologie et de contenu local dans les contrats miniers/industriels ; (3) utiliser la ZLECAf pour créer des \emph{chaînes de valeur régionales} (fèves Cameroun $\to$ pâte Cameroun $\to$ chocolat Nigeria/Ghana $\to$ marché CEDEAO/CEMAC) ; (4) éduquer le consommateur africain (fierté du « Made in Africa »). C'est là le rôle de la \emph{conscience citoyenne} : transformer le citoyen en acteur économique souverain.
\end{corrige}

\begin{savaistu}[Le Cameroun, pilote de la ZLECAf ?]
Le Cameroun a ratifié la ZLECAf dès 2019 (parmi les premiers). Yaoundé abrite le \textbf{Secrétariat permanent de la CEMAC} et le \textbf{Siège de la BEAC}. Le port de Kribi (eau profonde) et le corridor ferroviaire Kribi-Ngaoundéré (en projet) sont des infrastructures PIDA (Programme de Développement des Infrastructures en Afrique) conçues pour l'intégration. Si le Cameroun réussit sa transformation structurelle (Plan National de Développement 2020-2030), il deviendra le \emph{moteur industriel} de l'Afrique centrale — réalisation concrète du rêve panafricain.
\end{savaistu}

\coursec{Méthode : rédiger et justifier une démarche argumentative (application au Bac)}

\begin{methode}[Structure type d'une dissertation philosophique sur l'unité africaine]
\begin{enumerate}
    \item \textbf{Analyse du sujet} : Identifiez les concepts clés (unité, développement, utopie, nécessité, endogène), les termes de liaison (est-elle, peut-elle), la tension (idéal vs réel).
    \item \textbf{Problématisation} : Transformez le sujet en question philosophique. Ex. : « L'unité africaine est-elle une condition \emph{suffisante} du développement, ou une condition \emph{nécessaire mais non suffisante} qui exige en plus une transformation endogène des structures économiques et mentales ? »
    \item \textbf{Plan dialectique (Thèse / Antithèse / Synthèse)} :
        \begin{itemize}
            \item \textbf{Thèse} : L'unité est une nécessité vitale (poids démographique, marché unique, sécurité collective, voix internationale). Appuyez sur Nkrumah, UA, ZLECAf, FMM, exemples camerounais (CEMAC, réfugiés).
            \item \textbf{Antithèse} : L'unité bute sur des obstacles structurels (souverainetés jalouses, hétérogénéité économique, infrastructures coloniales, ingérences extérieures, déficit de légitimité démocratique). L'unité proclamée masque des divisions réelles (crises anglophones, coups d'État, fermeture frontières).
            \item \textbf{Synthèse} : L'unité est un \emph{processus historique} (consciencisme) : elle se construit \emph{par le bas} (intégration régionale effective, transformation locale, éducation panafricaine, conscience citoyenne) et \emph{par le haut} (institutions supranationales dotées de pouvoirs réels : Parlement africain législatif, Cour de justice contraignante, Banque centrale unique). Le Cameroun, par sa double culture et sa position géographique, est le test grandeur nature de cette synthèse.
        \end{itemize}
    \item \textbf{Rédaction} : Une idée = un paragraphe. Amorce (accroche + définition + problématique + annonce du plan). Développement (arguments étayés d'exemples précis, chiffrés, datés, nommés). Conclusion (réponse à la problématique + ouverture sur l'actualité : prochain sommet UA, entrée en vigueur phase 2 ZLECAf services/investissement).
    \item \textbf{Relecture} : Vérifiez la cohérence logique, l'orthographe des noms propres (Nkrumah, Cheikh Anta Diop, FMM, ZLECAf, CEMAC, BEAC), la pertinence des exemples camerounais.
\end{enumerate}
\end{methode}

\begin{attention}[Erreurs fatales au Probatoire]
\begin{itemize}
    \item \textbf{Confondre panafricanisme et négritude :} La négritude (Senghor, Césaire, Damas) est un mouvement \emph{littéraire et culturel} d'affirmation de l'identité noire. Le panafricanisme (Du Bois, Garvey, Nkrumah, Nyerere) est un mouvement \emph{politique et économique} d'union des États africains. Ne les mélangez pas.
    \item \textbf{Oublier le Cameroun :} En Première technique, \emph{toute} dissertation doit s'ancrer dans le réel camerounais. Un plan purement théorique sur « l'Afrique » sans citer la CEMAC, la ZLECAf au Cameroun, le bilinguisme, la FMM, le cacao/bois, est \emph{hors sujet} (note éliminatoire).
    \item \textbf{Généraliser à outrance :} « L'Afrique est pauvre », « Les Africains sont solidaires ». Non. Parlez du \emph{Cameroun}, du \emph{Nigeria}, de la \emph{CEMAC}, des \emph{femmes commerçantes de Kousseri}, des \emph{usines de transformation de Douala }. La philosophie technique exige le \emph{concret}.
\end{itemize}
\end{attention}

\coursec{Méthode : rédiger et justifier une démarche argumentative}

Après avoir analysé les applications concrètes de l'unité africaine au Cameroun — de la CEMAC à la ZLECAf en passant par le barrage de Lom Pangar — nous disposons désormais de matière solide pour construire une argumentation philosophique rigoureuse. La dissertation au Probatoire technique ne demande pas de réciter l'histoire, mais de \cle{problématiser} : transformer un thème en question qui appelle une discussion structurée, appuyée sur les doctrines étudiées (Nkrumah, Du Bois, Senghor, Wiredu, Cabral) et sur des exemples précis. Cette section vous donne la méthode pas à pas, avec un sujet type Probatoire entièrement travaillé.

\courssub{Rappel de la structure canonique de la dissertation philosophique}

\begin{methode}[Structure type de la dissertation au Probatoire Technique]
\textbf{Introduction (entonnoir) :}
\begin{enumerate}
    \item \textbf{Accroche concrète} : une citation, un fait d'actualité camerounaise, un paradoxe, une statistique (ex. : « En 2023, le commerce intra-africain ne représente que 15\% du commerce total du continent, contre 60\% en Europe »).
    \item \textbf{Définition des termes clés} : préciser le sens philosophique de « unité » (politique, économique, culturelle), « développement » (humain, durable, endogène), « condition » (nécessaire, suffisante, facilitatrice).
    \item \textbf{Problématique} : la question unique qui guide tout le devoir, issue de la tension entre les termes (ex. : « L'unité politique est-elle le levier indispensable d'un développement autonome ou une diversion face aux urgences nationales ? »).
    \item \textbf{Annonce du plan} : « Nous examinerons d'abord... (thèse), puis nous verrons que... (antithèse), pour enfin montrer que... (synthèse). »
\end{enumerate}

\textbf{Développement (corps du devoir) :}
\begin{itemize}
    \item \textbf{Plan dialectique (thèse / antithèse / synthèse)} : recommandé pour les sujets « Est-ce que... ? », « Faut-il... ? ».
    \item \textbf{Plan thématique (causes / conséquences / limites)} : utile pour « Quelles sont les conditions de... ? ».
    \item \textbf{Chaque grande partie = 1 idée directrice} formulée en phrase affirmative, subdivisée en 2 ou 3 sous-parties (argument + exemple + référence auteur).
    \item \textbf{Transitions logiques} : « Toutefois », « En outre », « Dès lors », « C'est pourquoi ».
\end{itemize}

\textbf{Conclusion :}
\begin{enumerate}
    \item \textbf{Réponse synthétique} à la problématique (une phrase claire).
    \item \textbf{Réouverture} : prolongement vers une question voisine, une perspective actuelle (ex. : ZLECAf, réforme de l'UA), ou une citation finale.
\end{enumerate}
\end{methode}

\begin{popfigure}
\begin{tikzpicture}[
    node distance=1.2cm and 1.5cm,
    box/.style={draw=popblue, thick, rounded corners, fill=popblueL, text width=3.2cm, align=center, minimum height=1.8cm},
    arrow/.style={-Stealth, thick, popdark},
    labelbox/.style={font=\small\bfseries, popdark}
]
% Introduction
\node[box, align=center] (intro){
    \textbf{INTRODUCTION}\\
    \textit{Entonnoir}\\
    Accroche $\rightarrow$ Définitions $\rightarrow$ \cle{Problématique} $\rightarrow$ Plan
};
% Thèse
\node[box, below left=of intro, xshift=-1cm, align=center] (these){
    \textbf{THÈSE}\\
    (Unité condition nécessaire)\\
    Arg. 1 : Nkrumah\\
    Arg. 2 : Marché unique\\
    Ex. : ZLECAf
};
% Antithèse
\node[box, below right=of intro, xshift=1cm, align=center] (antithese){
    \textbf{ANTITHÈSE}\\
    (Unité insuffisante / risques)\\
    Arg. 1 : Senghor / Wiredu\\
    Arg. 2 : Hétérogénéités\\
    Ex. : CEMAC, coups d'État
};
% Synthèse
\node[box, below=2.5cm of intro, align=center] (synthese){
    \textbf{SYNTHÈSE}\\
    (Unité différenciée, pragmatique)\\
    Intégration variable\\
    Ex. : Communautés régionales\\
    Réf. : Nkrumah révisé
};
% Conclusion
\node[box, below=1.2cm of synthese, align=center] (concl){
    \textbf{CONCLUSION}\\
    Réponse + Réouverture
};
% Flèches
\draw[arrow] (intro.south west) -- (these.north);
\draw[arrow] (intro.south east) -- (antithese.north);
\draw[arrow] (these.south) -- (synthese.north west);
\draw[arrow] (antithese.south) -- (synthese.north east);
\draw[arrow] (synthese.south) -- (concl.north);
\end{tikzpicture}
\legende{Schéma de la structure dialectique : l'entonnoir de l'introduction conduit à la problématique ; le développement oppose thèse et antithèse avant de les dépasser dans la synthèse ; la conclusion répond et rouvre.}
\end{popfigure}

\courssub{Justifier une réponse : l'argumentation philosophique}

\begin{definition}[Argumentation philosophique]
\cle{Justifier} une affirmation, c'est en donner les raisons logiques et les appuis factuels. En philosophie, tout énoncé doit reposer sur un triplet : \textbf{Thèse (affirmation) $\rightarrow$ Argument (raison logique) $\rightarrow$ Illustration (exemple concret ou référence textuelle)}.
\end{definition}

\begin{propriete}[Les trois piliers de la justification au Probatoire Technique]
\begin{enumerate}
    \item \textbf{La référence doctrinale} : citer explicitement un auteur au programme (Nkrumah, Du Bois, Senghor, Wiredu, Cabral) et résumer sa thèse en une phrase. \emph{Exemple : « Comme l'écrit Nkrumah dans \textit{Consciencism}, "l'indépendance politique n'est que le prélude à l'indépendance économique" ».}
    \item \textbf{L'exemple concret camerounais ou africain} : ancrer l'argument dans le réel (barrage de Lom Pangar, ZLECAf, CETA, échec de la CEMAC, politique de l'import-substitution au Cameroun dans les années 80).
    \item \textbf{Le lien logique} : expliciter \emph{pourquoi} l'exemple ou la citation prouve l'idée directrice. Ne jamais laisser une citation « flotter » sans analyse.
\end{enumerate}
\end{propriete}

\begin{attention}[Piège majeur : l'exposé d'histoire déguisé]
Ne racontez pas l'histoire du panafricanisme (conférences de 1900, 1919, 1945, OUA 1963, UA 2002) pour elle-même. \textbf{Chaque fait historique doit servir d'argument} pour ou contre la thèse en discussion.
\begin{itemize}
    \item \textbf{Mauvais} : « En 1963, l'OUA est créée à Addis-Abeba. En 2002, elle devient l'UA. » (Récit chronologique).
    \item \textbf{Bon} : « La Charte de l'OUA (1963) sacralise l'intangibilité des frontières coloniales (principe \textit{uti possidetis}), ce qui a freiné l'unité politique prônée par Nkrumah au profit d'une coopération interétatique ; c'est un argument pour l'antithèse : l'unité institutionnelle a primé sur l'unité fédérale. »
\end{itemize}
\end{attention}

\courssub{Construire une problématique : du sujet à la question philosophique}

\begin{methode}[Transformer un sujet en problématique (3 étapes)]
\begin{enumerate}
    \item \textbf{Analyser les termes} : identifier les concepts clés, leurs sens possibles, leurs présupposés.
    \item \textbf{Repérer la tension} : le sujet contient une affirmation implicite ou une opposition (condition nécessaire vs suffisante, unité vs diversité, politique vs économique, idéal vs réel).
    \item \textbf{Formuler la question} : elle doit être \emph{ouverte} (pas de réponse oui/non évidente), \emph{philosophique} (porte sur des concepts), \emph{unique} (une seule question directrice).
\end{enumerate}
\end{methode}

\begin{exemplebox}[Analyse guidée du sujet type Probatoire]
\textbf{Sujet :} « L'unité africaine est-elle une condition du développement de l'Afrique ? »

\textbf{Étape 1 : Analyse des termes}
\begin{itemize}
    \item \cle{Unité africaine} : sens politique (État fédéral panafricain \textit{à la} Nkrumah), économique (marché unique ZLECAf), culturelle (négritude, consciencisme), régionale (CEMAC, CEDEAO).
    \item \cle{Condition} : nécessaire (sans elle, pas de développement), suffisante (elle garantit le développement), ou facilitatrice (elle aide mais ne suffit pas).
    \item \cle{Développement} : croissance du PIB ? IDH ? Souveraineté alimentaire ? Industrialisation ? Développement humain (Sen) ? Développement endogène (Cabral) ?
\end{itemize}

\textbf{Étape 2 : Tensions repérées}
\begin{itemize}
    \item Nkrumah : unité politique \textbf{condition nécessaire} de l'indépendance économique réelle.
    \item Senghor / Wiredu : unité culturelle oui, mais unité politique fédérale \textbf{risque de tyrannie} ; le développement passe d'abord par la \textbf{démocratie locale} et la \textbf{gouvernance}.
    \item Réalité : 54 États, frontières coloniales, disparités linguistiques (anglophone/francophone/lusophone/arabophone), conflits, coups d'État. L'unité \textbf{est-elle réaliste} à court terme ?
\end{itemize}

\textbf{Étape 3 : Problématiques possibles (choisir la plus féconde)}
\begin{enumerate}
    \item « L'unité politique fédérale est-elle le \emph{seul} chemin vers un développement autonome, ou la coopération régionale variable (ZLECAf) suffit-elle ? »
    \item « L'unité africaine est-elle une \emph{condition nécessaire} du développement, ou une \emph{idéalisation} qui masque l'urgence de réformes nationales (gouvernance, éducation, industrie) ? »
    \item \textbf{Problématique retenue (exemple) :} \textbf{« L'unité politique et économique de l'Afrique est-elle la voie indispensable d'un développement souverain, ou une utopie face aux réalités nationales et aux intérêts géopolitiques ? »}
\end{enumerate}
\end{exemplebox}

\courssub{Rédaction guidée : introduction complète sur le sujet type}

Voici une introduction modèle, rédigée comme l'attend le correcteur du Probatoire technique. Chaque phrase a sa fonction.

\begin{exoresolu}[Introduction type Probatoire — Sujet : « L'unité africaine est-elle une condition du développement de l'Afrique ? »]
\textbf{Énoncé :} Sujet : « L'unité africaine est-elle une condition du développement de l'Afrique ? » — Rédigez une introduction complète (accroche, définitions, problématique, annonce de plan).

\tcblower

\textbf{Solution détaillée :}

\textbf{Accroche :} En 1963, à Addis-Abeba, Kwame Nkrumah lance aux chefs d'État : « Nous devons nous unir ou périr », affirmant que l'indépendance politique fragmentée est un leurre sans union fédérale. Soixante ans plus tard, la Zone de libre-échange continentale africaine (ZLECAf) entre en vigueur, mais le commerce intra-africain stagne à 15\% et le continent compte encore 33 pays parmi les moins avancés (PMA).

\textbf{Définitions :} L'\cle{unité africaine} désigne ici le projet d'intégration politique, économique et culturelle du continent, allant de la fédération panafricaine (Nkrumah) à l'intégration régionale variable (UA, ZLECAf). Le \cle{développement} s'entend au sens large : non seulement croissance économique, mais autonomie industrielle, souveraineté alimentaire, justice sociale et émancipation culturelle (Sen, Cabral). Une \cle{condition} est ce sans quoi un résultat ne peut advenir (nécessaire), ou ce qui le garantit (suffisante).

\textbf{Problématique :} Dès lors, \textbf{l'unité continentale est-elle le levier indispensable d'un développement souverain, ou une condition surestimée qui détourne de la refondation des États nationaux et de la gouvernance démocratique ?}

\textbf{Annonce du plan :} Nous soutiendrons d'abord, avec Nkrumah et le consciencisme, que l'unité politique est la \textbf{condition nécessaire} pour briser la dépendance néocoloniale (Thèse). Nous montrerons ensuite, avec Senghor, Wiredu et l'expérience des communautés régionales, que l'unité fédérale comporte des \textbf{risques de centralisme autoritaire} et que le développement suppose d'abord des États forts et légitimes (Antithèse). Nous dégagerons enfin une \textbf{synthèse pragmatique} : une unité \emph{différenciée}, économique d'abord (ZLECAf), politique par étapes, respectant les diversités, comme le préfigure le \textit{consciencisme} révisé.
\end{exoresolu}

\begin{aretenir}[Check-list d'une introduction réussie]
\begin{itemize}
    \item [$\square$] Accroche \textbf{concrète, datée, localisée} (Cameroun, Afrique, citation).
    \item [$\square$] Définitions \textbf{philosophiques} (pas le dictionnaire courant), distinguant les sens.
    \item [$\square$] Problématique \textbf{unique}, \textbf{ouverte}, issue de la tension des termes.
    \item [$\square$] Plan \textbf{annoncé} avec les verbes d'action (examiner, montrer, dégager) et l'indication des auteurs/exemples clés.
\end{itemize}
\end{aretenir}

\courssub{Exemple de plan dialectique rempli (tableau de révision)}

Ce tableau sert de \cle{fiche de révision} : il montre comment remplir chaque case du développement avec l'argument, la référence auteur et l'exemple concret.

\begin{popfigure}
\begin{center}
\begin{tabularx}{\linewidth}{|>{\bfseries}l|X|X|X|}
\hline
\rowcolor{popblueL}
\textbf{Partie} & \textbf{Idée directrice (phrase affirmative)} & \textbf{Arguments + Références (Auteurs au programme)} & \textbf{Exemples concrets (Cameroun / Afrique)} \\
\hline
\textbf{THÈSE}\newline(L'unité condition nécessaire) & L'unité politique fédérale est le seul rempart contre le néocolonialisme et le levier d'un développement endogène. & \textbf{Nkrumah (\textit{Consciencism}, \textit{Africa Must Unite})} : « L'indépendance politique sans unité économique est illusoire. »\newline\textbf{Du Bois} : panafricanisme comme réponse à la fragmentation de Berlin (1885).\newline\textbf{Cabral} : l'unité comme arme de libération nationale. & Projet d'\textbf{États-Unis d'Afrique} (Nkrumah).\newline\textbf{ZLECAf} (2021) : marché unique de 1,3 Md hab.\newline\textbf{Barrage de Lom Pangar} (Cameroun) : financement panafricain (BAD), énergie pour l'industrie locale.\newline Monnaie unique \textbf{Eco} (CEDEAO) : réduire coûts de transaction. \\
\hline
\textbf{ANTITHÈSE}\newline(Unité insuffisante / risques) & L'unité fédérale centralisée menace les libertés, ignore les réalités ethno-linguistiques et masque l'urgence de la gouvernance nationale. & \textbf{Senghor} : « La fédération est un leurre » ; priorité à la \textbf{Négritude} (unité culturelle) et à la démocratie locale.\newline\textbf{Wiredu} : consensus démocratique \textit{vs} majoritarisme ; décentralisation.\newline\textbf{Mazrui} : « Pax Africana » par coopération, pas fusion. & Échec de la \textbf{CEMAC} : intégration monétaire sans convergence budgétaire.\newline Crise \textbf{anglophone} au Cameroun : unité imposée \textit{vs} reconnaissance diversité.\newline Coups d'État (Mali, Burkina, Niger, Gabon) : États fragiles \textit{avant} union.\newline Dépendance persistante : matières premières exportées brutes. \\
\hline
\textbf{SYNTHÈSE}\newline(Unité différenciée, pragmatique) & L'unité est une \textbf{condition facilitatrice} mais non suffisante : elle doit être \emph{économique d'abord}, \emph{politique par étapes}, \emph{démocratique toujours}. & \textbf{Nkrumah révisé} : intégration par cercles concentriques (régional $\rightarrow$ continental).\newline\textbf{Consciencisme} : synthèse \textit{tradition + modernité + islam/chrétienté} $\rightarrow$ base culturelle d'une unité \emph{non uniformisante}.\newline\textbf{Agenda 2063 (UA)} : « L'Afrique que nous voulons » : intégration variable, libre circulation, passeport africain. & \textbf{ZLECAf} : suppression droits de douane 90\% lignes tarifaires.\newline\textbf{Corridor Douala-Ndjaména} : infrastructure régionale concrète.\newline\textbf{Passeport biométrique CEMAC} : libre circulation personnes.\newline\textbf{Cameroun} : hub logistique (port de Kribi, autoroute Yaoundé-Douala) au service intégration. \\
\hline
\end{tabularx}
\end{center}
\legende{Plan dialectique détaillé pour le sujet « L'unité africaine est-elle une condition du développement de l'Afrique ? ». Chaque ligne correspond à un grand paragraphe du développement.}
\end{popfigure}

\courssub{Critères d'évaluation au Probatoire technique}

\begin{definition}[Grille d'évaluation officielle (adaptée MINESEC)]
Le correcteur note sur 20 points répartis en quatre compétences :
\begin{enumerate}
    \item \textbf{Pertinence et problématisation (5 pts) :} Le sujet est-il compris ? La problématique est-elle philosophique, issue du sujet, et guide-t-elle vraiment le devoir ?
    \item \textbf{Connaissance des doctrines et auteurs (5 pts) :} Nkrumah, Du Bois, Senghor, Wiredu, Cabral sont-ils cités \emph{à bon escient} ? Les concepts (consciencisme, négritude, panafricanisme, néocolonialisme, développement endogène) sont-ils définis et utilisés correctement ?
    \item \textbf{Qualité de l'argumentation et de la structure (6 pts) :} Plan dialectique ou thématique clair ? Idées directrices affirmées ? Arguments logiques ? Exemples concrets (Cameroun/Afrique) \emph{analysés} et non juste listés ? Transitions fluides ?
    \item \textbf{Qualité de l'expression écrite (4 pts) :} Orthographe, syntaxe, vocabulaire philosophique précis (condition nécessaire/suffisante, aliénation, hétéronomie, autonomie, endogène, exogène), connecteurs logiques (donc, cependant, par conséquent, en somme).
\end{enumerate}
\end{definition}

\begin{attention}[Erreurs fatales à éviter absolument]
\begin{itemize}
    \item \textbf{Paraphraser le sujet} en introduction au lieu de le problématiser.
    \item \textbf{Accumuler les exemples} (dates, noms de conférences) sans les \emph{lier} à une thèse.
    \item \textbf{Confondre panafricanisme et consciencisme} : le premier est un mouvement historique/idéologique ; le second est une philosophie (épistémologie + éthique + politique) fondée par Nkrumah.
    \item \textbf{Oublier la conclusion} ou la réduire à « En conclusion, on a vu que... ». La conclusion \textbf{doit répondre} à la problématique.
    \item \textbf{Mélanger dissertation et commentaire de texte} : pas de citation longue sans analyse ; pas de « l'auteur dit que... » hors sujet.
    \item \textbf{Langage familier} : « Nkrumah a dit que... » $\rightarrow$ « Nkrumah soutient que... » / « Selon Nkrumah... ».
\end{itemize}
\end{attention}

\courssub{Exercice d'entraînement corrigé : de l'analyse du sujet à l'introduction}

\begin{exercice}[Sujet d'entraînement]
\textbf{Sujet :} « Le consciencisme peut-il être le fondement philosophique d'un développement africain authentique ? »
\begin{enumerate}
    \item Identifiez les trois concepts clés et leurs tensions.
    \item Proposez \textbf{deux} problématiques possibles. Choisissez la meilleure et justifiez votre choix en une phrase.
    \item Rédigez l'\textbf{introduction complète} (accroche, définitions, problématique choisie, annonce de plan) en vous appuyant sur le tableau de la section 5 et les cours précédents.
\end{enumerate}
\end{exercice}

\begin{corrige}[Exercice d'entraînement]
\textbf{1. Concepts \& tensions :}
\begin{itemize}
    \item \cle{Consciencisme} : philosophie de Nkrumah (matérialisme dialectique + humanisme traditionnel + réconciliation islam/chrétienté) $\rightarrow$ idéologie de la décolonisation mentale et de l'unité.
    \item \cle{Fondement philosophique} : base théorique, justification éthique/épistémologique, boussole pour l'action politique.
    \item \cle{Développement authentique} : endogène (Cabral), humain (Sen), souverain, non mimétique de l'Occident.
    \item \textbf{Tension} : Le consciencisme est-il \emph{opératoire} (guide pour politiques publiques, éducation, économie) ou \emph{utopique} (trop abstrait, ignoré par les élites, contredit par le néolibéralisme) ?
\end{itemize}

\textbf{2. Problématiques possibles :}
\begin{enumerate}
    \item « La synthèse philosophique du consciencisme (tradition + modernité) offre-t-elle une boussole efficace pour des politiques de développement concrètes, ou reste-t-elle un horizon théorique déconnecté des réalités de la gouvernance africaine ? »
    \item « Le consciencisme, en réconciliant les trois courants de la conscience africaine, peut-il fonder une éthique du développement qui dépasse l'alternative capitalisme/socialisme importée ? »
\end{enumerate}
\textbf{Choix :} La \textbf{problématique 1} est meilleure : elle oppose directement \emph{théorie} et \emph{pratique}, \emph{philosophie} et \emph{politique publique}, ce qui structure naturellement un plan dialectique (Thèse : oui, boussole éthique/épistémique ; Antithèse : non, gap théorie/pratique, élites néocoloniales ; Synthèse : oui, à condition d'être \emph{institutionnalisé} dans l'éducation, la constitution, la planification).

\textbf{3. Introduction rédigée :}
« "La philosophie ne doit pas seulement interpréter le monde, elle doit le transformer", écrit Marx dans ses \textit{Thèses sur Feuerbach}. Kwame Nkrumah, en fondant le consciencisme, entend précisément doter l'Afrique d'une philosophie \emph{guide} pour transformer la réalité coloniale en développement authentique. Pourtant, soixante ans après l'indépendance, le continent peine à traduire ses richesses en bien-être partagé.

Le \textbf{consciencisme} est la doctrine philosophique élaborée par Nkrumah (\textit{Consciencism}, 1964) qui synthétise trois pôles de la conscience africaine : l'héritage traditionaliste (communautarisme, spiritualité), l'héritage islamo-arabe (fraternité, justice sociale) et l'héritage euro-chrétien (rationnalisme technique, droits de l'homme), sur une base matérialiste dialectique. Un \textbf{fondement philosophique} est le socle théorique qui légitime et oriente l'action politique. Le \textbf{développement authentique}, chez Cabral comme chez Sen, désigne un processus endogène, auto-centré, libérateur, qui ne se réduit pas à la croissance du PIB.

Dès lors, \textbf{le consciencisme, en tant que synthèse philosophique de l'identité africaine, constitue-t-il une boussole opératoire pour des politiques de développement concrètes, ou reste-t-il un horizon théorique noble mais inopérant face aux réalités du néolibéralisme et de la gouvernance prédatrice ?}

Nous montrerons d'abord que le consciencisme offre une \textbf{anthropologie philosophique} (homme communautaire, matière-esprit) et une \textbf{éthique de la libération} (anti-impérialisme, justice sociale) qui fondent la légitimité d'un développement endogène (Thèse, Nkrumah, Cabral). Nous analyserons ensuite les \textbf{obstacles structurels} : absence d'institutionnalisation (éducation, droit), acculturation des élites, dépendance épistémique, qui vident le consciencisme de sa force transformatrice (Antithèse, Wiredu, Mazrui, réalité camerounaise). Nous conclurons que le consciencisme redevient fondement \textbf{s'il est traduit en institutions} : charte de l'éducation, planification stratégique, contrôle citoyen, comme le préfigure l'Agenda 2063. »
\end{corrige}

\courssub{Méthode de révision active pour l'épreuve}

\begin{methode}[Préparer le Probatoire en 3 séances de 1h]
\begin{enumerate}
    \item \textbf{Séance 1 — Fiches auteurs (30 min) + Fiches concepts (30 min).} Une fiche A5 par auteur (Nkrumah, Du Bois, Senghor, Wiredu, Cabral) : \emph{Thèse centrale} + \emph{Citation clé} + \emph{Exemple d'utilisation en dissertation}. Une fiche par concept (Unité, Développement, Négritude, Consciencisme, Néocolonialisme, Souveraineté) : \emph{Définition philosophique} + \emph{Auteurs liés} + \emph{Exemple camerounais}.
    \item \textbf{Séance 2 — Entraînement « Sujet $\rightarrow$ Problématique $\rightarrow$ Plan » (1h).} Prenez 3 sujets types (annales Probatoire, sujets probables). Pour chacun : analyse termes (5 min), 2 problématiques (5 min), choix + justification (5 min), plan détaillé avec arguments/auteurs/exemples (30 min), introduction rédigée (15 min). \textbf{Chronométrez-vous.}
    \item \textbf{Séance 3 — Rédaction d'un development partiel (1h).} Choisissez \emph{une} partie (Thèse ou Antithèse) d'un des plans de la séance 2. Rédigez-la \emph{entièrement} (2 grandes parties, 2 arguments chacune, transitions, conclusion partielle). Relisez avec la grille d'évaluation (pertinence, auteurs, exemples, expression).
\end{enumerate}
\end{methode}

\begin{savaistu}[Le consciencisme dans le système éducatif camerounais ?]
Le Cameroun a intégré l'\emph{éducation à la citoyenneté} et la \emph{culture nationale} dans les programmes (Loi d'orientation 1998, révisée 2023). Mais le \textbf{consciencisme} comme \emph{philosophie officielle} n'est pas enseigné en tant que tel. En revanche, la \textbf{philosophie en classe de Première Technique} (ce cours) en fait un chapitre central. Certains chercheurs (ex. Pr. Fabien Eboussi Boulaga, Pr. Jean-Marc Ela) ont plaidé pour une \emph{philosophie africaine} ancrée dans les réalités locales (palabre, consensus, terre) comme fondement du droit et de l'économie. La ZLECAf et l'Agenda 2063 réactivent ce besoin : une \emph{éthique panafricaine} pour réguler le marché unique. Votre maîtrise de ce chapitre est donc aussi une préparation à la citoyenneté active.
\end{savaistu}

\begin{aretenir}[Récapitulatif ultime pour l'épreuve]
\begin{itemize}
    \item \textbf{Problématiser} = transformer le sujet en question philosophique ouverte.
    \item \textbf{Argumenter} = Idée directrice + Argument logique + Auteur du programme + Exemple concret (Cameroun/Afrique) + Lien explicite.
    \item \textbf{Structurer} = Thèse / Antithèse / Synthèse (ou plan thématique) avec transitions.
    \item \textbf{Conclure} = Réponse directe à la problématique + Réouverture.
    \item \textbf{Vocabulaire} : condition nécessaire / suffisante / facilitatrice ; endogène / exogène ; aliénation / émancipation ; néocolonialisme ; consciencisme / négritude / panafricanisme.
    \item \textbf{Zéro tolérance} : récit historique pur, citations sans analyse, conclusion absente, hors-sujet.
\end{itemize}
\end{aretenir}

\newpage

\coursec{Activité d'intégration}

\coursec{Leçon 11 : Consciencisme, panafricanisme et développement de l'Afrique}

\courssub{Introduction : Problématique et définitions}

\begin{definition}[Consciencisme]
Philosophie politique élaborée par \textbf{Kwame Nkrumah} (1909--1972), premier président du Ghana, exposée dans \emph{Consciencism} (1964). Elle propose une \cle{synthèse dialectique} de trois courants culturels présents en Afrique pour fonder l'unité et le développement :
\begin{itemize}
    \item Le \cle{matérialisme dialectique} (apport socialiste/marxiste : analyse des infrastructures économiques, lutte contre l'impérialisme).
    \item L'\cle{humanisme traditionnel africain} (valeurs communautaires, \emph{Ubuntu}, primauté de l'humain sur la matière).
    \item Le \cle{christianisme} (éthique de la charité, justice sociale, dignité de la personne).
\end{itemize}
Le consciencisme n'est pas un syncrétisme statique mais une \emph{méthode d'analyse} pour transformer la société africaine par une planification continentale socialiste.
\end{definition}

\begin{definition}[Panafricanisme]
Mouvement idéologique et politique visant à l'\cle{unité africaine} (politique, économique, culturelle) et à la solidarité entre les peuples d'Afrique et de sa diaspora. Il traverse quatre phases historiques majeures (voir §5).
\end{definition}

\begin{aretenir}[Problématique de la leçon]
Comment le consciencisme (théorie) et le panafricanisme (processus historique) fondent-ils une stratégie de développement endogène pour le Cameroun et l'Afrique face à la persistance des structures extractives (exportation brute, commerce intra-africain faible) ? L'unité africaine (ZLECAf) est-elle une condition \emph{nécessaire et suffisante} du développement ?
\end{aretenir}

\courssub{Contexte historique : Colonisation, indépendances et quête d'identité}

\begin{savaistu}[Du partage à la réunification camerounaise]
La conférence de Berlin (1884--1885) balkanise l'Afrique. Le Cameroun, \emph{Kamerun} allemand, est partagé en 1919 entre la France (4/5) et la Grande-Bretagne (1/5) sous mandat SDN puis tutelle ONU. Cette double colonisation laisse un double héritage (droit, langue, administration). L'indépendance (1er janv. 1960 pour le Cameroun français) et la réunification (1er oct. 1961, plébiscite ONU au British Southern Cameroons) forgent une identité nationale fragile que le panafricanisme doit aider à consolider sans la dissoudre.
\end{savaistu}

\courssub{Les fondements idéologiques du consciencisme}

\begin{propriete}[Les trois piliers idéologiques (Nkrumah, \emph{Consciencism}, chap. 3--5)]
\begin{enumerate}
    \item \textbf{Matérialisme dialectique} : Outil d'analyse scientifique de la société. L'impérialisme est le stade suprême du capitalisme (Lénine) ; l'Afrique est maintenue en position de \cle{périphérie} (fournisseur de matières premières, marché captif). $\rightarrow$ Exige une \cle{planification continentale} (industrialisation, marché commun, monnaie unique) pour rompre la division internationale du travail.
    \item \textbf{Humanisme traditionnel africain} : \emph{« Je suis parce que nous sommes »} (Ubuntu). La communauté prime sur l'individu ; la propriété collective de la terre ; consensus décisionnel. $\rightarrow$ Base éthique d'un socialisme africain non aliénant, refusant l'individualisme capitaliste et le totalitarisme d'État.
    \item \textbf{Christianisme} : Non comme dogme théologique, mais comme \cle{révolution sociale} (libération des opprimés, éthique de responsabilité). Nkrumah retient le message de justice et d'égalité (Actes des Apôtres : mise en commun des biens).
\end{enumerate}
\textbf{Synthèse} : Le consciencisme = \emph{Matérialisme dialectique (cerveau) + Humanisme traditionnel (cœur) + Éthique chrétienne (conscience)}. Il rejette le capitalisme néocolonial et le marxisme-léninisme dogmatique (athéisme d'État, collectivisme forcé).
\end{propriete}

\courssub{Les fondements philosophiques du consciencisme}

\begin{aretenir}[Catégories philosophiques nkrumahistes]
\begin{itemize}
    \item \textbf{Ontologie} : L'être africain est \emph{devenir} (processus historique), non essence figée. La matière est première mais l'esprit (conscience) la transforme.
    \item \textbf{Epistémologie} : Connaissance = \cle{pratique révolutionnaire} (théorie + action). Le savoir doit servir la libération.
    \item \textbf{Éthique politique} : La fin (libération totale) justifie les moyens (État fort, parti unique \emph{transitoire}, planification). Risque dérive autoritaire (critique de Wiredu, Gyekye).
    \item \textbf{Unité} : Non juxtaposition d'États, mais \cle{fédération continentale} (Gouvernement continental, Armée commune, Monnaie unique, Citoyenneté africaine). Condition \emph{sine qua non} de la viabilité économique.
\end{itemize}
\end{aretenir}

\courssub{Historique du panafricanisme : des précurseurs à la ZLECAf}

\begin{exemplebox}[Les quatre phases du panafricanisme (Repères pour le Probatoire)]
\begin{center}
\begin{tabularx}{\linewidth}{|l|X|X|}\hline
\rowcolor{popblueL}
\textbf{Phase} & \textbf{Période / Acteurs clés} & \textbf{Objectifs / Réalisations} \\ \hline
\textbf{1. Diaspora \& Précurseurs} & 1900--1945 \newline Sylvester Williams, W.E.B. Du Bois, Marcus Garvey, George Padmore & Congrès panafricains (Londres 1900, Paris 1919, Londres 1921/23, New York 1927, Manchester 1945). Revendication : droits civiques, fin du colonialisme. \textbf{Manchester 1945} : basculement vers l'action de masse sur le continent (Nkrumah, Kenyatta, Nyerere). \\ \hline
\textbf{2. Indépendances \& Blocs rivaux} & 1957--1963 \newline Nkrumah (Casablanca) vs Houphouët-Boigny/Ahidjo (Monrovia) & \textbf{Groupe de Casablanca} (1961) : Unité politique \emph{immédiate} (État fédéral, armée commune). \textbf{Groupe de Monrovia} (1961) : Coopération économique \emph{progressive}, respect souveraineté. \textbf{Compromis d'Addis-Abeba (1963)} : Création OUA (Charte : souveraineté, non-ingérence, décolonisation). \\ \hline
\textbf{3. OUA : Souveraineté \& Décolonisation} & 1963--1999 & Lutte contre apartheid, soutien mouvements libération (Angola, Mozambique, Namibie, Zimbabwe). \textbf{Échec économique} : Traité de Lagos (CEDEAO 1975), Traité d'Abuja (Communauté Économique Africaine 1991) peu appliqués. Principe non-ingérence empêche gestion conflits internes. \\ \hline
\textbf{4. Union Africaine : Intégration \& Développement} & 2002--aujourd'hui \newline Kadhafi, Mbeki, Obasanjo $\rightarrow$ Agenda 2063 & \textbf{Sirte 1999 / Durban 2002} : OUA $\rightarrow$ UA. Principes : droit d'ingérence (crimes de guerre), intégration économique, paix/sécurité. \textbf{Outils} : NEPAD, PIDA (infrastructures), \textbf{ZLECAf} (Zone de Libre-Échange Continentale Africaine, entrée en vigueur mai 2019, échanges janv. 2021). Objectif : Marché unique 1,4 Md consommateurs, PIB \$3,4 Md. \\ \hline
\end{tabularx}
\end{center}
\end{exemplebox}

\begin{popfigure}
\centering
\begin{tikzpicture}[scale=0.85, every node/.style={font=\footnotesize}]
    % Timeline axis
    \draw[popdark, thick] (0,0) -- (18,0);
    % Years scale: 1900 to 2021 = 121 years. 18cm width. 1 year = 0.1487 cm.
    % Key dates positions (cm from 1900):
    % 1900:0, 1945:6.7, 1957:8.5, 1961:9.0, 1963:9.3, 1975:11.1, 1991:13.5, 1999:14.7, 2002:15.1, 2018:17.5, 2021:18.0
    % Phase backgrounds
    \fill[poporange!20] (0,-0.5) rectangle (6.7, 2.5);
    \fill[popgreen!20] (6.7,-0.5) rectangle (9.3, 2.5);
    \fill[popblue!20] (9.3,-0.5) rectangle (14.7, 2.5);
    \fill[poppurple!20] (14.7,-0.5) rectangle (18, 2.5);
    % Phase labels
    \node[poporange, rotate=90, anchor=south] at (0.3, 1.0) {Phase 1 : Diaspora};
    \node[popgreen, rotate=90, anchor=south] at (7.0, 1.0) {Phase 2 : Blocs};
    \node[popblue, rotate=90, anchor=south] at (10.0, 1.0) {Phase 3 : OUA};
    \node[poppurple, rotate=90, anchor=south] at (15.5, 1.0) {Phase 4 : UA / ZLECAf};
    % Events
    \def\ev#1#2#3{ % pos, label, desc
        \draw[popdark] (#1,0) -- (#1,-0.3);
        \node[anchor=north, align=center, text width=2.2cm] at (#1,-0.5) {\textbf{#2}\\\emph{#3}};
    }
    \ev{0}{1900}{Conf. panafricaine\\Londres (Williams, Du Bois)}
    \ev{6.7}{1945}{5\textsuperscript{e} Congrès\\Manchester (Nkrumah, Kenyatta)}
    \ev{8.5}{1957}{Indépendance\\Ghana (Nkrumah)}
    \ev{9.0}{1961}{Conf. Casablanca /\newline Monrovia}
    \ev{9.3}{1963}{Création OUA\\Addis-Abeba}
    \ev{11.1}{1975}{Traité Lagos\\CEDEAO}
    \ev{13.5}{1991}{Traité Abuja\\CEA}
    \ev{14.7}{1999}{Déclaration Sirte\\OUA $\rightarrow$ UA}
    \ev{15.1}{2002}{Lancement UA\\Durban}
    \ev{17.5}{2018}{Signature ZLECAf\\Kigali}
    \ev{18.0}{2021}{Début échanges\\ZLECAf}
    % Scale bar
    \draw[popdark, thick] (0,-1.8) -- (18,-1.8);
    \foreach \y/\x in {1900/0, 1950/7.4, 2000/14.9, 2021/18} {
        \draw[popdark] (\x,-1.8) -- (\x,-1.9);
        \node[anchor=north] at (\x,-1.95) {\x};
    }
    \node[anchor=north] at (9,-2.3) {Échelle : 1 cm $\approx$ 6,7 ans (121 ans = 18 cm)};
\end{tikzpicture}
\legende{Frise chronologique synthétique du panafricanisme (1900--2021). Les phases sont colorées : \textcolor{poporange}{■} Diaspora/Précurseurs, \textcolor{popgreen}{■} Indépendances/Blocs, \textcolor{popblue}{■} OUA, \textcolor{poppurple}{■} UA/ZLECAf.}
\end{popfigure}

\courssub{Unité africaine et développement : Applications concrètes au Cameroun}

\begin{popfigure}
\centering
\begin{tikzpicture}[scale=1.1, every node/.style={font=\scriptsize}]
    % Simplified map coordinates (approximate relative positions)
    % Cameroon center ~ (0,0). North up.
    % Neighbors: Chad N, CAR E, Congo S, Gabon SW, Eq. Guinea W.
    % Country shapes (very simplified polygons)
    \draw[popdark, thick, fill=popblue!10] (-2.5, 1.5) -- (2.5, 1.5) -- (1.5, -2) -- (-1.5, -2) -- cycle; % Cameroon
    \node[popblue, align=center] at (0, 0.2) {\textbf{CAMEROUN}\\(Yaoundé)};
    \draw[popdark, thick, fill=poporange!10] (-2, 3.5) -- (2, 3.5) -- (1.5, 1.5) -- (-1.5, 1.5) -- cycle; % Chad
    \node[poporange, align=center] at (0, 2.8) {\textbf{TCHAD}\\(N'Djaména)};
    \draw[popdark, thick, fill=popgreen!10] (1.5, 1) -- (4, 1) -- (3.5, -2.5) -- (1.5, -2) -- cycle; % CAR
    \node[popgreen, align=center] at (2.8, -0.5) {\textbf{RCA}\\(Bangui)};
    \draw[popdark, thick, fill=poppurple!10] (-1, -2) -- (3, -2) -- (2.5, -4) -- (-1.5, -4) -- cycle; % Congo
    \node[poppurple, align=center] at (0.5, -3) {\textbf{CONGO}\\(Brazzaville)};
    \draw[popdark, thick, fill=poppink!10] (-2.5, -2) -- (0, -2) -- (-0.5, -3.5) -- (-2.5, -3.5) -- cycle; % Gabon (partial)
    \node[poppink, align=center] at (-1.3, -2.7) {\textbf{GABON}\\(Libreville)};
    \draw[popdark, thick, fill=popgold!10] (-3, 0.5) -- (-1.5, 0.5) -- (-1.5, -1) -- (-3, -1) -- cycle; % Eq Guinea
    \node[popgold, align=center] at (-2.2, -0.2) {\textbf{GUINÉE ÉQUATORIALE}\\(Malabo)};
    % Ports
    \node[anchor=east, popblue] at (-2.6, -0.5) {\textbullet\ \textbf{Douala}};
    \node[anchor=north, popblue] at (-1.5, -2.1) {\textbullet\ \textbf{Kribi (eau profonde)}};
    \node[anchor=south, popblue] at (-2.5, 1.6) {\textbullet\ \textbf{Limbé}};
    % Corridors (Double lines)
    \draw[popdark, double distance=2pt, thick] (0,0.2) -- (0, 2.8); % Douala - N'Djamena
    \node[popdark, rotate=90, anchor=south] at (0.3, 1.5) {Corridor Douala--N'Djaména};
    \draw[popdark, double distance=2pt, thick] (0,0.2) -- (2.8, -0.5); % Douala - Bangui
    \node[popdark, rotate=30, anchor=south] at (1.4, 0) {Corridor Douala--Bangui};
    % Export Flows Cameroon -> CEMAC (Medium arrows)
    \draw[popblue, thick, -{Stealth[length=3mm]}, double distance=1pt] (-2.6, -0.5) -- (-2.2, 2.8); % Douala -> Chad (Petroleum products via pipeline/service)
    \draw[popblue, thick, -{Stealth[length=3mm]}, double distance=1pt] (-2.6, -0.5) -- (2.8, -0.5); % Douala -> CAR (Agro-industry)
    \draw[popblue, thick, -{Stealth[length=3mm]}, double distance=1pt] (-2.6, -0.5) -- (0.5, -3); % Douala -> Congo
    \draw[popblue, thick, -{Stealth[length=3mm]}, double distance=1pt] (-2.6, -0.5) -- (-1.3, -2.7); % Douala -> Gabon
    \draw[popblue, thick, -{Stealth[length=3mm]}, double distance=1pt] (-2.6, -0.5) -- (-2.2, -0.2); % Douala -> Eq Guinea
    % Import Flows CEMAC -> Cameroon (Thin arrows)
    \draw[poporange, thick, -{Stealth[length=2mm]}] (-2.2, 2.8) -- (-2.6, -0.5); % Chad -> Douala (Cattle, Cotton)
    \draw[popgreen, thick, -{Stealth[length=2mm]}] (2.8, -0.5) -- (-2.6, -0.5); % CAR -> Douala (Timber, Cotton)
    \draw[poppurple, thick, -{Stealth[length=2mm]}] (0.5, -3) -- (-2.6, -0.5); % Congo -> Douala (Timber)
    \draw[poppink, thick, -{Stealth[length=2mm]}] (-1.3, -2.7) -- (-2.6, -0.5); % Gabon -> Douala (Timber)
    % Legend
    \begin{scope}[shift={(4.5, 2.5)}]
        \node[anchor=west] at (0,0) {\textbf{Légende}};
        \draw[popdark, thick, fill=popblue!10] (0,-0.4) rectangle (0.4,0); \node[anchor=west] at (0.5,-0.2) {Cameroun};
        \draw[popdark, thick, fill=poporange!10] (0,-0.8) rectangle (0.4,-0.4); \node[anchor=west] at (0.5,-0.6) {Tchad};
        \draw[popdark, thick, fill=popgreen!10] (0,-1.2) rectangle (0.4,-0.8); \node[anchor=west] at (0.5,-1.0) {RCA};
        \draw[popdark, thick, fill=poppurple!10] (0,-1.6) rectangle (0.4,-1.2); \node[anchor=west] at (0.5,-1.4) {Congo};
        \draw[popdark, thick, fill=poppink!10] (0,-2.0) rectangle (0.4,-1.6); \node[anchor=west] at (0.5,-1.8) {Gabon};
        \draw[popdark, thick, fill=popgold!10] (0,-2.4) rectangle (0.4,-2.0); \node[anchor=west] at (0.5,-2.2) {Guinée Éq.};
        \draw[popblue, double distance=1pt, -{Stealth[length=2mm]}] (0,-2.8) -- (1.5,-2.8); \node[anchor=west] at (1.6,-2.8) {Exports Cameroun (8,5\% total) : Pétrole, Agro-ind.};
        \draw[poporange, -{Stealth[length=2mm]}] (0,-3.2) -- (1.5,-3.2); \node[anchor=west] at (1.6,-3.2) {Imports CEMAC $\rightarrow$ Cameroun : Bois, Bétail, Coton};
        \draw[popdark, double distance=2pt] (0,-3.6) -- (1.5,-3.6); \node[anchor=west] at (1.6,-3.6) {Corridors routiers structurants};
        \node[anchor=west] at (0,-4.0) {\textbullet\ Ports maritimes};
    \end{scope}
    % North Arrow
    \node at (-3.5, 3.5) {\begin{tikzpicture}\draw[popdark, thick, -{Stealth[length=4mm]}] (0,0) -- (0,0.8); \node[anchor=south] at (0,0.8) {N};\end{tikzpicture}};
    
    % Scale
    \draw[popdark, thick] (-3.5, -4.5) -- (-2, -4.5);
    \node[anchor=north] at (-2.75, -4.7) {$\approx$ 500 km};
\end{tikzpicture}
\legende{Croquis schématique : Le Cameroun dans l'espace CEMAC (2022). Flèches bleues doubles = exports camerounais vers CEMAC (8,5\% des exports totaux : produits pétroliers, farine, savon, bière, ciment). Flèches fines colorées = imports depuis pays CEMAC (bois, bétail, coton). Traits doubles noirs = corridors Douala--N'Djaména et Douala--Bangui. Sources : Douanes Cameroun, BEAC, CNUCED 2023.}
\end{popfigure}

\courssub{Méthode : Rédiger et justifier une démarche argumentative}

\begin{methode}[Dissertation / Commentaire type Probatoire : 4 étapes]
\begin{enumerate}
    \item \textbf{Problématiser} : Analyser le sujet (termes, présupposés, tension). Formuler une \cle{problématique} (question philosophique ouverte). Exemple : \emph{« L'unité africaine est-elle le levier principal du développement ou une condition nécessaire mais insuffisante sans rupture structurelle interne ? »}
    \item \textbf{Théser} : Prendre position clairement (\emph{Thèse}). Exemple : \emph{« Le panafricanisme (ZLECAf) est une condition nécessaire au développement camerounais car il élargit le marché et permet la planification continentale (consciencisme), mais il reste insuffisant sans transformation de la gouvernance et des infrastructures nationales. »}
    \item \textbf{Argumenter} (3 à 4 arguments) : Structure par paragraphe : \textbf{Idée directrice} + \textbf{Explication philosophique} (notion, auteur) + \textbf{Preuve concrète} (chiffre, document, exemple camerounais). Utiliser connecteurs logiques (\emph{En premier lieu, De surcroît, Cependant, En définitive}).
    \item \textbf{Conclure} : \textbf{Bilan synthétique} (réponse à la problématique) + \textbf{Ouverture} (perspective, nouvelle question, limite). Ne pas introduire de nouvel argument.
\end{enumerate}
\end{methode}

\begin{attention}[Pièges fréquents au Probatoire]
\begin{itemize}
    \item \textbf{Confusion Consciencisme / Panafricanisme} : Le consciencisme est la \emph{philosophie fondatrice} (Nkrumah) ; le panafricanisme est le \emph{mouvement historique et institutionnel} (OUA/UA/ZLECAf).
    \item \textbf{Anachronisme} : Ne pas prêter à Nkrumah la ZLECAf (2018). Dire : \emph{« La ZLECAf réalise l'exigence nkrumahiste de planification continentale... »}
    \item \textbf{Angélisme} : Ne pas dire « l'unité suffit ». Nuancer : souveraineté vs intégration, rente extractive vs transformation, gouvernance.
    \item \textbf{Données non sourcées} : Toujours citer la source (Doc 3, ONCC 2023, CNUCED 2022).
\end{itemize}
\end{attention}

\courssub{Activité d'intégration : Le Cameroun face à la ZLECAf (2 séances)}

\courssub{Objectifs de l'activité}
\begin{itemize}
    \item Mobiliser les connaissances sur le \cle{consciencisme} (fondements idéologiques et philosophiques) et l'\cle{historique du panafricanisme} pour analyser un corpus documentaire hétérogène.
    \item Construire des outils de synthèse visuelle : une \cle{frise chronologique} (1900--2021) et un \cle{croquis géographique} du Cameroun dans l'espace \cle{CEMAC} avec flux commerciaux.
    \item Exercer la \cle{démarche argumentative} : formuler une thèse, la soutenir par au moins trois arguments tirés des documents et du cours, et rédiger une conclusion justifiée (15--20 lignes).
    \item Appliquer les notions d'unité africaine à une situation concrète de développement (ZLECAf, transformation locale du cacao au Cameroun).
\end{itemize}

\courssub{Situation}
\textbf{Contexte authentique.} En janvier 2024, le Cameroun célèbre la première année de mise en œuvre effective des échanges sous la \cle{ZLECAf} (Zone de libre-échange continentale africaine, entrée en vigueur mai 2019, début des échanges janvier 2021). À Douala, le port autonome enregistre une hausse de 12 \% du trafic de conteneurs vers le Nigeria et le Gabon. Pourtant, dans les plantations du Sud (Ebolowa, Sangmélima), les planteurs de cacao vendent encore 85 \% de leur production à l'état brut vers l'Europe (données ONCC 2023). L'usine de transformation de \emph{Chococam} (Douala) et la nouvelle unité \emph{Neo Industry} (Kribi) peinent à s'approvisionner localement à cause de la logistique et des normes sanitaires.

\textbf{Dossier documentaire (quatre documents).}

\begin{enumerate}
    \item \textbf{Extrait de \emph{Consciencism} (Kwame Nkrumah, 1964, chap. 5).} « L'unité africaine n'est pas un rêve romantique ; c'est une nécessité pratique pour la libération économique. Tant que l'Afrique sera balkanisée en petits États non viables, elle restera le fournisseur de matières premières bon marché et le marché captif des produits manufacturés d'autrui. Le consciencisme exige une planification continentale : une seule monnaie, un marché commun, une défense commune, une citoyenneté africaine. »
    
    \item \textbf{Carte schématique : Groupes de Casablanca et de Monrovia (1961).} La carte montre l'Afrique divisée en deux blocs : le \textbf{Groupe de Casablanca} (Ghana, Guinée, Mali, Maroc, Égypte, Algérie provisoire/GPRA, Libye) prônant l'unité politique immédiate (État fédéral) ; le \textbf{Groupe de Monrovia} (Nigeria, Libéria, Sierra Leone, Éthiopie, Somalie, Togo, Cameroun, etc.) privilégiant la coopération économique progressive et la souveraineté nationale. Le Cameroun (indépendance 1960, réunification 1961) adhère au Groupe de Monrovia. La carte indique les capitales et les flèches des premières conférences (Casablanca janv. 1961, Monrovia mai 1961, Addis-Abeba 1963).
    
    \item \textbf{Données chiffrées : Commerce intra-régional (Source : CNUCED, Rapport 2023 ; BCEAO/BEAC 2022).}
    \begin{center}
    \begin{tabularx}{\linewidth}{|l|X|X|X|}\hline
    \rowcolor{popblueL}
    \textbf{Zone} & \textbf{Part du commerce intra-zonal (\% du total exports+imports)} & \textbf{Produits dominants} & \textbf{Année de référence} \\ \hline
    Afrique (total continental) & 15 \% & Matières premières (pétrole, minerais, cacao, café), faible transformation & 2022 \\ \hline
    Union européenne (UE-27) & 60 \% & Produits manufacturés, services, haute technologie & 2022 \\ \hline
    CEMAC (6 pays) & 4,2 \% & Bois, pétrole, coton, bétail ; barrières non tarifaires élevées & 2022 \\ \hline
    Cameroun vers CEMAC & 8,5 \% des exports camerounais & Produits agro-industriels (farine, savon, bière), hydrocarbures & 2022 \\ \hline
    \end{tabularx}
    \end{center}
    
    \item \textbf{Article de presse : « Cacao camerounais : la ZLECAf peut-elle briser la malédiction de l'export brut ? » (\emph{La Nouvelle Expression}, 12 mars 2024).} L'article rapporte que le Cameroun, 5\textsuperscript{e} producteur mondial (300 000 t/an), ne transforme localement que 18 \% de ses fèves (objectif officiel 50 \% à l'horizon 2030). La ZLECAf prévoit l'élimination des droits de douane sur 90 \% des lignes tarifaires d'ici 2030. Le directeur de \emph{Neo Industry} (Kribi) témoigne : « Nous importons du beurre de cacao de Côte d'Ivoire et du Ghana car la logistique intérieure (routes, électricité) coûte plus cher que le fret maritime d'Abidjan. » L'article mentionne le \emph{Programme de développement des chaînes de valeur cacao} (PDCVA) soutenu par la BAD (45 milliards FCFA sur 5 ans).
\end{enumerate}

\courssub{Consignes}
\textbf{Travail en groupes (4 élèves) – Durée : 2 séances (2 h).}

\begin{enumerate}
    \item \textbf{Frise chronologique (1900--2021).} Sur une feuille A3 paysage, construisez une frise à l'échelle (1 cm $\approx$ 6,7 ans / 121 ans = 18 cm). Positionnez et datez \textbf{10 événements majeurs} du panafricanisme : Conférence panafricaine de Londres (1900), 5\textsuperscript{e} Congrès de Manchester (1945), Indépendance du Ghana (1957), Conférences de Casablanca/Monrovia (1961), Création de l'OUA (1963), Charte de Lagos/CEDEAO (1975), Traité d'Abuja/CEA (1991), Sirte/OUA$\rightarrow$UA (1999/2002), Lancement ZLECAf (2018), Entrée en vigueur ZLECAf (2019) / Début échanges (2021). Ajoutez une bande colorée pour chaque « phase » (pré-indépendance, indépendances/blocs, OUA, UA/intégration économique).
    
    \item \textbf{Croquis : Le Cameroun dans la CEMAC – Flux commerciaux.} Sur un fond de carte simplifié du Cameroun et des 5 pays voisins de la CEMAC (Tchad, RCA, Congo, Guinée équatoriale, Gabon), représentez :
    \begin{itemize}
        \item Les 3 grands ports (Douala, Kribi, Limbé) et les 2 corridors routiers (Douala–N'Djaména ; Douala–Bangui).
        \item Les flux d'exportation du Cameroun vers la CEMAC (flèches proportionnelles : 8,5 \% du total exports) : produits pétroliers, agro-industrie.
        \item Les flux d'importation depuis la CEMAC vers le Cameroun (bois, bétail, coton).
        \item Une légende organisée (signes ponctuels, linéaires, surfaciques) et une flèche Nord.
    \end{itemize}
    
    \item \textbf{Analyse guidée des documents (réponses écrites dans le cahier).}
    \begin{enumerate}
        \item Doc 1 : Identifiez la thèse centrale de Nkrumah. Quel concept du consciencisme (matérialisme dialectique, humanisme traditionnel, christianisme) sous-tend l'exigence d'une « planification continentale » ?
        \item Doc 2 : Expliquez pourquoi le Cameroun a choisi le Groupe de Monrovia en 1961. Quel lien avec la réunification de 1961 ?
        \item Doc 3 : Calculez le ratio \emph{Commerce intra-africain / Commerce intra-européen}. Interprétez l'écart au regard de la thèse du Doc 1.
        \item Doc 4 : Relevez deux obstacles concrets à la transformation locale du cacao cités dans l'article. En quoi la ZLECAf peut-elle lever (ou non) ces obstacles ?
    \end{enumerate}
    
    \item \textbf{Rédaction individuelle (évaluation).} Rédigez une conclusion argumentée de \textbf{15 à 20 lignes} répondant à la problématique : \emph{« Le panafricanisme est-il encore utile au développement du Cameroun ? »} Votre texte doit comporter :
    \begin{itemize}
        \item Une introduction (accroche, définition du panafricanisme, thèse claire).
        \item Au moins \textbf{trois arguments distincts} justifiés par des références explicites aux documents (Doc 1 à 4) et aux notions du cours (consciencisme, étapes historiques, ZLECAf).
        \item Une conclusion ouverte (perspective).
    \end{itemize}
    
    \item \textbf{Restitution orale.} Chaque groupe désigne un rapporteur (3 min) pour présenter sa frise, son croquis et sa synthèse des arguments. Débat collectif animé par l'enseignant.
\end{enumerate}

\courssub{Ressources à mobiliser}
\begin{itemize}
    \item \textbf{Consciencisme (Nkrumah) :} Synthèse de trois courants : \cle{matérialisme dialectique} (analyse économique, lutte des classes), \cle{humanisme traditionnel africain} (communautarisme, Ubuntu), \cle{christianisme} (éthique de la charité, justice sociale). $\rightarrow$ Exigence d'unité politique \emph{et} économique.
    \item \textbf{Phases du panafricanisme :} (1) Diaspora/Précurseurs (1900--1945) ; (2) Indépendances/Blocs rivaux (1957--1963) ; (3) OUA : souveraineté/décolonisation (1963--1999) ; (4) Union Africaine : intégration économique, paix, ZLECAf (2002--aujourd'hui).
    \item \textbf{Indicateurs de développement :} Taux de transformation locale, part du commerce intra-régional, IDH, balance commerciale.
    \item \textbf{Méthode dissertation/commentaire :} \emph{Problématiser $\rightarrow$ Théser $\rightarrow$ Argumenter (exemple, citation, chiffre) $\rightarrow$ Conclure.}
\end{itemize}

\courssub{Démarche et corrigé commenté}
\begin{exoresolu}{Corrigé type de l'activité d'intégration}
\textbf{1. Frise chronologique (attendus).} 
\begin{itemize}
    \item \textbf{Echelle respectée :} 1900--2021 = 121 ans $\approx$ 18 cm. Repères : 1900 (0 cm), 1945 (6,7 cm), 1957 (8,5 cm), 1961 (9,0 cm), 1963 (9,3 cm), 1975 (11,1 cm), 1991 (13,5 cm), 1999 (14,7 cm), 2002 (15,1 cm), 2018 (17,5 cm), 2021 (18,0 cm).
    \item \textbf{Phases colorées :} \textcolor{poporange}{1900--1945 : Diaspora} ; \textcolor{popgreen}{1957--1963 : Indépendances/Blocs} ; \textcolor{popblue}{1963--1999 : OUA} ; \textcolor{poppurple}{2002--... : UA/ZLECAf}.
    \item \textbf{Justification :} Le 5\textsuperscript{e} Congrès de Manchester (1945) marque le basculement vers l'action de masse sur le continent. La double conférence 1961 (Casablanca/Monrovia) illustre la fracture stratégique résolue à Addis-Abeba (1963) par le compromis OUA.
\end{itemize}

\textbf{2. Croquis Cameroun/CEMAC (attendus).}
\begin{itemize}
    \item \textbf{Fond de carte :} Contour simplifié du Cameroun (triangle pointe au Nord), pays voisins positionnés relativement (Tchad N, RCA E, Congo S, Gabon S-O, Guinée équatoriale O).
    \item \textbf{Signes :} Ports (ancres) : Douala (principal), Kribi (eau profonde), Limbé. Corridors (traits doubles) : Douala–N'Djaména (N'Djaména au Nord-Est), Douala–Bangui (Bangui à l'Est).
    \item \textbf{Flux :} Flèches sortantes Cameroun $\rightarrow$ CEMAC (épaisseur moyenne) : hydrocarbures (pipeline Tchad–Kribi compté comme export service), farine, savon, bière, ciment. Flèches entrantes CEMAC $\rightarrow$ Cameroun (fines) : bois (Congo, Gabon), bétail (Tchad, RCA), coton (Tchad).
    \item \textbf{Légende :} Titre « Cameroun dans la CEMAC : flux commerciaux 2022 », échelle graphique, Nord, sources (Douanes, BEAC).
\end{itemize}

\textbf{3. Analyse guidée (réponses types).}
\begin{enumerate}
    \item \textbf{Doc 1 :} Thèse : L'unité africaine est une \emph{nécessité économique} (sortir de la division balkanisée) et non un idéal abstrait. Concept mobilisé : le \cle{matérialisme dialectique} du consciencisme (analyse des rapports de production, impérialisme = stade suprême du capitalisme) combiné à l'\cle{humanisme traditionnel} (solidarité communautaire). La « planification continentale » = application du matérialisme à l'échelle africaine.
    \item \textbf{Doc 2 :} Le Cameroun (réunifié en 1961 après plébiscite ONU) privilégie la \cle{souveraineté nationale} et la \cle{coopération progressive} (Groupe de Monrovia) par peur qu'un fédéralisme immédiat (Casablanca) ne dilue son identité fragile (double héritage franco-britannique) et ne favorise une hégémonie ghanéenne/nigériane. Ahidjo et Nkrumah s'opposent sur la méthode : État unitaire fort vs États-Unis d'Afrique.
    \item \textbf{Doc 3 :} Ratio = $\num{15} / \num{60} = \num{0,25}$ (soit 4 fois moins). L'écart confirme Nkrumah : l'Afrique trade 4 fois moins avec elle-même que l'Europe, preuve de la « balkanisation » économique. Le chiffre CEMAC (4,2 \%) montre l'échec régional malgré l'union monétaire (FCFA).
    \item \textbf{Doc 4 :} Obstacles : (1) \emph{Logistique intérieure défaillante} (routes, électricité) rendant le coût local $>$ coût import ; (2) \emph{Normes sanitaires/qualité} non harmonisées. ZLECAf : levier potentiel sur les droits de douane (90 \% lignes à 0 \%) et harmonisation normes (protocole commerce), mais \emph{insuffisante} sans investissements infrastructures (corridors) et énergie (barrages Nachtigal, Lom Pangar).
\end{enumerate}

\textbf{4. Rédaction individuelle – Proposition de corrigé (18 lignes).}
\textit{« Le panafricanisme reste indispensable au développement du Cameroun, non comme slogan mais comme levier structurel. D'abord, le consciencisme de Nkrumah (Doc 1) démontre que la balkanisation économique (Doc 3 : 15 \% commerce intra-africain vs 60 \% UE) maintient le Cameroun en position de fournisseur de cacao brut (Doc 4 : 82 \% non transformé). L'unité continentale via la ZLECAf offre un marché de 1,4 milliard de consommateurs pour absorber la transformation locale visée à 50 \%. Ensuite, l'histoire du panafricanisme (Doc 2) montre que le choix monrovien du Cameroun (souveraineté) n'exclut pas l'intégration progressive : la CEMAC (malgré 4,2 \% d'échanges) et la ZLECAf en sont les héritières institutionnelles. Enfin, les obstacles concrets (logistique, énergie, normes) soulignés par l'industriel de Kribi (Doc 4) ne seront levés que par une planification continentale (routes transafricaines, harmonisation normes) que seul le panafricanisme institutionnalisé (UA, PIDA) peut porter. Cependant, l'utilité du panafricanisme suppose un État camerounais capable de négocier ses intérêts (clause de sauvegarde, protection industries naissantes) et d'investir massivement dans l'humain et l'infrastructure. Le panafricanisme est donc une condition nécessaire, non suffisante : il faut y ajouter la gouvernance démocratique et la rupture avec la rente extractive. »}

\textbf{5. Grille d'évaluation (indicative).}
\begin{center}
\begin{tabularx}{\linewidth}{|l|X|X|}\hline
\rowcolor{popblueL}
\textbf{Critère} & \textbf{Indicateurs de réussite} & \textbf{Points} \\ \hline
Frise (exactitude, échelle, phases) & 10 événements bien datés, 4 phases colorées, lisibilité & 4 \\ \hline
Croquis (localisation, flux, légende) & 3 ports, 2 corridors, flèches proportionnelles, légende complète & 4 \\ \hline
Analyse docs (justesse, vocabulaire cours) & 4 réponses correctes, termes : consciencisme, matérialisme, souveraineté, ZLECAf & 4 \\ \hline
Rédaction (structure, 3 arguments, 15--20 lignes) & Intro/Thèse, 3 arguments doc+cours, conclusion ouverte, langue soignée & 8 \\ \hline
Restitution orale & Synthèse claire, réponse aux questions, esprit critique & 5 \\ \hline
\textbf{Total} & & \textbf{25} \\ \hline
\end{tabularx}
\end{center}
\end{exoresolu}

\courssub{Bilan de l'activité}
\begin{aretenir}[Bilan pédagogique]
Cette activité d'intégration a permis de :
\begin{itemize}
    \item \textbf{Ancrer les concepts abstraits} (consciencisme, panafricanisme) dans des données réelles camerounaises (cacao, ZLECAf, corridors CEMAC), rendant la philosophie opératoire pour l'analyse de développement.
    \item \textbf{Articuler histoire et actualité :} la frise montre la continuité (Manchester $\rightarrow$ ZLECAf) et la rupture (blocs 1961 $\rightarrow$ compromis OUA $\rightarrow$ UA économique).
    \item \textbf{Développer des compétences transverselles :} cartographie (croquis flux), traitement de données (ratio $\num{0,25}$), lecture critique de presse, rédaction argumentée contrainte.
    \item \textbf{Révéler les tensions réelles :} souveraineté vs intégration (choix Monrovia 1961), rente extractive vs transformation locale (cacao), infrastructures physiques vs cadres juridiques (ZLECAf nécessaire mais insuffisante sans routes/énergie).
    \item \textbf{Préparer l'épreuve du Probatoire :} dissertation type « Sujet d'actualité philosophique » ou commentaire de texte (extrait Nkrumah), avec exigence de thèse, d'arguments documentés et de conclusion nuancée.
\end{itemize}
L'enseignant veillera, lors de la correction collective, à souligner que le \cle{consciencisme} n'est pas une idéologie figée mais une \emph{méthode d'analyse} (matérialisme + humanisme) pour penser l'unité africaine aujourd'hui, et que le \cle{panafricanisme} est un \emph{processus historique} dont la ZLECAf est l'actuelle traduction juridique et économique.
\end{aretenir}

\newpage

\coursec{Méthodes et exercices résolus}

\courssub{Méthodes et exercices résolus}

\begin{methode}[Analyser et définir un concept philosophique africain (Consciencisme, Panafricanisme)]
\textbf{Objectif :} Maîtriser la définition rigoureuse des notions du programme pour les mobiliser en dissertation ou en commentaire.
\begin{enumerate}
    \item \textbf{Identifier le genre et l'espèce :} De quel type de notion s'agit-il ? (Idéologie, philosophie, mouvement politique, concept métaphysique).
    \item \textbf{Circonscrire l'étymologie et le sens courant :} D'où vient le mot ? Que signifie-t-il dans le langage commun ?
    \item \textbf{Exposer la définition technique (auteur de référence) :} Citer explicitement \cle{Kwame Nkrumah} pour le consciencisme ; \cle{Du Bois}, \cle{Padmore}, \cle{Nkrumah}, \cle{Nyerere} pour le panafricanisme.
    \item \textbf{Dégager les attributs essentiels (critères) :} Quelles sont les conditions nécessaires et suffisantes ? (Ex: Pour le consciencisme : synthèse tripartite, base matérialiste, visée humaniste).
    \item \textbf{Illustrer par un contre-exemple ou une distinction :} Différencier panafricanisme politique et panafricanisme culturel ; distinguer consciencisme et socialisme scientifique classique.
\end{enumerate}
\textbf{Réflexe Bac :} Toute définition doit être \emph{attribuée} (« Selon Nkrumah, le consciencisme est... ») et \emph{problématisée} (en quoi cette définition répond-elle à un problème historique ?).
\end{methode}

\begin{exoresolu}[Définir et problématiser le consciencisme]
\textbf{Énoncé :} « Le consciencisme n'est pas une simple synthèse éclectique, c'est une philosophie de l'action pour la décolonisation totale. » À partir de cette affirmation, définissez le consciencisme et montrez en quoi il constitue une réponse spécifique à la situation africaine post-coloniale.

\tcblower
\textbf{Solution détaillée}

\textbf{1. Analyse du sujet et délimitation}
\begin{itemize}
    \item \textbf{Concept clé :} Consciencisme (philosophie/idéologie de Kwame Nkrumah).
    \item \textbf{Citation à expliciter :} « Pas une simple synthèse éclectique » $\rightarrow$ rejette le syncrétisme superficiel. « Philosophie de l'action » $\rightarrow$ praxis, matérialisme. « Décolonisation totale » $\rightarrow$ politique, économique, mentale.
    \item \textbf{Problématique :} Comment le consciencisme, par sa structure philosophique originale, permet-il de penser l'unité et le développement africains au-delà de l'indépendance formelle ?
\end{itemize}

\textbf{2. Définition rigoureuse (Étape 1 à 3 de la méthode)}
Le \cle{consciencisme} est l'idéologie philosophique fondée par \cle{Kwame Nkrumah} (exposée dans \emph{Consciencism}, 1964). Il se définit comme une \textbf{synthèse dialectique} entre trois courants constitutifs de l'identité africaine contemporaine :
\begin{enumerate}
    \item L'héritage \textbf{traditionnel} (communalisme, humanisme, spiritualité non théiste).
    \item L'apport \textbf{islamique} (fraternité, discipline, monothéisme éthique).
    \item L'héritage \textbf{euro-chrétien/colonial} (technique, rationalité scientifique, droits de l'homme, mais aliénation).
\end{enumerate}
Cette synthèse repose sur une base \textbf{matérialiste dialectique} (héritage marxiste assumé comme outil d'analyse) mais refuse l'athéisme d'État : la matière est \emph{esprit en devenir}. La catégorie centrale est le \textbf{humanisme} : l'homme africain est la fin et le moyen du développement.

\textbf{3. Argumentation : En quoi est-ce une réponse à la situation post-coloniale ? (Étape 4-5)}

\textbf{Thèse 1 : Dépassement de l'aliénation culturelle (Contre l'éclectisme).}
L'éclectisme juxtapose sans hiérarchiser. Le consciencisme \emph{intègre} par la dialectique : il conserve le noyau humaniste de la tradition (valeur suprême : la communauté), critique l'aliénation coloniale (technique sans humanisme) et récupère l'universel (science, raison). \textbf{Exemple concret :} Au Cameroun, le bilinguisme officiel (héritage colonial) ne doit pas être une simple coexistence (éclectisme) mais une synthèse créant une identité nationale nouvelle (consciencisme appliqué).

\textbf{Thèse 2 : Fondement philosophique du développement endogène (Philosophie de l'action).}
Le matérialisme dialectique fournit l'outil d'analyse des structures néocoloniales (dépendance économique). Mais le but n'est pas la dictature du prolétariat, c'est l'\textbf{humanisme africain} (le « \cle{communalisme} » modernisé). L'action vise la \textbf{réhabilitation de la personnalité africaine} (Fanon, Senghor) par la maîtrise des sciences et techniques.
\textbf{Application Cameroun :} La politique d'import-substitution ou la ZLECAf (Zone de Libre-Échange Continentale Africaine) relèvent de cette « philosophie de l'action » : transformer la matière première localement (base matérialiste) pour l'épanouissement des populations (fin humaniste).

\textbf{Thèse 3 : Décolonisation totale (Mentale et Économique).}
L'indépendance politique (années 60) est incomplète sans décolonisation des mentalités (éducation, langues, valeurs) et des structures économiques (néocolonialisme). Le consciencisme lie \textbf{libération} et \textbf{création}.

\textbf{4. Conclusion}
Le consciencisme n'est pas un syncrétisme culturel mais une \textbf{arme théorique} : il donne à l'Afrique les catégories philosophiques (matérialisme dialectique + humanisme communal) pour penser son développement \emph{par elle-même et pour elle-même}, résolvant la tension entre tradition et modernité.

\textbf{Phrase de conclusion type :} « En synthétisant dialectiquement les trois âmes de l'Afrique sur une base matérialiste orientée vers l'humanisme, le consciencisme offre le cadre théorique indispensable à une décolonisation intégrale, condition \emph{sine qua non} d'un développement endogène et solidaire. »
\end{exoresolu}

\begin{methode}[Construire une explication historique structurée (Genèse du Panafricanisme)]
\textbf{Objectif :} Restituer l'historique d'un mouvement complexe sans faire un simple catalogue de dates, en articulant \textbf{contexte}, \textbf{acteurs} et \textbf{évolution des idées}.
\begin{enumerate}
    \item \textbf{Périodiser en grandes phases logiques (pas seulement chronologiques) :}
        \begin{itemize}
            \item Phase 1 : \textbf{Précurseurs / Diaspora} (XVIII\textsuperscript{e}-XIX\textsuperscript{e} s.) : Révolte contre l'esclavage, quête d'identité (\cle{Olaudah Equiano}, \cle{Marcus Garvey}).
            \item Phase 2 : \textbf{Institutionnalisation intellectuelle} (1900-1945) : Congrès panafricains (\cle{Du Bois}, \cle{Padmore}), passage du « retour en Afrique » à la « libération de l'Afrique ».
            \item Phase 3 : \textbf{Appropriation continentale / Indépendances} (1945-1963) : Congrès de Manchester (1945), rôle des \cle{Pères des Indépendances} (Nkrumah, Nyerere, Kenyatta, Houphouët-Boigny), création de l'OUA (1963).
            \item Phase 4 : \textbf{De l'OUA à l'UA / Renaissance} (1963-à nos jours) : Développement économique (NEPAD, UA 2002, ZLECAf 2018), panafricanisme populaire et culturel.
        \end{itemize}
    \item \textbf{Pour chaque phase :} Contexte historique $\rightarrow$ Acteurs clés $\rightarrow$ Thèses centrales $\rightarrow$ Formes d'organisation $\rightarrow$ Limites/Échecs.
    \item \textbf{Identifier les lignes de force (invariant) :} Unité politique, solidarité diaspora/continent, libération totale, développement endogène.
    \item \textbf{Problématiser les tensions internes :} \cle{Monrovia} (intégration progressive, souveraineté) vs \cle{Casablanca} (fédéralisme immédiat, rupture radicale) $\rightarrow$ Synthèse de l'OUA puis UA.
\end{enumerate}
\textbf{Formule à retenir :} « Contexte $\rightarrow$ Acteur/Thèse $\rightarrow$ Action/Institution $\rightarrow$ Héritage/Limite ».
\end{methode}

\begin{exoresolu}[Expliquer l'évolution du panafricanisme de la diaspora à l'Union Africaine]
\textbf{Énoncé :} Retracez les grandes étapes de l'histoire du panafricanisme en montrant comment le passage de la diaspora au continent a transformé ses objectifs et ses moyens d'action.

\tcblower
\textbf{Solution détaillée}

\textbf{Introduction}
Le panafricanisme n'est pas un bloc monolithique mais un \textbf{mouvement historique en devenir}, né de la rencontre traumatique de l'Afrique avec l'Europe (traite, colonisation). Son histoire suit le déplacement de son centre de gravité : de la \textbf{diaspora} (Amériques/Caraïbes) vers le \textbf{continent africain}.

\textbf{Phase 1 : La matrice diasporique (Fin XVIII\textsuperscript{e} – 1945) : La quête de dignité et de retour}
\begin{itemize}
    \item \textbf{Contexte :} Traite négrière, esclavagisme, racisme scientifique, hégémonie coloniale totale (Conférence de Berlin 1885).
    \item \textbf{Acteurs/Thèses :} \cle{Marcus Garvey} (UNIA, « L'Afrique aux Africains », retour physique, fierté noire - \cle{Black Pride}) ; \cle{W.E.B. Du Bois} (Congrès panafricains 1900, 1919, 1921, 1923, 1927 : approche intellectuelle, pétitionnaire, droits civiques) ; \cle{George Padmore} (lien avec le mouvement ouvrier international, \emph{How Britain Rules Africa}).
    \item \textbf{Évolution :} Du « Back to Africa » (Garvey) à la revendication d'autonomie politique pour les colonies (Du Bois/Padmore).
    \item \textbf{Limite :} Déconnecté des réalités politiques concrètes du continent ; élitisme intellectuel (Du Bois) ou échec organisationnel (Garvey).
\end{itemize}

\textbf{Phase 2 : Le tournant de Manchester 1945 : L'appropriation par les élites coloniales}
\begin{itemize}
    \item \textbf{Contexte :} Fin de la Seconde Guerre mondiale, affaiblissement des métropoles, montée des nationalismes (Inde 1947).
    \item \textbf{Acteur clé :} \cle{Kwame Nkrumah} (organisateur avec Padmore). Participation massive de futurs chefs d'État (Kenyatta, Nyerere, Kaunda).
    \item \textbf{Rupture majeure :} La résolution finale proclame le \textbf{droit à l'autodétermination par tous les moyens} (y compris la lutte armée). Le panafricanisme devient \textbf{programme d'action politique immédiate} pour l'indépendance.
    \item \textbf{Moyen :} Création de partis de masse (CPP au Ghana, TANU en Tanzanie), grève, désobéissance civile.
\end{itemize}

\textbf{Phase 3 : L'institutionnalisation étatique (1957-1963) : L'unité dans la diversité (Casablanca vs Monrovia)}
\begin{itemize}
    \item \textbf{Ghana 1957 :} Première indépendance au sud du Sahara. Nkrumah : « L'indépendance du Ghana n'a de sens que si elle est liée à la libération totale de l'Afrique ». Base arrière des mouvements de libération (ANC, PAIGC, FRELIMO).
    \item \textbf{Clivage 1961-63 :}
    \begin{center}
    \begin{tabularx}{\linewidth}{|l|X|X|}\hline
    \rowcolor{popblueL} \textbf{Groupe} & \textbf{Thèse (Moyens)} & \textbf{Objectif (Fin)} \\ \hline
    \textbf{Casablanca} (Ghana, Guinée, Mali, Maroc, Égypte, Algérie prov.) & Fédéralisme \textbf{immédiat} : Armée commune, marché commun, citoyenneté unique. Rupture néocoloniale. & \textbf{États-Unis d'Afrique} \textbf{maintenant}. \\ \hline
    \textbf{Monrovia} (Nigeria, Liberia, Côte d'Ivoire, Sénégal, Cameroun, etc.) & Coopération \textbf{progressive} : Respect souveraineté, non-ingérence, développement étape par étape. & \textbf{Unité} \textbf{éventuelle} par le fonctionnalisme. \\ \hline
    \end{tabularx}
    \end{center}
    \item \textbf{Synthèse de l'OUA (Addis-Abeba, 25 mai 1963) :} Charte de l'OUA. Compromis : \cle{Principes de Monrovia} (souveraineté, non-ingérence, frontières intangibles \emph{uti possidetis juris}) + \cle{Objectifs de Casablanca} (libération totale, coordination politique). \textbf{Paradoxe :} L'OUA libère le continent (Comité de libération) mais fige les frontières coloniales et tolère les dictatures (non-ingérence).
\end{itemize}

\textbf{Phase 4 : De l'OUA à l'Union Africaine (2002) et la ZLECAf (2018) : Le panafricanisme du développement}
\begin{itemize}
    \item \textbf{Constat :} Indépendance politique acquise (sauf Sahara Occidental), mais dépendance économique, conflits internes, marginalisation mondiale.
    \item \textbf{Rupture UA (Durban 2002) :} Droit d'ingérence dans les crimes de guerre/génocides (Art. 4h), Parlement panafricain, Cour de justice, \cle{NEPAD} (Nouveau Partenariat pour le Développement de l'Afrique) : appropriation du développement par les Africains (gouvernance, paix, investissement).
    \item \textbf{Actuel :} \textbf{ZLECAf} (entrée en vigueur 2021) = marché unique 1,3 milliard de consommateurs. \textbf{Agenda 2063} : « L'Afrique que nous voulons ». Passage du panafricanisme \emph{politique} (libération) au panafricanisme \emph{économique et social} (intégration par le bas).
\end{itemize}

\textbf{Conclusion}
L'histoire du panafricanisme est celle d'une \textbf{territorialisation} : né hors sol (diaspora), il s'enracine dans la lutte anticoloniale (Manchester), s'institutionnalise dans le compromis (OUA), et se réinvente aujourd'hui comme levier de \textbf{puissance économique collective} (UA/ZLECAf). Le Cameroun, pays charnière (Monrovia puis UA), incarne cette tension permanente entre souveraineté nationale et impératif d'intégration.

\textbf{Phrase de conclusion :} « Du rêve garveyiste de retour à la réalité de la ZLECAf, le panafricanisme a opéré une mutation historique décisive : il est passé d'une \emph{idéologie de la libération} portée par la diaspora à une \emph{stratégie de puissance collective} portée par les États africains eux-mêmes. »
\end{exoresolu}

\begin{methode}[Comparer des thèses philosophiques / idéologiques (Dialectique : Thèse - Antithèse - Synthèse)]
\textbf{Objectif :} Réussir le développement d'une dissertation (plan dialectique) ou une question de comparaison explicite au Bac.
\begin{enumerate}
    \item \textbf{Identifier le point de friction (le problème) :} Sur quoi les auteurs divergent-ils ? (Ex: Priorité : Unité politique immédiate vs Intégration économique progressive ? Rôle de l'État vs Marché ? Tradition vs Modernité ?).
    \item \textbf{Construire le tableau comparatif (Brouillon) :}
    \begin{center}
    \begin{tabular}{|l|c|c|}\hline
    \rowcolor{popblueL} \textbf{Critère de comparaison} & \textbf{Thèse A (ex: Nkrumah/Casablanca)} & \textbf{Thèse B (ex: Senghor/Monrovia / Nyerere/Ujamaa)} \\ \hline
    Ontologie (Vision de l'homme) & & \\ \hline
    Épistémologie (Source du savoir) & & \\ \hline
    Politique (Stratégie unité) & & \\ \hline
    Économique (Modèle dév.) & & \\ \hline
    \end{tabular}
    \end{center}
    \item \textbf{Rédiger en trois temps (Plan dialectique) :}
        \begin{itemize}
            \item \textbf{Thèse (Moment 1) :} Exposer la position la plus radicale/novatrice (souvent Nkrumah pour l'unité). Arguments FORCE. \emph{« Les partisans de... arguent que... »}
            \item \textbf{Antithèse (Moment 2) :} Exposer la position critique/réaliste (souvent Monrovia/Nyerere/Tshisekedi). Arguments FORCE. \emph{« En revanche, d'autres observent que... »} \textbf{Ne pas faire un catalogue : confronter point par point.}
            \item \textbf{Synthèse (Moment 3) :} Dépasser l'opposition. Montrer la complémentarité ou la condition de possibilité. \emph{« La tension entre... et... révèle que l'unité africaine nécessite à la fois la volonté politique (Casablanca) et la viabilité institutionnelle (Monrovia), synthétisées aujourd'hui dans l'UA et la ZLECAf. »}
        \end{itemize}
    \item \textbf{Utiliser les connecteurs logiques :} \emph{En revanche, toutefois, cependant, si... toutefois, dès lors, par conséquent, en définitive.}
\end{enumerate}
\textbf{Piège à éviter :} Le plan « A puis B » (catalogue). Il faut un plan « A \textbf{vs} B \textbf{sur le critère C} ».
\end{methode}

\begin{exoresolu}[Comparer les stratégies d'unité africaine : Nkrumah (Casablanca) vs Nyerere (Ujamaa/Monrovia modéré)]
\textbf{Énoncé :} « L'unité africaine ne se décrète pas, elle se construit. » Confrontrez la thèse de l'unité politique immédiate (Kwame Nkrumah) à la thèse de l'unité par la base économique et sociale (Julius Nyerere). Quelle synthèse pour l'Afrique d'aujourd'hui ?

\tcblower
\textbf{Solution détaillée}

\textbf{1. Problématisation}
L'opposition Nkrumah / Nyerere structure le débat panafricain : \textbf{Volontarisme étatique centralisé} (Casablanca) vs \textbf{Constructivisme sociologique décentralisé} (Ujamaa/Arusha). Le problème : comment passer de l'indépendance formelle à la souveraineté réelle ?

\textbf{2. Thèse : Nkrumah — L'unité politique comme condition \emph{sine qua non} (Le « Gouvernement de l'Union » maintenant)}
\begin{itemize}
    \item \textbf{Fondement :} Consciencisme + Marxisme-léninisme (analyse de classe, impérialisme). L'État colonial est l'instrument de l'exploitation. Seul un \textbf{État fédéral panafricain} (Haut Commandement, Marché commun, Monnaie unique, Citoyenneté) peut briser le néocolonialisme (balkanisation = faiblesse).
    \item \textbf{Argument force :} L'histoire (USA, URSS, Chine) prouve que l'unité politique précède l'intégration économique. « \emph{Seek ye first the political kingdom} ». L'OUA (Monrovia) a échoué à créer l'unité réelle (conflits, dépendance).
    \item \textbf{Application Cameroun :} Nkrumah aurait soutenu une fédération Cameroun-Nigeria-Tchad-RCA pour mutualiser les ressources (pétrole, eau, agriculture) et peser face aux multinationales.
\end{itemize}

\textbf{3. Antithèse : Nyerere — L'unité par le développement endogène et le consensus (Ujamaa / Arusha)}
\begin{itemize}
    \item \textbf{Fondement :} Socialisme africain (\cle{Ujamaa} = « vie de famille »), philosophie bantoue, christianisme social. L'unité imposée d'en haut (État fort) risque la dictature et l'effondrement économique (expérience ghanéenne post-Nkrumah, ou Zaïre/Mobutu).
    \item \textbf{Stratégie :} \textbf{Priorité au développement rural, à l'autosuffisance alimentaire, à l'éducation (Kiswahili), à la démocratie participative.} L'unité continentale ne sera durable que si les peuples en sont les artisans (intégration par le bas : routes, télécoms, marchés régionaux - CEEAC, SADC).
    \item \textbf{Critique de Nkrumah :} « L'unité de l'Afrique ne sera pas le cadeau d'un seul homme, fut-il un géant ». Risque de « super-État » bureaucratique déconnecté des villageois.
    \item \textbf{Application Cameroun :} Décentralisation, promotion des langues nationales, agriculture paysanne, intégration sous-régionale (CEMAC) comme brique de base.
\end{itemize}

\textbf{4. Synthèse : La dialectique de l'État et de la Société (L'UA et la ZLECAf comme dépassement)}
\begin{itemize}
    \item \textbf{Dépassement :} Nkrumah a raison sur la \textbf{nécessité politique} (sans souveraineté partagée, pas de politique commerciale commune, pas de défense commune). Nyerere a raison sur la \textbf{condition sociale} (sans adhésion des peuples, l'État panafricain est une coquille vide).
    \item \textbf{Synthèse historique :} L'\textbf{Union Africaine (2002)} et la \textbf{ZLECAf} réalisent cette synthèse :
        \begin{itemize}
            \item \textbf{Niveau macro (Nkrumah) :} Institutions supranationales (Parlement panafricain, Cour, Commissaire au Commerce), Agenda 2063, passeport africain, monnaie unique (projet Eco/Eco-6).
            \item \textbf{Niveau micro/méso (Nyerere) :} Intégration régionale effective (CEMAC, CEDEAO, SADC, COMESA) = « building blocks » ; société civile, secteur privé, femmes, jeunesse (CIMA, YALI) ; développement des corridors (autoroute Lagos-Mombasa, barrage Grand Inga).
        \end{itemize}
    \item \textbf{Nuance camerounaise :} Le Cameroun (Ahidjo/Biya) a historiquement choisi Monrovia (souveraineté), mais subit la contrainte Nkrumah (insécurité Boko Haram, crise anglophone, pression ZLECAf). La synthèse impose une \textbf{réforme de l'État} (gouvernance, décentralisation effective) pour que l'intégration profite aux populations.
\end{itemize}

\textbf{Conclusion}
La formule « L'unité ne se décrète pas, elle se construit » (proche de Nyerere) ne doit pas masquer l'urgence politique (Nkrumah). La synthèse opératoire aujourd'hui est : \textbf{Intégration politique différenciée (géométrie variable) + Intégration économique contraignante (ZLECAf) + Légitimité démocratique (Gouvernance)}.

\textbf{Phrase de conclusion :} « L'histoire a tranché : ni l'État fédéral instantané (utopie Nkrumah), ni le seul développement national autarcique (risque Nyerere) ne suffisent. L'Afrique de la ZLECAf invente une \emph{troisième voie} : une confédération d'États souverains réformés, liés par un marché unique et des institutions communes à géométrie variable. »
\end{exoresolu}

\begin{methode}[Appliquer les concepts à une situation concrète camerounaise (Savoir-faire : Application)]
\textbf{Objectif :} Mobiliser le consciencisme et le panafricanisme pour analyser une politique publique, un fait de société ou un défi de développement au Cameroun (Examen oral/écrit, Grand Oral).
\begin{enumerate}
    \item \textbf{Identifier le concept opératoire :} Quel concept éclaire la situation ? (Ex: \cle{Aliénation}, \cle{Néocolonialisme}, \cle{Communalisme}, \cle{Intégration régionale}, \cle{Souveraineté monétaire}).
    \item \textbf{Décrire le fait social / la politique publique (Données factuelles) :} Chiffres, textes de loi, institutions, acteurs locaux (MINADER, MINEPAT, Collectivités, ONG, Diaspora).
    \item \textbf{Établir le lien conceptuel (L'argument) :} \emph{« Cette situation illustre le concept de X dans la mesure où... »} ou \emph{« Cette politique contredit le principe de Y car... »}.
    \item \textbf{Évaluer (Esprit critique) :} Succès ? Échecs ? Contradictions ? Résistances ? Proposition d'ajustement (Praxis conscienciste).
    \item \textbf{Conclure par l'ouverture :} Lien avec l'Agenda 2063 ou la ZLECAf.
\end{enumerate}
\textbf{Exemples de « situations concrètes » types au Cameroun :}
\begin{itemize}
    \item Politique d'import-substitution (riz, poisson, blé) $\leftrightarrow$ Consciencisme (base matérialiste, autosuffisance).
    \item Décentralisation / Régionalisation (Code des collectivités) $\leftrightarrow$ Panafricanisme (Nyerere/Ujamaa, subsidiarité) / Consciencisme (humanisme, participation).
    \item Crise anglophone / Vivre-ensemble $\leftrightarrow$ Consciencisme (synthèse dialectique des héritages) / Panafricanisme (unité dans la diversité).
    \item ZLECAf / Corridor Douala-Ndjamena $\leftrightarrow$ Panafricanisme économique (Casablanca/Monrovia synthèse).
    \item FCFA / Monnaie unique CEMAC/ECO $\leftrightarrow$ Souveraineté monétaire (Nkrumah) vs Intégration monétaire (Monrovia).
\end{itemize}
\end{methode}

\begin{exoresolu}[Analyser la politique d'import-substitution au Cameroun à la lumière du consciencisme]
\textbf{Énoncé :} Le Cameroun lance la « Stratégie Nationale de Développement 2030 » (SND30) avec pour axe majeur l'import-substitution (riz, maïs, poisson, blé, médicaments). En vous appuyant sur les fondements idéologiques et philosophiques du consciencisme, analysez la pertinence et les limites de cette stratégie.

\tcblower
\textbf{Solution détaillée}

\textbf{1. Identification des concepts opératoires (Consciencisme)}
\begin{itemize}
    \item \textbf{Base matérialiste :} La libération économique prime. « Qui vous nourrit vous contrôle ». Maîtriser les moyens de production (terre, semences, transformation).
    \item \textbf{Synthèse dialectique (Tradition/Islam/Occident) :} Moderniser l'agriculture paysanne (tradition) par la recherche scientifique (Occident) et la solidarité/coopératives (Islam/Communalisme).
    \item \textbf{Humanisme / Fin :} L'épanouissement de l'homme africain (sécurité alimentaire, emploi jeunes, revenu paysan).
    \item \textbf{Praxis :} Théorie + Action. Pas de conscience sans transformation matérielle.
\end{itemize}

\textbf{2. Description factuelle de la SND30 / Import-substitution (Données 2023-2024)}
\begin{itemize}
    \item \textbf{Contexte :} Déficit commercial structurel (~ 2000 Mrds FCFA/an), dépendance riz (350 000 t importées), poisson (200 000 t), blé (1 M t). Crise Ukraine/Covid = choc prix.
    \item \textbf{Mesures :} Subventions intrants (engrais, semences certifiées), mécanisation (SODERAM), transformation locale (usines riz Ndé, Yagoua ; farine de manioc/plantain), financement agricole (Banque Agricole annoncée), ZLECAf (export futur).
    \item \textbf{Résultats mitigés :} Production riz en hausse (300 000 t 2023 vs 150 000 t 2015) mais rendements faibles (2,5 t/ha vs 6 t/ha potentiel), financement encore faible (< 5\% crédit bancaire), foncier bloqué, changement climatique.
\end{itemize}

\textbf{3. Analyse conscienciste : Pertinence (Thèse)}

\textbf{A. Rupture avec l'aliénation néocoloniale (Matérialisme).}
L'import-substitution vise la \textbf{souveraineté alimentaire}, condition de la souveraineté politique (Nkrumah : « Indépendance économique »). Réduire la facture alimentaire libère des devises pour l'industrialisation (base matérialiste). \textbf{Exemple :} L'usine de transformation de manioc en farine (Moungo) substitue l'import de blé (occidental) par une ressource locale (tradition) + technologie (occident) = \textbf{Synthèse conscienciste}.

\textbf{B. Réhabilitation du communalisme paysan (Humanisme/Tradition).}
La stratégie cible les \textbf{petits producteurs} (70\% population active). Organisation en coopératives (GIC, Coop-CA) = forme moderne du \cle{communalisme} (entraide, mise en commun moyens). \textbf{Exemple :} Projet PIDMA (Productivité Agricole) finance des groupements. Cela rejoint l'idéal nkrumahiste : l'homme (paysan) est la fin, pas l'actionnaire étranger.

\textbf{C. Valorisation du savoir endogène (Synthèse culturelle).}
Promotion du \textbf{riz Nerica} (croisement Oryza glaberrima africain $\times$ O. sativa asiatique) + farines composites (manioc, maïs, soja, patate). C'est l'application littérale de la \textbf{synthèse tripartite} : génome africain (tradition) + biotechnologie (occident) + diffusion par vulgarisation communautaire (islam/tradition).

\textbf{4. Analyse conscienciste : Limites et Contradictions (Antithèse)}

\textbf{A. Risque de « capitalisme d'État » bureaucratique (Trahison du matérialisme).}
Subventions massives captées par élites urbaines / « agriculteurs du dimanche » (fonctionnaires, militaires) $\rightarrow$ \textbf{Néocolonialisme interne}. Nkrumah dénonçait la bourgeoisie compradore. Si l'État achète les engrais à prix fort aux multinationales (Yara, OCP) pour les donner aux riches, la base matérialiste est détournée.

\textbf{B. Négligence de la transformation locale / Industrie (Praxis incomplète).}
Produire du paddy (riz non décortiqué) sans usine de décorticage moderne = vendre de la matière première. Le consciencisme exige la \textbf{maîtrise de la chaîne de valeur} (semence $\rightarrow$ production $\rightarrow$ transformation $\rightarrow$ commercialisation $\rightarrow$ consommation). Usines de Ndé/Yagoua souvent en panne ou sous-capacité.

\textbf{C. Foncier et Genre : Contradiction avec l'Humanisme.}
Accès à la terre difficile pour femmes/jeunes (droit coutumier patriarcal). Le consciencisme (humanisme) exige l'égalité. Sans réforme foncière inclusive, l'import-substitution renforce les inégalités traditionnelles (aliénation interne).

\textbf{D. Dépendance technologique persistante.}
Semences hybrides (non reproductibles), engrais chimiques, machines importées. Pas encore de \textbf{souveraineté technologique} (instituts de recherche IRAD sous-financés, brevets étrangers). Nkrumah voulait une « Académie des Sciences africaine ».

\textbf{5. Synthèse et Proposition (Praxis conscienciste actualisée)}
La SND30 est une \textbf{application partielle et nécessaire} du consciencisme. Pour être pleinement conscienciste, elle doit opérer un \textbf{double mouvement} :
\begin{enumerate}
    \item \textbf{Basculer de la subvention à l'investissement structurel :} Irrigation (barrages Logone, Lom Pangar), routes rurales, énergie (solaire), recherche publique (semences paysannes reproductibles).
    \item \textbf{Institutionnaliser le pouvoir paysan :} Droits fonciers sécurisés (femmes/jeunes), gouvernance coopérative démocratique, prix planchers garantis (État régulateur, pas seulement pourvoyeur).
    \item \textbf{Intégrer la ZLECAf :} Produire pour le marché camerounais \textbf{ET} pour le marché africain (Tchad, RCA, Gabon, Nigeria) $\rightarrow$ Échelle critique = industrialisation.
\end{enumerate}

\textbf{Conclusion}
L'import-substitution est l'\textbf{traduction concrète} de l'impératif conscienciste de « décolonisation économique ». Mais sans \textbf{démocratie économique} (pouvoir aux producteurs) et \textbf{souveraineté technologique}, elle risque de reproduire les structures d'aliénation qu'elle combat. Le consciencisme exige que la SND30 ne soit pas un « projet » de plus, mais le levier d'une \textbf{refondation de l'État camerounais} au service de l'homme.

\textbf{Phrase de conclusion :} « La SND30 incarne la \emph{volonté} conscienciste (souveraineté), mais son efficacité suppose la \emph{méthode} conscienciste : dialectique de l'État stratège et de la paysannerie organisée, synthèse de la tradition agraire et de la science moderne, le tout orienté vers l'humanisme intégral. »
\end{exoresolu}

\begin{methode}[Rédiger une introduction de dissertation philosophique (Standard Bac Probatoire)]
\textbf{Objectif :} Produire une introduction « modèle » en 4 étapes strictes (Accroche $\rightarrow$ Définitions/Problématisation $\rightarrow$ Problématique $\rightarrow$ Annonce du plan).
\begin{enumerate}
    \item \textbf{L'Accroche (1 phrase max) :} Citation d'auteur au programme (Nkrumah, Nyerere, Fanon, Senghor, Wiredu, Mudimbe), fait historique précis (Manchester 1945, OUA 1963, UA 2002, ZLECAf 2021), ou paradoxe saisissant. \textbf{Interdit :} « Depuis la nuit des temps... », « De tout temps les hommes... ».
    \item \textbf{La Définition des termes du sujet (2-3 phrases) :} Définir \textbf{tous} les termes clés du sujet en les reliant au programme (Consciencisme, Panafricanisme, Développement, Unité, Aliénation, Souveraineté, Communalisme). Citer l'auteur de référence.
    \item \textbf{La Problématisation (Le cœur) :} Montrer la \textbf{tension}, la \textbf{contradiction} ou le \textbf{problème} caché derrière les définitions.
        \begin{itemize}
            \item \emph{Structure type :} « Si [Terme 1] signifie [Définition A], et [Terme 2] signifie [Définition B], alors comment [Articulation difficile] ? »
            \item \emph{Exemple :} « Si le panafricanisme vise l'unité politique (Nkrumah), mais que les États protègent farouchement leur souveraineté (Monrovia), \textbf{comment l'unité africaine peut-elle être effective sans dissoudre les États nations ?} »
        \end{itemize}
    \item \textbf{L'Annonce du Plan (1 phrase, 3 parties) :} Utiliser les connecteurs : \emph{Dans un premier temps... Dans un second temps... Enfin...}
        \begin{itemize}
            \item Partie I : Thèse (Analyse du concept / Argument principal).
            \item Partie II : Antithèse (Limites / Critique / Réalité contraignante).
            \item Partie III : Synthèse (Dépassement / Perspective actuelle / Solution concrète).
        \end{itemize}
\end{enumerate}
\textbf{Contrainte de forme :} Un seul paragraphe (ou deux max). Pas de « Je vais étudier... », mais « Nous analyserons... puis nous montrerons... enfin nous verrons... ».
\end{methode}

\begin{exoresolu}[Rédiger l'introduction : « Le panafricanisme est-il encore pertinent pour le développement de l'Afrique ? »]
\textbf{Énoncé :} Sujet de dissertation type Bac : « Le panafricanisme est-il encore pertinent pour le développement de l'Afrique ? » Rédigez l'introduction complète.

\tcblower
\textbf{Solution détaillée}

\textbf{Introduction rédigée :}

« \textbf{L'unité africaine n'est pas un rêve romantique, c'est une nécessité vitale} », affirmait \cle{Kwame Nkrumah} à Addis-Abeba en 1963, fondant l'OUA sur l'impératif d'une libération totale du continent. Soixante ans plus tard, alors que l'\cle{Union Africaine} lance la \cle{ZLECAf} (Zone de Libre-Échange Continentale Africaine) et l'\cle{Agenda 2063}, le \cle{panafricanisme} — défini historiquement comme l'idéologie prônant la solidarité et l'unité politique, économique et culturelle des peuples africains et de leur diaspora (\cle{Du Bois}, \cle{Padmore}, \cle{Nkrumah}) — et le \cle{développement} — entendu ici non comme simple croissance du PIB mais comme \emph{développement humain intégral} (\cle{Sen}, \cle{Nyerere}) : autonomie, dignité, maîtrise endogène — semblent plus que jamais liés. \textbf{Pourtant, une tension majeure persiste :} si le panafricanisme a remporté la bataille de la décolonisation politique, il peine à s'imposer comme le moteur effective de la transformation structurelle des économies africaines, fragmentées par des frontières héritées du colonialisme, des souverainetés nationales jalousement gardées et une dépendance extravertie structurelle. \textbf{Dès lors, le panafricanisme constitue-t-il encore un opératoire théorique et pratique pertinent pour penser et réaliser le développement africain, ou est-il devenu un mythe mobilisateur déconnecté des réalités géopolitiques et économiques contemporaines ?} \textbf{Nous analyserons d'abord en quoi le panafricanisme reste le cadre théorique indispensable pour penser l'autonomie africaine (I), avant d'examiner les obstacles structurels et politiques qui en limitent l'effectivité concrète (II), pour enfin montrer comment la refondation institutionnelle actuelle (UA, ZLECAf, Agenda 2063) opère une synthèse pragmatique entre idéal unitaire et réalisme étatique (III).}
\end{exoresolu}

\begin{methode}[Structurer le développement dialectique (Corps de la dissertation)]
\textbf{Objectif :} Rédiger des parties équilibrées, argumentées, illustrées, sans répétition ni hors-sujet.
\begin{enumerate}
    \item \textbf{Une idée directrice par paragraphe (1 paragraphe = 1 argument).} Première phrase = \textbf{Phrase thèse} (l'argument). Dernière phrase = \textbf{Transition} vers l'argument suivant.
    \item \textbf{Structure interne du paragraphe (A.E.I.) :}
        \begin{itemize}
            \item \textbf{A - Affirmation :} L'argument clair (théorique).
            \item \textbf{E - Explication/Preuve :} Référence auteur (\cle{Nkrumah}, \cle{Nyerere}, \cle{Fanon}, \cle{Mudimbe}), concept précis (\cle{néocolonialisme}, \cle{communalisme}, \cle{intégration régionale}), mécanisme causal.
            \item \textbf{I - Illustration concrète (Cameroun / Afrique) :} Chiffre, politique publique, fait historique, institution (UA, CEMAC, ZLECAf, SND30, PIDMA, Décentralisation). \textbf{Obligatoire au Bac Tech.}
        \end{itemize}
    \item \textbf{Équilibrer les parties :} Même nombre de paragraphes (2 ou 3) par grande partie (I, II, III). Même niveau d'exigence.
    \item \textbf{Transitions inter-parties :} Expliciter le mouvement dialectique.
        \begin{itemize}
            \item I $\rightarrow$ II : « Cependant, cette nécessité théorique se heurte à la réalité du terrain... »
            \item II $\rightarrow$ III : « C'est pour dépasser ce blocage que l'Afrique contemporaine invente de nouvelles formes d'unité... »
        \end{itemize}
    \item \textbf{Vocabulaire de liaison (Cohérence) :} \emph{En effet, plus précisément, c'est ainsi que, par conséquent, toutefois, en revanche, de surcroît, corrélativement, in fine.}
\end{enumerate}
\textbf{Règle d'or :} \textbf{Zéro affirmation sans justification.} « L'Afrique est riche » $\rightarrow$ Faux (généralisation). « L'Afrique dispose de 30\% des réserves minières mondiales (Source UA 2023) mais ne transforme que 5\% localement » $\rightarrow$ Vrai (argumenté).
\end{methode}

\begin{exoresolu}[Rédiger la Partie I d'une dissertation : « Le panafricanisme, cadre théorique indispensable de l'autonomie »]
\textbf{Énoncé :} Rédigez la première partie (2 paragraphes) de la dissertation dont l'introduction a été faite ci-dessus. Sujet : « Le panafricanisme est-il encore pertinent pour le développement de l'Afrique ? »

\tcblower
\textbf{Solution détaillée}

\textbf{PARTIE I : LE PANAFRICANISME, CADRE THÉORIQUE INDISPENSABLE POUR PENSER L'AUTONOMIE AFRICAINE}

\textbf{Paragraphe 1 : La rupture épistémologique : décoloniser le développement (Fanon / Nkrumah / Mudimbe).}
\textbf{[Affirmation]} Le panafricanisme est d'abord pertinent car il opère une \textbf{rupture épistémologique} radicale avec le paradigme du « développement » imposé par l'Occident (modernisation, ajustement structurel, mimétisme). \textbf{[Explication]} Pour \cle{Frantz Fanon} (\emph{Les Damnés de la Terre}), le développement ne peut être que \textbf{auto-centré} : « Chaque génération doit, dans une relative opacité, découvrir sa mission, la remplir ou la trahir ». \cle{Kwame Nkrumah} (\emph{Néocolonialisme, dernier stade de l'impérialisme}) démontre que l'indépendance politique sans unité économique laisse le continent otage du « néocolonialisme » : contrôle des prix matières premières, dette, termes de l'échange défavorables. \cle{Valentin Mudimbe} (\emph{L'Invention de l'Afrique}) ajoute que le concept même de « développement » est une construction coloniale ; le panafricanisme permet de \textbf{réinventer les catégories} (richesse, pauvreté, progrès) à partir des réalités africaines (communalisme, rapport à la nature, oralité). \textbf{[Illustration]} Au Cameroun, les \textbf{Programmes d'Ajustement Structurel (PAS)} des années 90 (FMI/BM) ont imposé la libéralisation brutale du café/cacao, ruinant les petits producteurs (coopératives dissoutes) au nom de l'efficacité marchande. Le retour actuel à la \textbf{régulation étatique} (Office National du Café et du Cacao - ONCC renaissant, prix plancher) et à la \textbf{transformation locale} (usines de cacao à Kribi, Douala) valide la thèse panafricaniste : pas de développement sans \textbf{souveraineté de la décision économique}.

\textbf{Paragraphe 2 : L'impératif d'échelle : le marché unique comme condition de l'industrialisation (Nkrumah / UA / ZLECAf).}
\textbf{[Affirmation]} La pertinence du panafricanisme tient ensuite à une \textbf{nécessité structurelle économique} : la balkanisation du continent en 54 États (dont 16 enclavés) rend impossible l'industrialisation par substitution aux importations à l'échelle nationale. \textbf{[Explication]} \cle{Nkrumah} l'avait formulé dès 1963 : « Aucun État africain n'est viable économiquement isolément ». La \cle{ZLECAf} (entrée en vigueur 2021), en créant un marché unique de \textbf{1,3 milliard de consommateurs} et un PIB cumulé de \textbf{3 400 milliards USD}, réalise l'objectif casablancais d'un \textbf{marché commun panafricain}. Elle permet les \textbf{économies d'échelle} indispensables pour attirer l'investissement productif, financer les infrastructures (PIDA - Programme de Développement des Infrastructures en Afrique) et négocier en bloc face aux puissances mondiales (UE, Chine, USA). \textbf{[Illustration]} Le \textbf{Cameroun}, « Afrique en miniature » et carrefour CEMAC/CEDEAO, en est l'illustration parfaite : son industrie naissante (aluminium Alucam, ciment Cimencam, textile CICAM, agro-industrie CDC/PAMOL/SOSUCAM) ne devient compétitive que si elle accède au marché nigérian (200M hab.), tchadien, centrafricain, gabonais sans droits de douane. Le corridor \textbf{Douala-Ndjamena} (route, rail projeté, pipeline, fibre optique) est une \textbf{application concrète du panafricanisme fonctionnaliste (Monrovia) au service de l'intégration profonde (Casablanca)}. Sans cadre panafricain, le Cameroun reste une économie d'enclave.

\textbf{[Transition vers Partie II]} \emph{Toutefois, cette pertinence théorique et structurelle se heurte à la persistance de freins politiques, institutionnels et sécuritaires majeurs qui vident souvent le panafricanisme de sa substance opératoire.}
\end{exoresolu}

\begin{methode}[Conclure une dissertation : Synthèse, Réponse, Ouverture]
\textbf{Objectif :} Clore le devoir en répondant explicitement à la problématique, en résumant la dialectique et en ouvrant sur une perspective précise (pas une banalité).
\begin{enumerate}
    \item \textbf{Reprendre la problématique (1 phrase) :} Rappeler la tension initiale.
    \item \textbf{Synthétiser le parcours dialectique (2-3 phrases) :} \emph{« Nous avons vu que... (I), mais que... (II), d'où l'émergence de... (III). »} Ne pas répéter les arguments, en donner la \textbf{portée philosophique}.
    \item \textbf{Apporter une réponse nuancée (Thèse finale) :} \emph{« En définitive, le panafricanisme n'est pas une solution magique, mais la \textbf{condition nécessaire} (pas suffisante) du développement. Sa pertinence dépend de la \textbf{volonté politique} de transformer la souveraineté formelle en puissance collective. »}
    \item \textbf{L'Ouverture (1 phrase précise) :} Lier à une autre notion du programme, un auteur non traité, ou une actualité brûlante.
        \begin{itemize}
            \item \emph{Notion :} « Cette refondation panafricaine suppose une refonte de la \textbf{démocratie} (gouvernance, reddition des comptes) et de l'\textbf{éducation} (philosophie africaine, langues nationales) — notions du programme. »
            \item \emph{Auteur :} « Comme le souligne \cle{Achille Mbembe} (\emph{Critique de la raison nègre}), l'avenir de l'Afrique se joue dans l'invention de nouvelles formes de \textbf{communalité planétaire}, au-delà du seul panafricanisme étatique. »
            \item \emph{Actualité :} « La réussite de la \textbf{ZLECAf} et de l'\textbf{Agenda 2063} sera le test ultime de la capacité du panafricanisme à passer de l'idéologie à l'ingénierie du développement. »
        \end{itemize}
\end{enumerate}
\textbf{Interdits :} « Pour conclure, je dirais que... », « J'espère que... », « L'avenir nous le dira », ouverture vague (« Il faudrait aussi étudier la place des femmes... »).
\end{methode}

\begin{exoresolu}[Rédiger la conclusion de la dissertation sur le panafricanisme et le développement]
\textbf{Énoncé :} Rédigez la conclusion de la dissertation dont l'introduction et la Partie I ont été traitées. Rappel du plan : I. Cadre théorique indispensable. II. Obstacles structurels (souveraineté, conflits, néocolonialisme). III. Synthèse : UA/ZLECAf/Agenda 2063 (Pragmatisme institutionnel).

\tcblower
\textbf{Solution détaillée}

\textbf{Conclusion rédigée :}

\textbf{[Reprise problématique]} Au terme de cette analyse, la question de la pertinence du panafricanisme pour le développement africain trouve une réponse dialectique : \textbf{le panafricanisme n'est pas une option parmi d'autres, il est la condition \emph{sine qua non} de tout développement authentique, mais sa pertinence effective dépend de sa capacité à se réinventer institutionnellement.}

\textbf{[Synthèse dialectique]} Nous avons montré que, \textbf{théoriquement (I)}, le panafricanisme reste l'unique cadre capable de penser l'autonomie africaine en brisant l'aliénation épistémologique (Fanon, Mudimbe) et en résolvant l'impossibilité structurelle de l'industrialisation nationale (Nkrumah, ZLECAf). \textbf{Pourtant (II)}, cette nécessité vitale se heurte au « réalisme » des États (souveraineté Monrovia), aux conflits endogènes (terrorisme, coups d'État), à la capture des ressources par les élites compradores et à la persistance des logiques néocoloniales (FCFA, dette, accords de défense). \textbf{C'est précisément pour dépasser cette contradiction (III)} que l'Union Africaine (2002) et ses instruments (ZLECAf, Agenda 2063, PAM, ACDC) opèrent une \textbf{synthèse pragmatique} : l'unité ne se décrète plus par le haut (utopie Casablanca), elle se construit par le bas (intégration régionale, marché unique, libre circulation, gouvernance partagée), tout en dotant le continent d'institutions supranationales à compétences réelles (Cour, Parlement, CDC Africa).

\textbf{[Réponse nuancée]} \textbf{En définitive, le panafricanisme est pertinent \emph{à la condition} de cesser d'être une \emph{idéologie d'État} pour devenir une \emph{ingénierie politique et économique} partagée : il n'y a pas de développement durable sans intégration, mais pas d'intégration durable sans \textbf{démocratie interne} et \textbf{reddition des comptes} des États aux peuples.}

\textbf{[Ouverture précise]} Cette exigence renvoie directement à la notion de \textbf{consciencisme} (Nkrumah) : le développement africain ne saurait se réduire à des indicateurs macroéconomiques ; il suppose une \textbf{révolution des consciences} (éducation panafricaine, valorisation des langues/savoirs endogènes, éthique de la responsabilité) pour que l'intégration économique (ZLECAf) ne reproduise pas, à l'échelle continentale, les inégalités qu'elle prétend combattre à l'échelle nationale.
\end{exoresolu}

\begin{attention}[Les 5 erreurs qui coûtent des points au Bac (Probatoire / Bac Tech)]
\begin{enumerate}
    \item \textbf{Confondre « Panafricanisme » et « Négritude » (Senghor/Césaire).}
    \emph{Erreur :} « Le panafricanisme de Senghor prône le métissage culturel... »
    \emph{Correction :} La \cle{Négritude} (Senghor, Césaire, Damas) est un mouvement \textbf{littéraire et culturel} (affirmation des valeurs noires). Le \cle{Panafricanisme} (Du Bois, Garvey, Nkrumah, Nyerere) est un mouvement \textbf{politique et idéologique} (unité, libération, développement). Ils se croisent (Senghor à l'OUA), mais ne se confondent pas. \textbf{Distinction programme :} Chapitre « Consciencisme, Panafricanisme » $\neq$ Chapitre « Négritude » (souvent en Terminale ou option).
    \item \textbf{Réduire le Consciencisme à un « syncrétisme religieux » ou un « mélange ».}
    \emph{Erreur :} « Nkrumah mélange tradition, islam et christianisme pour faire plaisir à tout le monde. »
    \emph{Correction :} Le consciencisme est une \textbf{synthèse dialectique} (Hegel/Marx) sur base \textbf{matérialiste}. Il \emph{nie} (critique l'aliénation), \emph{conserve} (humanisme communal) et \emph{élève} (science, raison). C'est une \textbf{philosophie de l'action} (praxis), pas un syncrétisme cultuel. \textbf{Mot-clé :} \cle{Matérialisme dialectique}, \cle{Humanisme}, \cle{Praxis}.
    \item \textbf{Traiter l'histoire du panafricanisme comme une « liste de congrès » (Catalogue).}
    \emph{Erreur :} « En 1900 congrès Londres, 1919 Paris, 1921 Londres/Bruxelles/Paris, 1923 Londres/Lisbonne, 1927 New York, 1945 Manchester, 1958 Accra, 1963 Addis-Abeba... »
    \emph{Correction :} \textbf{Périodiser par phases logiques} (Diaspora $\rightarrow$ Manchester $\rightarrow$ Indépendances/Casablanca-Monrovia $\rightarrow$ OUA $\rightarrow$ UA/ZLECAf). Pour chaque phase : \textbf{Contexte $\rightarrow$ Acteur $\rightarrow$ Thèse $\rightarrow$ Rupture/Continuité}. \textbf{Analyser}, ne pas lister.
    \item \textbf{Oublier le « Cameroun » dans l'application concrète (Spécificité Bac Tech).}
    \emph{Erreur :} Dissertation générale sur l'Afrique sans aucun exemple camerounais (SND30, ZLECAf/Douala-Ndjamena, Décentralisation, Crise NOSO, Bilinguisme, FCFA/CEMAC, PIDMA, ONCC).
    \emph{Correction :} Le programme exige : « \textbf{Appliquer les notions de l'unité à des situations concrètes de la vie courante} ». \textbf{Chaque grande thèse doit être illustrée par un fait camerounais daté/chiffré.} C'est la marque de l'enseignement \textbf{technique} (ancrage territorial).
    \item \textbf{Rédiger une « fausse synthèse » en Partie III (Moyen terme / Compromis mou).}
    \emph{Erreur :} « Il faut un peu de Nkrumah et un peu de Nyerere, les deux ont raison. »
    \emph{Correction :} La synthèse doit \textbf{dépasser} la contradiction en en expliquant la \textbf{résolution historique actuelle}. Ex: « La ZLECAf réalise la synthèse : elle instaure le \textbf{marché unique} (objectif Casablanca) par la \textbf{libéralisation progressive et la coopération régionale} (méthode Monrovia), sous l'égide d'institutions supranationales renforcées (UA) ». \textbf{La synthèse = Nouvel objet d'analyse (UA/ZLECAf/Agenda 2063), pas juste un mélange.}
\end{enumerate}
\end{attention}

\newpage

\coursec{Exercices}

\courssub{Je vérifie mes connaissances}

\begin{exercice}[QCM : Repères chronologiques et figures]
Parmi les propositions suivantes, une seule est exacte. La recopier en la justifiant brièvement.
\begin{enumerate}
\item Le 5\textsuperscript{e} Congrès panafricain s'est tenu à Manchester en 1945 sous la présidence de Kwame Nkrumah.
\item L'Organisation de l'Unité Africaine (OUA) a été fondée à Addis-Abeba en 1963, sur une initiative conjointe de Nkrumah et de Sékou Touré.
\item Le concept de \emph{consciencisme} a été théorisé par Julius Nyerere dans son ouvrage \emph{Le Socialisme africain} (1967).
\item Ruben Um Nyobè a prononcé un discours historique devant la IV\textsuperscript{e} Commission de l'ONU en 1952 pour réclamer l'indépendance immédiate du Cameroun.
\end{enumerate}
\end{exercice}

\begin{exercice}[Vrai / Faux justifié : Notions clés]
Dire si les affirmations sont vraies ou fausses et justifier la réponse en une phrase précise.
\begin{enumerate}
\item Le communalisme, tel que défini par Nkrumah, correspond à la propriété privée des moyens de production dans la société traditionnelle africaine.
\item Le néocolonialisme désigne, pour Nkrumah, la poursuite de la domination économique et politique par d'autres moyens après l'accession formelle à l'indépendance.
\item La \emph{conscience révolutionnaire} est, dans le consciencisme, l'adhésion passive aux traditions ancestrales sans esprit critique.
\item Le panafricanisme politique né au congrès de Manchester (1945) revendique l'unité fédérale des États africains comme condition de la libération totale.
\end{enumerate}
\end{exercice}

\begin{exercice}[Définitions attendues au Bac]
Définir rigoureusement les notions suivantes en deux lignes maximum chacune, en précisant leur auteur ou leur courant de référence.
\begin{enumerate}
\item Consciencisme
\item Panafricanisme
\item Néocolonialisme (selon Nkrumah)
\item Communalisme (selon Nkrumah)
\end{enumerate}
\end{exercice}

\begin{exercice}[Correspondances : Dates, Congrès, Figures]
Établir la correspondance exacte entre les éléments de la colonne A et ceux de la colonne B.
\begin{center}
\begin{tabularx}{\linewidth}{|c|X|}\hline
\textbf{Colonne A} & \textbf{Colonne B} \\ \hline
1. Congrès de Manchester & A. 1900 (Londres) \\ \hline
2. Fondation de l'OUA & B. 1945 \\ \hline
3. Premier Congrès panafricain (Du Bois) & C. 1957 (Indépendance du Ghana) \\ \hline
4. Indépendance du Ghana & D. 1963 (Addis-Abeba) \\ \hline
5. Création de l'Union Africaine (UA) & E. 2002 (Durban) \\ \hline
\end{tabularx}
\end{center}
\end{exercice}

\courssub{Je m'entraîne}

\begin{exercice}[Analyse de thèse : Extrait de \emph{Consciencism} (Nkrumah)]
On lit dans \emph{Consciencism} : « La philosophie du consciencisme est une philosophie de la révolution. Elle ne se contente pas d'interpréter le monde, elle vise à le transformer. »
\begin{enumerate}
\item Dégager la thèse principale de cet extrait.
\item Expliquer en quoi le consciencisme se distingue d'une philosophie purement contemplative ou académique.
\item Identifier les trois sources idéologiques que Nkrumah synthétise dans le consciencisme et dire comment elles s'articulent pour servir le développement.
\end{enumerate}
\end{exercice}

\begin{exercice}[Le tournant de Manchester (1945) : Rupture ou continuité ?]
À partir de vos connaissances sur les congrès panafricains (1900, 1919, 1921, 1923, 1927, 1945), répondre aux questions :
\begin{enumerate}
\item Quelles sont les deux ruptures majeures introduites par le 5\textsuperscript{e} Congrès de Manchester par rapport aux congrès précédents (Du Bois) ?
\item Citer trois figures africaines présentes à Manchester qui deviendront chefs d'État et préciser leur pays.
\item Expliquer pourquoi la résolution finale de Manchester marque la naissance du panafricanisme \emph{de masse} et \emph{révolutionnaire}.
\end{enumerate}
\end{exercice}

\begin{exercice}[Application camerounaise : Um Nyobè et l'unité]
Ruben Um Nyobè, figure majeure du nationalisme camerounais, a défendu l'idée d'une indépendance réelle et d'une unité africaine par la base.
\begin{enumerate}
\item Rappeler le contenu essentiel de son intervention à l'ONU (1952, 1953, 1954).
\item Montrer en quoi sa conception de l'unité (fédéralisme, solidarité des peuples) rejoint ou diffère du panafricanisme d'État (OUA, 1963).
\item Donner un exemple concret actuel (infrastructure, institution, politique) qui illustre la réalisation ou l'échec de cette unité au Cameroun.
\end{enumerate}
\end{exercice}

\begin{exercice}[Débat structuré : Le Franc CFA]
Sujet : « Le Franc CFA : instrument de solidarité africaine ou frein à la souveraineté monétaire du Cameroun ? »
\begin{enumerate}
\item Présenter les deux thèses opposées en deux arguments chacune (thèse A : solidarité/stabilité ; thèse B : frein/souveraineté).
\item Mobiliser la notion de néocolonialisme (Nkrumah) pour analyser la thèse B.
\item Rédiger une conclusion personnelle (5 lignes) tranchant le débat en s'appuyant sur le critère du \emph{développement endogène}.
\end{enumerate}
\end{exercice}

\courssub{Je me prépare au Bac}

\begin{exercice}[Sujet de dissertation type Baccalauréat Technique]
\textbf{Sujet :} « L'unité africaine est-elle une condition nécessaire du développement de l'Afrique ? »
\begin{enumerate}
\item Analyser le sujet : définir les termes clés, dégager la problématique, annoncer le plan en deux ou trois parties.
\item Rédiger l'introduction complète (amorce, définitions, problématique, annonce du plan).
\item Traiter la première partie du développement (argumentaire principal avec exemples précis : OUA/UA, ZLECAf, CEMAC, infrastructures transfrontalières).
\item Rédiger une conclusion synthétique ouvrant sur une perspective.
\end{enumerate}
\end{exercice}

\begin{exercice}[Commentaire de texte : \emph{Consciencism}, Kwame Nkrumah]
\textbf{Texte :} « La société traditionnelle africaine était égalitaire. Mais l'égalitarisme traditionnel ne suffit pas. Il faut une conscience socialiste, c'est-à-dire une conscience révolutionnaire, pour transformer la structure de la société et la mettre au service du développement rapide. Le consciencisme est la philosophie qui permet cette transformation en synthétisant l'héritage traditionnel, l'apport islamo-arabe et l'héritage euro-chrétien. »
\begin{enumerate}
\item \textbf{Dégager la thèse} de l'auteur en une phrase.
\item \textbf{Analyser les arguments} : identifier les trois piliers de la synthèse conscienciste et expliquer la fonction de chacun dans le projet de développement.
\item \textbf{Discuter la pertinence actuelle} : cette synthèse est-elle encore opératoire face aux défis de la mondialisation, de la dette et de la gouvernance au Cameroun ? (Argumenter en 15 lignes minimum).
\end{enumerate}
\end{exercice}

\begin{exercice}[Question à réponse construite : Fondements et Développement]
\textbf{Sujet :} « Présentez les fondements idéologiques et philosophiques du consciencisme et montrez comment ils visent le développement de l'Afrique. »
\begin{enumerate}
\item \textbf{Idéologie :} Expliciter le communalisme (base traditionnelle), le socialisme (objectif moderne) et l'anti-impérialisme/anti-néocolonialisme (lutte politique).
\item \textbf{Philosophie :} Expliquer le matérialisme dialectique adapté (matière-esprit, conscience révolutionnaire, catharsis) et la synthèse des trois courants (traditionnel, islamo-arabe, euro-chrétien).
\item \textbf{Articulation au développement :} Démontrer, exemples camerounais à l'appui (politique agricole, industrialisation, éducation, ZLECAf), comment ce cadre théorique entend passer de l'indépendance formelle à la souveraineté économique réelle.
\item \textbf{Conclusion :} Évaluer les limites de cette application dans le contexte camerounais actuel.
\end{enumerate}
\end{exercice}

\newpage

\coursec{Corrigés des exercices}

\begin{corrige}[Exercice 1 — QCM : Repères chronologiques et figures]
\textbf{Proposition exacte : la proposition 4.}

\begin{enumerate}
\item \textbf{Fausse.} Le 5\textsuperscript{e} Congrès panafricain s'est bien tenu à Manchester en octobre 1945, mais il était présidé par W. E. B. Du Bois ; Kwame Nkrumah y était secrétaire d'organisation (avec George Padmore). Confondre le rôle de président et celui de secrétaire est une erreur classique.
\item \textbf{Fausse.} L'OUA a bien été fondée à Addis-Abeba le 25 mai 1963, mais l'initiative n'est pas une initiative « conjointe » de Nkrumah et Sékou Touré : elle résulte de la réunion des 32 États indépendants, et Nkrumah y défendait (en vain) les « États-Unis d'Afrique » contre le groupe de Casablanca et celui de Monrovia. La charte adoptée est un compromis qui consacre la souveraineté des États.
\item \textbf{Fausse.} Le consciencisme a été théorisé par Kwame Nkrumah dans \emph{Consciencism: Philosophy and Ideology for Decolonization} (1964). Julius Nyerere est l'auteur du \emph{Ujamaa} (socialisme communautaire tanzanien), théorisé notamment dans la Déclaration d'Arusha (1967).
\item \textbf{Vraie.} Ruben Um Nyobè, secrétaire général de l'UPC, s'est exprimé devant la IV\textsuperscript{e} Commission de l'ONU à New York en décembre 1952 (puis en 1953 et 1954) pour réclamer l'indépendance du Cameroun et la réunification des deux Cameroun (sous tutelle française et britannique).
\end{enumerate}

\textbf{Points de vigilance :} au Bac, une justification de QCM doit citer le fait exact qui invalide les autres propositions (date, lieu, auteur, fonction). Ne jamais se contenter de recopier la bonne réponse.
\end{corrige}

\begin{corrige}[Exercice 2 — Vrai / Faux justifié : Notions clés]
\begin{enumerate}
\item \textbf{Faux.} Pour Nkrumah, le communalisme désigne au contraire la propriété \emph{collective} de la terre et des moyens de production dans la société traditionnelle africaine, fondée sur la famille élargie et la solidarité lignagère : la terre appartient à la communauté, pas à l'individu.
\item \textbf{Vrai.} Dans \emph{Le Néocolonialisme, dernier stade de l'impérialisme} (1965), Nkrumah définit le néocolonialisme comme une domination qui se poursuit par des moyens économiques, financiers et culturels (aide conditionnée, sociétés multinationales, monnaies contrôlées de l'extérieur) alors même que l'indépendance politique est formellement acquise : « l'État néocolonial est, en théorie, indépendant ; en réalité, son système économique et sa politique sont dirigés de l'extérieur ».
\item \textbf{Faux.} La conscience révolutionnaire est une prise de conscience \emph{active et critique} de la situation de domination, qui engage la transformation pratique de la société ; l'adhésion passive aux traditions serait précisément ce que le consciencisme refuse (il opère une sélection critique dans l'héritage traditionnel).
\item \textbf{Vrai.} La résolution finale de Manchester (1945) réclame l'indépendance immédiate des colonies, appelle les masses (travailleurs, paysans, intellectuels) à l'action directe et fait de l'unité des peuples africains l'instrument de la libération totale du continent.
\end{enumerate}

\textbf{Points de vigilance :} la justification doit être une phrase \emph{précise} contenant l'auteur, l'ouvrage ou le concept exact ; une justification vague (« c'est évident », « tout le monde le sait ») est sanctionnée.
\end{corrige}

\begin{corrige}[Exercice 3 — Définitions attendues au Bac]
\begin{enumerate}
\item \textbf{Consciencisme} (Kwame Nkrumah, \emph{Consciencism}, 1964) : philosophie et idéologie de la décolonisation qui synthétise l'héritage traditionnel africain (communalisme), l'apport islamo-arabe et l'héritage euro-chrétien pour fonder une conscience révolutionnaire au service du socialisme africain et du développement.
\item \textbf{Panafricanisme} : mouvement politique et culturel né dans la diaspora noire (Du Bois, Garvey, Padmore) qui affirme la solidarité de tous les peuples d'origine africaine et revendique l'unité du continent comme condition de sa libération et de son émancipation.
\item \textbf{Néocolonialisme} (Nkrumah, 1965) : forme de domination par laquelle une puissance étrangère continue de contrôler l'économie et la politique d'un État formellement indépendant, par des moyens indirects (monnaie, dette, multinationales, aide conditionnée).
\item \textbf{Communalisme} (Nkrumah) : organisation socio-économique de l'Afrique traditionnelle caractérisée par la propriété collective de la terre, le travail coopératif et l'égalitarisme, présentée comme le fondement « socialiste » naturel sur lequel peut s'édifier le socialisme moderne.
\end{enumerate}

\textbf{Points de vigilance :} une définition « type Bac » comporte trois éléments : le genre (philosophie, mouvement, forme de domination, organisation), la différence spécifique (ce qui distingue la notion) et la référence (auteur, ouvrage, date). Deux lignes maximum : toute digression est pénalisée.
\end{corrige}

\begin{corrige}[Exercice 4 — Correspondances : Dates, Congrès, Figures]
\begin{center}
\begin{tabularx}{\linewidth}{|X|X|}\hline
\textbf{Colonne A} & \textbf{Correspondance exacte} \\ \hline
1. Congrès de Manchester & \textbf{B.} 1945 (5\textsuperscript{e} Congrès, présidé par Du Bois, avec Nkrumah, Kenyatta, Padmore) \\ \hline
2. Fondation de l'OUA & \textbf{D.} 25 mai 1963, Addis-Abeba (32 États signataires) \\ \hline
3. Premier Congrès panafricain (Du Bois) & \textbf{A.} 1900, Londres (organisé par Henry Sylvester Williams, discours de Du Bois : « le problème du XX\textsuperscript{e} siècle est le problème de la ligne de partage des couleurs ») \\ \hline
4. Indépendance du Ghana & \textbf{C.} 6 mars 1957, première colonie d'Afrique subsaharienne à accéder à l'indépendance, sous la conduite de Nkrumah \\ \hline
5. Création de l'Union Africaine (UA) & \textbf{E.} 2002, Durban (Afrique du Sud), succédant à l'OUA \\ \hline
\end{tabularx}
\end{center}

\textbf{Réponses condensées :} 1-B ; 2-D ; 3-A ; 4-C ; 5-E.

\textbf{Points de vigilance :} ne pas confondre le \emph{premier congrès} (Londres, 1900) et les congrès dits « de Du Bois » (Paris 1919, Londres-Bruxelles-Paris 1921, Londres-Lisbonne 1923, New York 1927) ; ne pas confondre OUA (1963) et UA (2002) : l'UA intègre des objectifs d'intégration économique et un droit d'intervention que l'OUA refusait.
\end{corrige}

\begin{corrige}[Exercice 5 — Analyse de thèse : Extrait de \emph{Consciencism} (Nkrumah)]
\begin{enumerate}
\item \textbf{Thèse principale.} Le consciencisme est une philosophie \emph{pratique et révolutionnaire} : sa finalité n'est pas la contemplation du monde mais sa transformation effective, au service de la libération et du développement de l'Afrique.
\item \textbf{Distinction d'avec la philosophie contemplative.} Nkrumah reprend ici la onzième thèse de Marx sur Feuerbach (« les philosophes n'ont fait qu'interpréter le monde, il s'agit de le transformer »). Une philosophie purement académique se contente de décrire, d'expliquer ou de spéculer ; le consciencisme, lui, est une \emph{idéologie engagée} : il part d'une situation concrète (la domination coloniale), forme une conscience critique de cette situation (conscience révolutionnaire) et débouche sur l'action politique (indépendance, socialisme, unité africaine). La vérité d'une philosophie se vérifie dans la praxis.
\item \textbf{Les trois sources idéologiques et leur articulation.}
\begin{itemize}
\item \emph{L'héritage traditionnel africain} : le communalisme (propriété collective, égalitarisme, humanisme) fournit le fondement éthique et social — il montre que le socialisme n'est pas une importation mais un approfondissement de valeurs endogènes.
\item \emph{L'apport islamo-arabe} : présent surtout en Afrique du Nord et de l'Ouest, il apporte des valeurs de discipline, de fraternité et de savoir ; il doit être intégré comme fait historique africain.
\item \emph{L'héritage euro-chrétien} : colonisation, christianisme, techniques et sciences occidentales ; il est un fait irréversible qu'il faut assumer \emph{critiquement}, en retenant la science et la technique tout en rejetant la domination.
\end{itemize}
\textbf{Articulation :} la synthèse n'est pas un simple mélange mais une \emph{harmonisation dialectique} au service d'un but unique : le développement rapide et indépendant. Le communalisme donne la base sociale, la science (issue de l'apport européen) donne les moyens techniques, la conscience révolutionnaire donne la volonté politique de transformation.
\end{enumerate}

\textbf{Points de vigilance :} la question 2 exige la référence à Marx (thèse 11 sur Feuerbach) ; la question 3 exige de montrer que la synthèse est \emph{critique} (Nkrumah ne prend pas tout de chaque courant) et \emph{finalisée} (orientée vers le développement). Éviter de présenter les trois courants comme juxtaposés.
\end{corrige}

\begin{corrige}[Exercice 6 — Le tournant de Manchester (1945) : Rupture ou continuité ?]
\begin{enumerate}
\item \textbf{Les deux ruptures majeures.}
\begin{itemize}
\item \emph{Rupture sociologique} : les congrès précédents (1900-1927) réunissaient une élite intellectuelle de la diaspora (Du Bois, Padmore) et formulaient des pétitions modérées ; Manchester accueille pour la première fois des représentants directs des masses africaines (syndicalistes, étudiants, futurs leaders politiques venus d'Afrique) : le panafricanisme devient un mouvement \emph{de masse} et \emph{africain}.
\item \emph{Rupture stratégique} : les congrès antérieurs demandaient des réformes et une représentation dans les institutions coloniales ; Manchester réclame l'\emph{indépendance immédiate et complète} des colonies et légitime l'action directe (grèves, boycotts, désobéissance civile) : le panafricanisme devient \emph{révolutionnaire}.
\end{itemize}
\item \textbf{Trois figures présentes devenues chefs d'État :} Kwame Nkrumah (Ghana, indépendance en 1957) ; Jomo Kenyatta (Kenya, indépendance en 1963) ; Hastings Kamuzu Banda (Malawi, indépendance en 1964). On peut ajouter Obafemi Awolowo (Nigeria).
\item \textbf{Naissance du panafricanisme de masse et révolutionnaire.} La résolution finale, rédigée notamment par Nkrumah, affirme que « la lutte pour le pouvoir politique par les masses coloniales est le premier pas vers la liberté complète ». Elle lie désormais trois éléments : la participation des travailleurs et paysans (masse), l'objectif d'indépendance totale et non de réforme (radicalité), et la solidarité continentale comme arme (unité). Manchester transforme ainsi le panafricanisme d'un club d'intellectuels en programme d'action politique, qui inspirera directement les indépendances des années 1957-1963.
\end{enumerate}

\textbf{Points de vigilance :} bien distinguer \emph{continuité} (l'idéal d'unité, présent depuis Du Bois) et \emph{rupture} (les moyens et les acteurs) ; une copie qui ne voit que la rupture oublie que Du Bois présidait encore Manchester : la transmission est assumée.
\end{corrige}

\begin{corrige}[Exercice 7 — Application camerounaise : Um Nyobè et l'unité]
\begin{enumerate}
\item \textbf{Contenu essentiel des interventions à l'ONU.} Devant la IV\textsuperscript{e} Commission (décembre 1952, puis 1953 et 1954), Ruben Um Nyobè, secrétaire général de l'UPC, dénonce la répression coloniale, réclame la fin du régime de tutelle, l'\emph{indépendance effective} du Cameroun (et non une indépendance octroyée et contrôlée) et la \emph{réunification} des deux Cameroun (sous tutelle française et britannique), séparés depuis 1916. Il exige aussi des élections libres et une amnistie générale.
\item \textbf{Convergences et différences avec le panafricanisme d'État.} \emph{Convergence} : comme Nkrumah, Um Nyobè pense que l'indépendance formelle sans souveraineté réelle est un leurre, et que les peuples africains doivent s'unir contre l'impérialisme. \emph{Différence} : Um Nyobè défend une unité \emph{par la base}, fondée sur la solidarité des peuples et la réunification des entités artificiellement divisées par la colonisation (les frontières héritées de Berlin), là où l'OUA de 1963 consacre au contraire le principe de l'\emph{intangibilité des frontières coloniales} et de la non-ingérence : c'est un panafricanisme d'États souverains, prudent et juridique, éloigné du fédéralisme radical des « États-Unis d'Afrique ».
\item \textbf{Exemple actuel.} \emph{Réalisation relative} : la CEMAC (Communauté économique et monétaire de l'Afrique centrale) institue une coopération régionale (marché commun, BEAC) ; le corridor Douala-N'Djamena et le pont sur le Ntem vers la Guinée équatoriale illustrent l'intégration matérielle. \emph{Échec ou limite} : la libre circulation des personnes reste incomplète, les échanges intra-africains demeurent faibles (moins de 20\,\% du commerce du continent), et la dépendance monétaire (Franc CFA) montre que la souveraineté réelle revendiquée par Um Nyobè n'est pas pleinement acquise.
\end{enumerate}

\textbf{Points de vigilance :} la question 2 est le cœur de l'exercice : il faut opposer explicitement \emph{unité des peuples} (Um Nyobè, Nkrumah) et \emph{unité des États} (OUA), avec le principe d'intangibilité des frontières comme argument décisif. L'exemple de la question 3 doit être précis et daté, non générique.
\end{corrige}

\begin{corrige}[Exercice 8 — Débat structuré : Le Franc CFA]
\begin{enumerate}
\item \textbf{Les deux thèses.}
\begin{itemize}
\item \emph{Thèse A — Solidarité et stabilité.} Argument 1 : la parité fixe avec l'euro garantit la stabilité monétaire, une inflation maîtrisée et la convertibilité, ce qui sécurise les investissements et le commerce extérieur du Cameroun. Argument 2 : la mutualisation des réserves au sein de la BEAC et de la CEMAC crée une solidarité entre États de la zone, les plus faibles bénéficiant de la crédibilité collective.
\item \emph{Thèse B — Frein et souveraineté.} Argument 1 : la politique monétaire est arrimée à des règles externes (parité fixe, centralisation d'une partie des réserves auprès du Trésor français), ce qui prive le Cameroun d'instruments de souveraineté (dévaluation, politique de crédit adaptée à l'industrialisation). Argument 2 : la surévaluation relative du CFA pénalise les exportations et l'industrie locale, favorisant les importations et perpétuant une économie de rente primaire.
\end{itemize}
\item \textbf{Analyse par le néocolonialisme.} Pour Nkrumah, le néocolonialisme est la domination poursuivie « par d'autres moyens » après l'indépendance formelle : le contrôle de la monnaie en est l'instrument parfait, car il détermine le crédit, les prix, les échanges, sans aucune présence militaire visible. La thèse B illustre exactement ce diagnostic : le Cameroun est politiquement indépendant depuis 1960, mais une partie de sa souveraineté économique s'exerce « de l'extérieur » — ce que Nkrumah appelait un État « indépendant en théorie, dirigé en réalité de l'extérieur ».
\item \textbf{Conclusion personnelle (modèle).} Au regard du critère du \emph{développement endogène} — capacité d'un pays à produire lui-même sa croissance à partir de ses ressources, de son industrie et de ses choix — le Franc CFA apparaît ambivalent : sa stabilité est réelle, mais elle est achetée au prix d'une dépendance qui bride la transformation structurelle (industrialisation, transformation locale des matières premières). Une monnaie au service du développement devrait être un instrument de politique nationale ou panafricaine, non un héritage administré. La voie conscienciste consisterait donc non pas en une sortie brutale et désordonnée, mais en la construction d'une souveraineté monétaire africaine négociée et mutualisée — par exemple une monnaie commune gouvernée par les seuls Africains — condition d'un développement réellement endogène.
\end{enumerate}

\textbf{Points de vigilance :} présenter les deux thèses avec la même force (impartialité du débat) avant de trancher ; la conclusion doit mobiliser \emph{explicitement} le critère demandé (développement endogène) et éviter la position extrême non argumentée.
\end{corrige}

\begin{corrige}[Exercice 9 — Sujet de dissertation type Baccalauréat Technique]
\textbf{Sujet : « L'unité africaine est-elle une condition nécessaire du développement de l'Afrique ? »}

\begin{enumerate}
\item \textbf{Analyse du sujet.}
\begin{itemize}
\item \emph{Termes clés.} « Unité africaine » : solidarité politique, économique et culturelle des États et peuples d'Afrique (panafricanisme, OUA/UA, intégration régionale). « Condition nécessaire » : ce sans quoi le développement est impossible (à distinguer de « condition suffisante »). « Développement » : transformation économique, sociale et humaine durable, au-delà de la seule croissance.
\item \emph{Problématique.} Si les États africains, pris isolément, sont trop faibles face à la mondialisation et au néocolonialisme, l'unité semble indispensable ; mais l'expérience (échecs de l'intégration, conflits, souverainetés jalouses) invite à demander si l'unité suffit, ou si le développement exige d'abord la bonne gouvernance interne de chaque État.
\item \emph{Plan.} I. L'unité africaine comme condition nécessaire du développement (thèse panafricaniste). II. Ses limites : l'unité ne suffit pas sans gouvernance, paix et volonté politique. III. Synthèse : une unité réelle, économique et populaire, condition nécessaire mais non suffisante.
\end{itemize}
\item \textbf{Introduction (modèle rédigé).} Lorsque le Ghana accède à l'indépendance en 1957, Kwame Nkrumah déclare que l'indépendance de son pays « n'a pas de sens si elle n'est pas liée à la libération totale de l'Afrique ». Cette affirmation pose l'une des grandes questions de la philosophie politique africaine : les cinquante-cinq États du continent, héritiers de frontières coloniales artificielles, peuvent-ils se développer séparément ? Par « unité africaine », on entend la solidarité politique et économique des peuples et États d'Afrique, héritière du panafricanisme ; par « développement », la transformation durable des conditions de vie matérielles et humaines. Le problème est alors de savoir si l'unité est une \emph{condition nécessaire} du développement — ce sans quoi il est impossible — ou seulement un facteur favorable parmi d'autres. Nous montrerons d'abord que la faiblesse des États isolés rend l'unité nécessaire ; nous verrons ensuite que l'unité seule ne suffit pas ; nous soutiendrons enfin qu'une unité réelle, économique et populaire, demeure la condition première, à compléter par la bonne gouvernance.
\item \textbf{Première partie (développement rédigé).} La thèse panafricaniste, de Du Bois à Nkrumah, repose sur un constat : la balkanisation de l'Afrique la condamne à la faiblesse. \emph{Premier argument} : la puissance économique exige des marchés vastes ; un État de quelques millions d'habitants ne peut financer seul son industrialisation. Nkrumah le formulait dès 1963 : « l'indépendance politique sans unité économique est un château de cartes ». \emph{Deuxième argument} : face aux multinationales et aux puissances extérieures, seule une négociation collective permet des rapports de force équilibrés — c'est le sens de la ZLECAf (Zone de libre-échange continentale africaine, entrée en vigueur en 2021), qui vise à créer un marché de plus d'un milliard de consommateurs et à porter le commerce intra-africain au-delà de son niveau actuel, inférieur à 20\,\%. \emph{Troisième argument} : les infrastructures de développement sont par nature transfrontalières : le corridor Douala-N'Djamena, les projets d'interconnexion électrique, la CEMAC et sa Banque centrale (BEAC) montrent qu'au Cameroun même, l'intégration régionale est déjà à l'œuvre. L'histoire institutionnelle confirme cette dynamique : de l'OUA (1963) à l'Union africaine (2002), le continent a progressivement renforcé ses instruments communs. L'unité apparaît donc bien comme le cadre sans lequel aucun développement autocentré n'est concevable.
\item \textbf{Conclusion (modèle).} L'analyse conduit à une réponse nuancée : l'unité africaine est une condition \emph{nécessaire} du développement, car aucun État isolé ne dispose de la taille, des marchés et du rapport de force requis ; mais elle n'est pas \emph{suffisante}, car elle suppose la paix, la bonne gouvernance et la volonté politique de céder une part de souveraineté. Le défi actuel n'est plus de proclamer l'unité, mais de la réaliser concrètement : faire de la ZLECAf un marché vivant, des frontières des ponts et non des murs. Reste ouverte la question que posait déjà Um Nyobè : l'unité se fera-t-elle par les États, ou par les peuples ?
\end{enumerate}

\textbf{Barème indicatif (sur 20) :} analyse du sujet et problématique : 4 pts ; introduction (amorce, définitions, problématique, plan) : 4 pts ; première partie argumentée avec exemples précis : 8 pts ; conclusion avec ouverture : 4 pts.

\textbf{Points de vigilance :} définir « condition nécessaire » (sinon toute la dissertation est faussée) ; citer au moins trois exemples précis et datés (OUA 1963, UA 2002, ZLECAf 2021, CEMAC/BEAC) ; l'ouverture de la conclusion doit être une \emph{question} ou une \emph{perspective}, non une répétition.
\end{corrige}

\begin{corrige}[Exercice 10 — Commentaire de texte : \emph{Consciencism}, Kwame Nkrumah]
\begin{enumerate}
\item \textbf{Thèse de l'auteur.} Le développement rapide de l'Afrique exige de dépasser le seul égalitarisme traditionnel par une conscience socialiste révolutionnaire, que le consciencisme fonde en synthétisant les trois héritages culturels du continent.
\item \textbf{Analyse des arguments : les trois piliers de la synthèse.}
\begin{itemize}
\item \emph{L'héritage traditionnel africain} (communalisme, égalitarisme, humanisme) : il fournit la \emph{base sociale et morale} du projet ; il prouve que le socialisme prolonge des valeurs endogènes et n'est pas une importation étrangère. Mais Nkrumah pose clairement sa limite : « l'égalitarisme traditionnel ne suffit pas », car la société traditionnelle est statique et ne possède ni la technique ni l'organisation d'une économie moderne.
\item \emph{L'apport islamo-arabe} : il est reconnu comme un fait historique constitutif de l'Afrique (savoir, discipline, fraternité) ; sa fonction est d'\emph{intégrer} toutes les composantes culturelles du continent dans un projet commun, refusant toute hiérarchie des cultures.
\item \emph{L'héritage euro-chrétien} : il apporte la \emph{science, la technique et les méthodes d'organisation moderne}, indispensables au « développement rapide » ; mais il doit être assumé \emph{critiquement}, en rejetant la domination coloniale et l'individualisme capitaliste qui l'accompagnent.
\end{itemize}
La « conscience révolutionnaire » est le ciment de la synthèse : elle est la prise de conscience active qui transforme la structure sociale au lieu de la subir. Le raisonnement est dialectique : ni retour pur au passé (traditionalisme), ni imitation de l'Occident (aliénation), mais dépassement créateur.
\item \textbf{Discussion de la pertinence actuelle (modèle, environ 20 lignes).} La synthèse conscienciste conserve une pertinence réelle face aux défis camerounais actuels, mais à condition d'être actualisée. D'abord, son diagnostic sur l'insuffisance de l'indépendance formelle reste valable : la question de la souveraineté économique se pose encore à travers la dette (qui contraint les politiques publiques via les programmes d'ajustement), la monnaie (Franc CFA) et la dépendance aux exportations de matières premières brutes (pétrole, cacao, bois). Le néocolonialisme que Nkrumah analysait en 1965 a changé de visage — institutions financières internationales, multinationales, accords commerciaux asymétriques — mais non de logique. Ensuite, l'idée d'une synthèse critique plutôt que d'un rejet ou d'une imitation est précieuse : le Cameroun a besoin de la science et de la technologie mondialisées (numérique, agro-industrie), mais il peut les mettre au service de valeurs endogènes de solidarité (tontines, travail coopératif, économie sociale), ce qui rejoint l'esprit du communalisme modernisé. Cependant, trois limites apparaissent. Premièrement, Nkrumah sous-estimait le problème de la \emph{gouvernance} : la conscience révolutionnaire ne suffit pas contre la corruption, la mauvaise gestion et l'ethnicisation du pouvoir, qui sont aujourd'hui des obstacles internes majeurs au développement. Deuxièmement, la mondialisation impose des contraintes (concurrence, normes) que le cadre national-étatique de Nkrumah ne prévoyait pas ; c'est pourquoi sa dimension panafricaniste (ZLECAf) est plus actuelle que jamais. Troisièmement, la « synthèse des trois courants » doit être élargie : l'Afrique d'aujourd'hui dialogue aussi avec l'Asie et construit ses propres voies. En définitive, le consciencisme demeure opératoire non comme recette toute faite, mais comme \emph{méthode} : assumer tous les héritages, choisir souverainement, transformer au service du peuple. C'est à ce prix — unité, souveraineté réelle, gouvernance responsable — que la philosophie de Nkrumah peut encore éclairer le développement du Cameroun.
\end{enumerate}

\textbf{Barème indicatif (sur 20) :} thèse dégagée : 3 pts ; analyse des trois piliers et de la conscience révolutionnaire : 8 pts ; discussion argumentée (au moins deux arguments pour, deux limites, exemples camerounais) : 9 pts.

\textbf{Points de vigilance :} dans un commentaire, toujours partir du texte (citer ses expressions) avant de mobiliser ses connaissances ; la discussion doit être \emph{argumentée} (thèse, antithèse, position personnelle justifiée), jamais un simple résumé ; les 15 lignes minimum sont une condition formelle sanctionnée.
\end{corrige}

\begin{corrige}[Exercice 11 — Question à réponse construite : Fondements et Développement]
\begin{enumerate}
\item \textbf{Fondements idéologiques.}
\begin{itemize}
\item \emph{Le communalisme (base traditionnelle)} : la société africaine précoloniale repose sur la propriété collective de la terre, le travail coopératif et l'égalitarisme lignager. Pour Nkrumah, cette organisation contient déjà en germe le socialisme : il n'y a donc pas à « importer » une idéologie étrangère, mais à \emph{moderniser} un fonds endogène.
\item \emph{Le socialisme (objectif moderne)} : l'égalitarisme traditionnel, statique, doit être dépassé en un socialisme scientifique capable d'industrialiser le pays, de planifier l'économie et de répartir équitablement les richesses. Le socialisme est la forme moderne du communalisme.
\item \emph{L'anti-impérialisme et l'anti-néocolonialisme (lutte politique)} : le développement suppose d'abord la libération ; or l'indépendance formelle peut masquer une domination économique continue (néocolonialisme). La lutte idéologique vise donc la souveraineté réelle : contrôle des ressources, de la monnaie, des choix économiques.
\end{itemize}
\item \textbf{Fondements philosophiques.}
\begin{itemize}
\item \emph{Le matérialisme dialectique adapté} : Nkrumah affirme le primat de la matière (la réalité économique et sociale détermine les idées), mais il l'assouplit : il refuse le matérialisme mécaniste et reconnaît à la conscience un rôle actif, capable de transformer la matière sociale. La dialectique (thèse-antithèse-synthèse) décrit l'histoire africaine : société traditionnelle, choc colonial, synthèse socialiste.
\item \emph{La conscience révolutionnaire} : elle est la prise de conscience collective de la domination et de la possibilité de la renverser ; elle convertit le peuple de spectateur en acteur de son histoire. Nkrumah y voit une forme de « catharsis » : la libération intérieure (mentale) précède et accompagne la libération extérieure (politique et économique).
\item \emph{La synthèse des trois courants} : traditionnel africain (éthique communautaire), islamo-arabe (discipline, savoir), euro-chrétien (science, technique). Le consciencisme les harmonise \emph{critiquement} : il retient de chacun ce qui sert le développement et rejette ce qui aliène (domination, individualisme, fatalisme).
\end{itemize}
\item \textbf{Articulation au développement : de l'indépendance formelle à la souveraineté réelle.} Le raisonnement conscienciste se déploie en trois temps. \emph{Premièrement}, la conscience révolutionnaire détruit le complexe d'infériorité colonial et restaure la confiance endogène : sans estime de soi, aucune politique de développement n'est appropriée par le peuple — d'où l'importance de l'\emph{éducation} (enseignement technique, valorisation des cultures et langues nationales au Cameroun). \emph{Deuxièmement}, le socialisme planifié oriente les ressources vers la transformation structurelle : politique agricole vivrière et de rente (cacao, café, coton), industrialisation par transformation locale (complexes agro-industriels, aluminium d'Édéa historiquement, projets de transformation du bois sur place plutôt qu'exportation de grumes), infrastructures (barrages, routes). \emph{Troisièmement}, l'anti-néocolonialisme impose l'unité continentale comme horizon : un État seul ne peut négocier face aux multinationales ; la ZLECAf (2021) et la CEMAC concrétisent aujourd'hui cette exigence de marché intégré. Le passage de l'indépendance formelle (1960) à la souveraineté réelle suppose donc simultanément la libération mentale, la transformation économique et l'unité africaine : c'est toute la cohérence du consciencisme.
\item \textbf{Conclusion : limites dans le contexte camerounais actuel.} L'application rencontre des limites réelles : la gouvernance (corruption, centralisation) détourne parfois les politiques publiques de leurs finalités sociales ; la dépendance monétaire et commerciale persiste (Franc CFA, exportations brutes, dette) ; l'intégration régionale avance lentement (libre circulation incomplète, faibles échanges intra-africains). Le consciencisme reste ainsi davantage un \emph{horizon normatif} — une boussole : souveraineté, unité, socialisme endogène — qu'un bilan accompli. Sa valeur est de rappeler que le développement n'est pas seulement une affaire de capitaux, mais d'abord de conscience, de volonté politique et d'unité.
\end{enumerate}

\textbf{Barème indicatif (sur 20) :} fondements idéologiques : 5 pts ; fondements philosophiques : 5 pts ; articulation au développement avec exemples camerounais précis : 7 pts ; conclusion évaluative : 3 pts.

\textbf{Points de vigilance :} bien distinguer le plan \emph{idéologique} (communalisme, socialisme, anti-néocolonialisme) du plan \emph{philosophique} (matérialisme dialectique, conscience révolutionnaire, synthèse des trois courants) — les confondre est l'erreur la plus fréquente ; chaque affirmation sur le développement doit être adossée à un exemple camerounais daté et vérifiable ; la conclusion doit \emph{évaluer} (porter un jugement nuancé), pas résumer.
\end{corrige}

\newpage

\coursec{Fiche bilan}

\begin{popfigure}
\begin{tikzpicture}[mindmap, grow cyclic, every node/.style={concept, font=\small},
  root concept/.append style={concept color=popblue, font=\bfseries, minimum size=3.2cm},
  level 1/.append style={level distance=3.6cm, sibling angle=51.4},
  level 1 concept/.append style={font=\scriptsize, minimum size=2.4cm}]
\node[root concept]{Consciencisme, panafricanisme et développement de l'Afrique}
  child[concept color=poporange] { node {Contexte historique : colonisation, indépendances, quête d'identité} }
  child[concept color=popgreen] { node {Fondements idéologiques du consciencisme : Nkrumah, socialisme africain} }
  child[concept color=poppurple] { node {Fondements philosophiques : conscience, liberté, personne humaine} }
  child[concept color=poppink] { node {Historique du panafricanisme : congrès, OUA, UA} }
  child[concept color=popgold] { node {Unité africaine et développement : applications au Cameroun} }
  child[concept color=popblueL] { node {Méthode : rédiger et justifier une démarche argumentative} }
  child[concept color=popgreenL] { node {Problématique et définitions clés} };
\end{tikzpicture}
\legende{Carte mentale de la leçon : les sept branches essentielles à maîtriser pour le Probatoire / Baccalauréat technique.}
\end{popfigure}

\begin{aretenir}[L'essentiel de la leçon]
\textbf{Définitions clés}
\begin{itemize}
  \item \cle{Consciencisme} : doctrine philosophique et politique de Kwame Nkrumah fondée sur la prise de conscience par les Africains de leur situation coloniale et de leur identité, en vue de la libération et de l'unité de l'Afrique.
  \item \cle{Panafricanisme} : mouvement politique et culturel visant l'unité et la solidarité de tous les peuples d'Afrique et de la diaspora africaine.
  \item \cle{Colonisation} : domination politique, économique et culturelle d'un territoire africain par une puissance étrangère.
  \item \cle{Développement} : transformation économique, sociale et culturelle d'un pays pour améliorer les conditions de vie des populations.
  \item \cle{Conscience} : connaissance que l'homme a de lui-même et du monde ; elle rend possible la liberté et l'action responsable.
\end{itemize}

\textbf{Repères à connaître par c\oe ur}
\begin{center}
\begin{tabularx}{\linewidth}{|l|X|}\hline
\rowcolor{popblueL} \textbf{Repère} & \textbf{Contenu} \\ \hline
Auteur central & Kwame Nkrumah (Ghana), \emph{Consciencisme} (1964) \\ \hline
Fondement idéologique & Socialisme africain, refus de l'aliénation coloniale, unité politique de l'Afrique \\ \hline
Fondement philosophique & La conscience comme source de la liberté et de la dignité de la personne humaine \\ \hline
Étapes du panafricanisme & Congrès panafricains (1900, 1919-1945), indépendances, création de l'OUA (1963), puis de l'UA (2002) \\ \hline
Application camerounaise & Indépendance (1960), réunification (1961), intégration dans l'OUA/UA, coopération sous-régionale (CEMAC) \\ \hline
\end{tabularx}
\end{center}

\textbf{Méthodes express}
\begin{itemize}
  \item \textbf{Dissertation} : définir les termes du sujet, poser une problématique (l'unité africaine est-elle une condition du développement ?), construire deux ou trois parties argumentées avec exemples camerounais et africains.
  \item \textbf{Justifier une réponse} : énoncer la thèse, donner l'argument, citer un auteur ou un fait historique, conclure clairement.
  \item \textbf{Appliquer à une situation concrète} : relier la notion (conscience, unité, développement) à un fait de la vie courante (coopération, intégration régionale, éducation citoyenne).
\end{itemize}
\end{aretenir}

\begin{aretenir}[Checklist avant le Bac]
\begin{itemize}
  \item $\square$ Je sais définir le consciencisme et citer son auteur principal.
  \item $\square$ Je sais définir le panafricanisme et distinguer ses grandes étapes.
  \item $\square$ Je connais les dates clés : indépendance du Cameroun (1960), réunification (1961), OUA (1963), UA (2002).
  \item $\square$ Je sais expliquer les fondements idéologiques du consciencisme (socialisme africain, libération, unité).
  \item $\square$ Je sais expliquer ses fondements philosophiques (conscience, liberté, dignité de la personne).
  \item $\square$ Je sais relier colonisation, indépendances et quête d'identité africaine.
  \item $\square$ Je sais montrer le lien entre unité africaine et développement.
  \item $\square$ Je sais donner des exemples camerounais concrets (CEMAC, coopération, éducation).
  \item $\square$ Je sais construire une problématique à partir d'un sujet de dissertation.
  \item $\square$ Je sais rédiger une démarche argumentative en trois temps : thèse, argument, exemple.
  \item $\square$ Je sais rédiger une introduction (accroche, définitions, problématique, annonce du plan) et une conclusion.
  \item $\square$ Je relis ma copie pour vérifier l'orthographe des noms propres et des dates.
\end{itemize}
\end{aretenir}

\findecours

\end{document}
